\section{Scope}\label{sec:scope}

An obstruction theorem earns its keep only if its boundary is drawn exactly.
This section states what each result does not cover.

\subsection{What \texorpdfstring{\Cref{thm:divisor}}{the divisor obstruction} does not cover}

\begin{remark}[Markman's construction]\label{rem:markmanscope}
The objects constructed by Markman \cite{Mar25a,Mar25b} are not obtained by
restricting from an ambient projective space, and they are not
$K$-equivariant. \Cref{sec:escape} analyses exactly which hypotheses they
escape and why, and \Cref{prop:whichescaped} lists the three. Nothing in this
paper bears on them, and nothing must: their conclusion is a theorem.
\end{remark}

\begin{remark}[Twisted objects and gerbes]\label{rem:twisted}
The hypothesis of \Cref{thm:divisor} is that $\cF$ is an object of
$D^{b}(\PP^{N-1})$, so a twisted sheaf on a gerbe over $\PP^{N-1}$, or an
object of a derived category of twisted sheaves, is not covered. What is
covered, and what the proof actually uses, is that the Chern character of the
restricted object lies in the subring of $H^{*}(\Av,\QQ)$ generated by
divisor classes. Any construction whose input has that property is subject to
the conclusion, whatever category it is phrased in.
\end{remark}

\begin{remark}[Other ambient varieties]\label{rem:otherambient}
Replacing $\PP^{N-1}$ by another ambient variety $P$ changes the conclusion
exactly when $H^{*}(P,\QQ)$ is not generated by divisor classes, since the
proof only uses $\ch(\cF)\in D^{*}(P)$ and the fact that $i^{*}$ carries
$D^{*}(P)$ into $D^{*}(\Av)$. For $P$ a Grassmannian, a flag variety, or any
variety whose cohomology is generated in degree two, nothing changes. For $P$
with exceptional Hodge classes of its own the argument does not apply, but
then the problem has been moved rather than solved.
\end{remark}

\subsection{What \texorpdfstring{\Cref{thm:character}}{the character obstruction} does not cover}

\begin{remark}[Constructions that are not equivariant]\label{rem:nonequivariant}
\Cref{thm:character} constrains eigenclasses. A construction whose output is
not a $\Kx$-eigenclass is not constrained, and by \Cref{rem:rightcharacter}
such a construction must break the naturality of the $K$-action in an
essential way: it must depend on a choice that is not $K$-invariant, and its
output must have components in both extreme summands of the grading at once.
We know of no systematic way to arrange this, and we regard the question of
whether one exists as the interesting one left by this paper.
\end{remark}

\begin{remark}[Abelian varieties that are not very general]\label{rem:special}
At special points of the moduli of Weil type the endomorphism algebra grows,
the Neron--Severi group grows with it, and the Weil line can meet the divisor
subring. This is the split case, where Weil classes are known to be algebraic
by elementary means, and it is consistent with everything here: our statements
about $\HW(A)\cap D^{2n}(A)$ assume $\NS(A)\otimes\QQ=\QQ\eta$, which fails
there. \Cref{thm:character} itself makes no such assumption and holds at every
point, but at special points a class can be a sum of eigenclasses for several
characters and the conclusion applies to each summand separately.
\end{remark}

\subsection{What \texorpdfstring{\Cref{cor:kernelfamily}}{the semiregularity statement} does not cover}

\begin{remark}[Objects of positive rank]\label{rem:positiverankscope}
\Cref{thm:rankzero} and \Cref{cor:kernelfamily} both require the Chern character to
lie on the Weil line, so that $z\mapsto z\cdot\ch(\cE)$ has no degree-two
term. An object of nonzero rank has $\rk(\cE)\,z$ in that degree, the map is
injective, and nothing here applies. That is the class Markman's sheaves
belong to, and \Cref{rem:familyuniform} records why the distinction is not
incidental to his argument.
\end{remark}

\begin{remark}[The weaker semiregularity hypothesis]\label{rem:weakerscope}
What \Cref{cor:kernelfamily} obstructs is the injectivity of $\sigma_{\cE}$ on the
whole of $\Ext^{2}(\cE,\cE)$. The variant proposed in
\cite[Question 11.4]{Mar25b}, injectivity on the image of an evaluation map,
asks less. What we prove about it is confined to \Cref{thm:weakfails},
\Cref{prop:evaluation}, \Cref{thm:linebundle} and \Cref{cor:noline}, where it
fails, to \Cref{thm:factor}, where for a secant object of positive rank it
is shown equivalent to two conditions on the factor, and to
\Cref{thm:candidatefails} and \Cref{cor:localobstruction}, where it fails
for Markman's candidate in dimension eight and for every secant object whose
support has two transversal branches of codimension two through a point. We
do not decide the variant for any other object of positive rank, and we do
not decide whether an affirmative answer to the question would follow from
it.
\end{remark}

\begin{remark}[Deformation of the class itself]\label{rem:notnondeform}
A kernel is not an obstruction class. \Cref{cor:kernelfamily} says that a sufficient
criterion for the Weil class to stay algebraic is unavailable at the objects
in question; it does not say the class fails to stay algebraic, and no
statement in this paper does. \Cref{q:residual} remains open in both
directions.
\end{remark}

\subsection{What \texorpdfstring{\Cref{thm:hyperplane}}{the hyperplane obstruction} does not cover}

\begin{remark}[Families that are not linear systems on \texorpdfstring{$X$}{X}]%
\label{rem:notlinear}
\Cref{thm:hyperplane} is about smooth members of an ample linear system on
$X$, and about the maps $j^{*}$ and $j_{*}$ they carry. It does not obstruct
an argument that deforms $X$ itself, that passes to a different variety by a
correspondence, or that uses a family in which $X$ appears as a fibre and the
other fibres have the same dimension. In the last case the inductive parameter
does not decrease, which is a different difficulty and the one recorded in
\Cref{rem:dimensionbookkeeping}.
\end{remark}

\begin{remark}[Coniveau]\label{rem:coniveauscope}
A Hodge class of positive coniveau is supported on a proper closed
subvariety, and the reduction to that subvariety is legitimate and lowers the
dimension. Nothing here obstructs it. The obstruction concerns classes of
coniveau zero, which is where the primitive middle cohomology of a general
variety sits, and the generalised Hodge conjecture predicts that these are
precisely the ones a coniveau argument cannot reach.
\end{remark}

\begin{remark}[The Kuga--Satake route]\label{rem:ksscope}
For a weight-two Hodge structure of K3 type the Kuga--Satake construction
produces an abelian variety $\mathrm{KS}(T)$ and an injection of $T$, twisted
once, into the endomorphisms of its first cohomology, carrying the Hodge class
in question to a class of type $(0,0)$, that is to an endomorphism, whose graph
is a cycle. The conclusion is transported back by a correspondence, not by a
restriction, so \Cref{thm:hyperplane} does not apply; this is
\Cref{prop:ksnotrestriction}. The route is limited by the requirement
$h^{2,0}=1$, and \Cref{prop:normpart} locates that limit exactly for abelian
varieties of Weil type: no rational sub-Hodge structure of $H^{2}(A,\QQ)$ is of
K3 type unless $n=1$. \Cref{sec:ks} treats the construction in full, including
what remains open about the algebraicity of the correspondence itself.
\end{remark}

\subsection{What stands between this and the conjecture}
\label{ssec:closuremap}

The results above close several routes and narrow one. We set down, once
and precisely, what remains, because the remainder is not a single
obstruction whose removal would finish the problem. It is a list, and we
record it with the exact statement of each entry and the exact place where
this paper leaves it; \Cref{prop:descent} records the one way in which the
entries on Weil classes are linked.

Fix an imaginary quadratic field $K=\QQ(\sqrt{-d})$, an integer $n\ge2$, and a
class $\delta\in\QQ^{\times}_{>0}/\Nm_{K/\QQ}(\Kx)$, the discriminant of the
hermitian form. Write
\begin{quote}
$\mathbf{W}(K,n,\delta)$: every abelian $2n$-fold of $(K,-1,n)$-Weil type with
discriminant $\delta$ has algebraic Weil classes.
\end{quote}

\begin{proposition}[The Weil statements, and how they are linked]\label{prop:separate}
For each triple $(K,n,\delta)$ the following hold.
\begin{enumerate}[label=\textup{(\roman*)},leftmargin=2.6em]
\item $\mathbf{W}(K,n,\delta)$ is equivalent to the algebraicity of the single
  class $\omega_{s}$ at every point $s$ of the period domain $\cD_{n,n}$, and
  that domain is connected; this is \Cref{cor:exactform}.
\item $\mathbf{W}(K,3,\delta_{0})$ holds for the trivial discriminant
  $\delta_{0}$ \cite{Mar25a,Mar25b}, and with it, by \Cref{prop:descent},
  $\mathbf{W}(K,2,\delta)$ for every $K$ and every $\delta$, which is how
  Markman obtains the fourfolds. Nothing else is known.
\item The period domains attached to distinct triples are disjoint, and
  \Cref{thm:residual} and \Cref{cor:exactform} are statements about one
  triple at a time. The triples are nevertheless not independent: by
  \Cref{prop:descent}, $\mathbf{W}(K,N,\delta'')$ implies
  $\mathbf{W}(K,n,\delta)$ for every $n<N$ and every $\delta$, so the families
  of trivial discriminant in infinitely many dimensions give every imaginary
  quadratic family, and the families of a CM field of degree at least four
  give those of the imaginary quadratic fields it contains.
\end{enumerate}
\end{proposition}

\begin{proof}
(i) is \Cref{cor:exactform} together with \Cref{prop:weildomain}. (ii) is the
cited work together with \Cref{prop:descent}(i); \Cref{rem:residualmap}
records that the construction of \Cref{ssec:split} reaches only $\delta_{0}$.
For (iii), the discriminant is a discrete invariant of the hermitian form and
is constant on $\cD_{n,n}$ by \Cref{prop:splitgeom}(ii), while $n$ and $K$ are
fixed by the endomorphism algebra, so the domains are disjoint; the links are
\Cref{prop:descent}, which passes between families through a product with a
Weil surface or through a scalar extension, and not through a deformation.
An earlier version of this proposition asserted that settling one triple
leaves every other exactly where it was. That was wrong: it considered only
the constructions of this paper, and the product argument of Schoen that
Markman uses to pass from sixfolds to fourfolds already links the triples.
\end{proof}
\begin{proposition}[Descent and scalar extension]\label{prop:descent}
Let $F$ be a CM field of degree $2m$ and write $\mathbf{W}(F,n,\delta)$ for
the statement that every abelian variety of $(F,n,\delta)$-Weil type has
algebraic Weil classes, so that $\mathbf{W}(K,n,\delta)$ is the statement
above when $m=1$.
\begin{enumerate}[label=\textup{(\roman*)},leftmargin=2.6em]
\item \emph{Descent.} Let $B$ be of $(F,n,\delta)$-Weil type and $Y$ of
  $(F,1,\delta')$-Weil type. Then $X=B\times Y$, with the diagonal action of
  $F$ and the product polarisation, is of $(F,n+1,\delta\delta')$-Weil type,
  and the correspondence
  \[
    P(x)=pr_{B*}\bigl(x\cdot pr_{Y}^{*}(\eta_{Y}^{2m-2}\,y')\bigr),
  \]
  with $\eta_{Y}$ the polarisation of $Y$ and $y'$ a nonzero Weil class of
  $Y$, maps the Weil classes of $X$ isomorphically onto those of $B$. Since
  $y'$ lies in $H^{2}(Y)$ it is algebraic, so $P$ is algebraic, and every
  discriminant, a totally positive class in
  $F_{0}^{\times}/\Nm_{F/F_{0}}(F^{\times})$, is that of some $Y$.
  Hence $\mathbf{W}(F,N,\delta'')$ for one $N$ and one $\delta''$ implies
  $\mathbf{W}(F,n,\delta)$ for every $n<N$ and every $\delta$.
\item \emph{Scalar extension.} Let $K\subset F$ be CM fields and $B$ of
  $(K,n,\delta)$-Weil type. Then $B_{F}=B\otimes_{\cO_{K}}\cO_{F}$ is of
  $(F,n,\iota(\delta))$-Weil type, and for $j\colon B\to B_{F}$,
  $b\mapsto b\otimes1$, one has
  $j^{*}\Tr_{F/\QQ}(a\det_{F})=\Tr_{K/\QQ}(\Tr_{F/K}(a)\det_{K})$, so $j^{*}$
  maps the Weil classes of $B_{F}$ onto those of $B$. Hence
  $\mathbf{W}(F,n,\iota(\delta))$ implies $\mathbf{W}(K,n,\delta)$; every
  imaginary quadratic $K$ lies in the quartic CM field $K(\sqrt2)$.
\end{enumerate}
\end{proposition}

\begin{proof}
(i) $H^{1}(X)=H^{1}(B)\oplus H^{1}(Y)$ as $F$-spaces and the hermitian form
of $X$ is the orthogonal sum, of signature $(n+1,n+1)$ at every real place;
its normalised discriminant is the product, and $H_{Y}=\mathrm{diag}(1,-t)$
with $t$ totally positive realises every class. The Weil classes of $X$ are
$\bigoplus_{\sigma}\bigwedge^{2n}V_{B,\sigma}\otimes\bigwedge^{2}V_{Y,\sigma}$,
because the top exterior power of a direct sum is the tensor product of the
top powers. On the $\sigma$-component, $P$ is multiplication by
$\int_{Y}w_{Y}^{\sigma}\,\eta_{Y}^{2m-2}\,y'^{\bar\sigma}$, which is nonzero
since $w_{Y}^{\sigma}\wedge y'^{\bar\sigma}$ fills the $(\sigma,\bar\sigma)$
block and $\eta_{Y}$ is nondegenerate on the others; the remaining pairings
vanish for weight reasons under the norm-one torus of $F$. So $P$ is
rational, diagonal over $\CC$ with nonzero entries, and algebraic, and
$\omega(B)$ is algebraic whenever $\omega(X)$ is; \Cref{lem:onecomponent}
passes from one Weil class to all of them. (ii) $H^{1}(B_{F})=H^{1}(B)\otimes_{K}F$,
each eigenspace of $F$ is an eigenspace of $K$ of signature $(n,n)$, the
hermitian form is $H_{B}\otimes F$ with determinant $\det H_{B}$, and a norm
from $K$ is a norm from $F$. The identity for $j^{*}$ is the multiplicativity
of the determinant under restriction of scalars, and $\Tr_{F/K}$ is onto.
Item (LIV) of \Cref{sec:verification} checks the linear algebra of both
statements exactly, over seven CM fields.
\end{proof}

The case $K$ imaginary quadratic and $N=3$ of (i) is the degeneration argument
of Schoen through which Markman obtains every Weil fourfold, at every
discriminant, from the sixfolds of trivial discriminant \cite{Mar25a}; what
(i) adds is that the same argument works in every dimension and for every CM
field. Its consequence for the structure of the problem is not small, and it
corrects an earlier version of this section, which asserted that the triples
$(K,n,\delta)$ are separate problems. They are not: the trivial discriminant
families in infinitely many dimensions give every imaginary quadratic family,
and by (ii) the split families of the CM fields of degree at least four give
the imaginary quadratic fields as well. In particular the base points of
nontrivial discriminant supplied by \Cref{thm:everyfamily} are not needed on
any route to the conjecture, since the split families already carry base
points \eqref{eq:splitclosed}; what \Cref{thm:everyfamily} shows is that
propagation can be started in every family, not that it has to be.

\noindent
With that said, here is the list.

\begin{enumerate}[label=\textup{(C\arabic*)},leftmargin=3.0em,itemsep=4pt]
\item \emph{One secant object on one fourfold.} There should exist a
  principally polarised abelian fourfold $(X,\Theta)$ and a perfect complex
  $\cF$ on it with $\ch(\cF)=u_{\Theta}+3v_{\Theta}$ satisfying
  \eqref{eq:heisenberg}. By \Cref{thm:factor} and \Cref{cor:markmancondition}
  the transform $\Phi(\cF\boxtimes\cF)$ would then satisfy the weakened
  criterion, and by \Cref{prop:markman8} the three conditions of Markman's
  strategy that are conditions on the Chern character hold at the point
  $(1,3)$ for any such $\cF$, since they are conditions on $\ch(\cF)$ alone.
  Turning that into algebraicity needs one thing more, and it is not ours to
  supply. The weakened criterion enters in place of the hypothesis Markman
  verifies, and whether the semiregularity theorem survives the substitution
  is exactly \cite[Question 11.4]{Mar25b}, which is open. So (C1) is a demand
  on a construction, and an affirmative answer to that question is a separate
  prerequisite standing beside it; \Cref{thm:closure} places both.
  What this paper contributes is the exact form of the demand. The local part
  of the criterion is a condition on
  depth and nothing else: if the support is Cohen--Macaulay there is no local
  obstruction at any point (\Cref{thm:cmvanishing}), and if it is not, at a
  point where two components meet only at that point or where there is an
  embedded point, the criterion fails outright
  (\Cref{thm:candidatefails}, \Cref{cor:localobstruction},
  \Cref{cor:cmdichotomy}). The candidate of \cite[Section 12]{Mar25b} is of the
  second kind, and it cannot be moved to the first: separating the branches is
  what would restore the depth, and \Cref{rem:normalisationcm} shows that doing
  so moves the Chern character off the secant plane. So what is needed is a
  Cohen--Macaulay support, and for one of those the whole criterion is the
  single vanishing of the image of $A_{\theta}A_{\theta'}$ in
  $H^{1}(Z,\cE xt^{1}(\cI_{Z},\cI_{Z}))$. Two shapes of such a support remain
  (\Cref{cor:disjointroutes}): a smooth one, where the invariants are forced
  outright and $d\le9$ (\Cref{thm:smoothinvariants},
  \Cref{cor:fourdiscriminants}), and a singular one; in neither may the
  Hilbert--Burch resolution split into powers of $\Theta$ or into simple
  semi-homogeneous bundles (\Cref{thm:nosplit}, \Cref{thm:noblocks}), which
  removes every candidate the search of
  \Cref{thm:burchsearch} found, and no complete intersection occurs
  (\Cref{thm:noci}). \Cref{rem:routemap} states the two remaining questions in
  the form a construction would have to answer.
\item \emph{The same for every $n$ and every discriminant.} A positive answer
  to (C1) reaches one triple, and the discriminant of that triple is the
  trivial one, because the object lives on $X\times\widehat X$ and every split
  member has trivial discriminant by \Cref{prop:splitgeom}(ii). The other
  triples are reached through \Cref{prop:descent}: a positive answer in
  infinitely many dimensions gives $\mathbf{W}(K,n,\delta)$ for every $n$ and
  every $\delta$, so only uniformity in $n$ is asked for, and (C1) at $n=4$,
  with Question~11.4, already gives every sixfold. What no answer of this kind reaches is
  the CM fields of degree at least four.
  What this paper contributes: three of the four
  ingredients of \Cref{q:higher} are supplied for every $n$ and every $K$
  (\Cref{thm:twoofthefour}), including the existence of the secant sheaves
  themselves (\Cref{thm:secantexists}); the fourth is untouched
  (\Cref{rem:deformation}).
\item \emph{CM fields of degree larger than two.} For a CM field of degree
  $2m>2$ acting on an abelian variety, the Weil classes span a space of
  dimension $2m$, and the secant framework of \Cref{sec:construction}, which
  is built on a plane meeting the spinor variety in a conjugate pair, does not
  apply. What this paper contributes: everything except the propagation. By
  \Cref{sec:cmfields} each such family has a connected period domain, a flat
  space of Weil classes, a dense set of rational orbits, a base point at a CM
  member where the classes are written as polynomials in divisor classes
  (\Cref{thm:cmbasepoint}), and the two propagation criteria
  (\Cref{thm:cmpropagation}); so the higher CM fields are the same
  propagation problem as the quadratic ones, and not a separate entry.
  \Cref{app:albert} places the imaginary quadratic case inside Albert's
  classification.
\item \emph{The non-split sixfolds.} $\mathbf{W}(K,3,\delta)$ is open for
  every $\delta\ne\delta_{0}$; the case $\delta=\delta_{0}$ is Markman's
  \cite{Mar25a}. By \Cref{prop:descent}(i) it would follow from
  $\mathbf{W}(K,4,\delta_{0})$, and by \Cref{prop:descent}(ii) from
  $\mathbf{W}(K(\sqrt2),3,\iota(\delta))$. What this paper contributes: a base
  point in every family, at a member with quaternionic multiplication
  (\Cref{thm:everyfamily}); the locus of such members is a Lagrangian
  Grassmannian of codimension $n(n-1)/2$ (\Cref{prop:lagrangian}), so
  propagation off it is needed and is not supplied.
\item \emph{Exceptional classes that are not of Weil type.} On an abelian
  variety the Hodge ring can contain exceptional classes that the Weil
  construction does not produce. Through dimension five it does not
  \cite{MZ99}, and in dimension eight it does: the self-product of an abelian
  fourfold of Mumford type carries two Hodge classes that are neither
  products of divisor classes nor Weil classes of any quotient
  (\Cref{prop:mumfordproduct}). So the Hodge conjecture for abelian varieties
  does not reduce to Weil classes, and the statement that remains is the
  algebraicity of these classes themselves. What this paper contributes: the
  example, made exact, and nothing towards their algebraicity. Every positive
  statement here is about classes on the Weil line.
\item \emph{Varieties that are not abelian.} No reduction of the Hodge
  conjecture for arbitrary smooth projective varieties to abelian varieties is
  known, and none of Hodge-theoretic kind can exist: the $H^{3}$ of a very
  general quintic threefold lies outside the tensor category generated by the
  cohomology of abelian varieties (\Cref{prop:noabeliantype}), so its Hodge
  structure cannot be realised inside that of an abelian variety by any
  correspondence. \Cref{thm:hyperplane} and \Cref{thm:returntrip} close the
  two routes through hyperplane sections that would be the natural candidates.
  What this paper contributes: those three obstructions, and nothing positive.
  Nor is this entry the complement of the others. The varieties that are not
  abelian include $A\times\PP^{1}$ for every abelian variety $A$, and the
  conjecture for $A\times\PP^{1}$ gives it for $A$, so the conjecture for
  the varieties that are not abelian is the whole conjecture
  (\Cref{prop:f3ishc}).
\end{enumerate}

\begin{remark}[The shape of the remainder]\label{rem:remaindershape}
Two features of this list deserve emphasis, because they are the reason the
paper does not claim more than it does.

The narrowing is real but bounded. \Cref{cor:fourdiscriminants} removes all
but four discriminants from the smooth case of (C1), and
\Cref{cor:localobstruction} removes an entire class of supports in every
dimension. Those are restrictions on the space of possible proofs, not steps
towards one; they say where a construction cannot be, and the set of places it
could be remains infinite in every direction except the one that
\Cref{cor:fourdiscriminants} bounds.

The equivalence of \Cref{cor:exactform} is not an obstruction. It says that
for a fixed triple the Hodge conjecture in every degree is the algebraicity of
one class at every point of one connected domain. That is a restatement of the
problem in its smallest form, and a restatement is worth having, but the
implication it isolates is the conjecture itself for those varieties, and
proving it would close those varieties only.

The third feature is the one a list cannot show, and it is the subject of
\Cref{ssec:closure}, once \Cref{ssec:bullseye} has supplied the implications
of the literature. The entries above are not of equal standing. Some of them lie
on every route to the conjecture and some lie on none, and which is which is
not a matter of judgement: it is decided by the implications the paper
actually proves.
\end{remark}

\subsection{The target common to the short routes}\label{ssec:bullseye}

The theorems of this paper are not the only implications into the
conjecture. The literature supplies others, and a frontier computed without
them overstates what is needed. This subsection proves those implications, or
quotes them where they are theorems of the literature, so that the
computation of \Cref{ssec:closure} runs over both. Nothing in it is new as
mathematics except the assembly; what is new is that, once assembled, it
changes the answer.

\begin{definition}[Three statements about all varieties]\label{def:lmv}
Here a variety is smooth, projective and defined over $\CC$, except the base
$S$ in \textup{(V)}.
\begin{enumerate}[leftmargin=3.0em,itemsep=3pt]
\item[\textup{(L)}] $B(X)$ holds for every variety $X$ (\Cref{def:Bconj}).
\item[\textup{(M)}] Every Hodge class on every variety $X$ is
  \emph{motivated} in the sense of Andr\'e \cite[\S2.1]{And96b}: it equals
  $\mathrm{pr}_{*}(\alpha\cup{*}\beta)$ for some variety $Y$, algebraic
  classes $\alpha,\beta$ on $X\times Y$, the projection
  $\mathrm{pr}\colon X\times Y\to X$, and the Lefschetz involution ${*}$ of
  $H^{*}(X\times Y,\QQ)$ with respect to some ample class.
\item[\textup{(V)}] For every smooth connected complex variety $S$, every
  smooth projective morphism $f\colon\cX\to S$ and every global section $\xi$
  of $R^{2p}f_{*}\QQ$ over $S^{\mathrm{an}}$: if $\xi_{s}$ is algebraic at one
  point $s\in S$, then $\xi_{t}$ is algebraic at every point $t\in S$.
\end{enumerate}
The Lefschetz involution of a variety $Z$ of dimension $e$ with ample class
$L$ sends $L^{r}x$, for $x$ primitive of degree $j$ and $0\le r\le e-j$, to
$\pm L^{e-j-r}x$, with a sign that depends only on $j$ \cite{And96b}; the
sign plays no role below.
\end{definition}

The statement (V) is the variational Hodge conjecture for algebraic classes;
it asks nothing about Hodge classes, only that algebraicity, once present,
survives deformation.

\begin{lemma}[Kleiman]\label{lem:lefinvolution}
Let $Z$ be a variety of dimension $e$ with ample class $L$, and suppose that
$B(Z)$ holds. Then the K\"unneth projectors $\pi_{k}$ of $H^{*}(Z,\QQ)$, the
projectors onto the summands $L^{r}P^{j}$ of the Lefschetz decomposition, and
the Lefschetz involution ${*}$ are all induced by algebraic cycles on
$Z\times Z$.
\end{lemma}

\begin{proof}
Operators induced by algebraic cycles on $Z\times Z$ are closed under
composition, sums and rational multiples, because composition of
correspondences is \cite{Kle68}. The operator $L$ is induced by the image of
a divisor of class $L$ under the diagonal, and $\Lambda$ by $B(Z)$. With
$[L,\Lambda]=H$ and $H$ acting on $H^{k}$ by $k-e$ (the proof of
\Cref{lem:LambdaL}), $H=L\Lambda-\Lambda L$ is induced by an algebraic cycle,
and since the numbers $k-e$ for $0\le k\le 2e$ are distinct,
\[
  \pi_{k}=\prod_{i\ne k}\frac{H-(i-e)}{k-i}
\]
is as well. Put $\Omega=H^{2}+2H+4\Lambda L$. For $x$ primitive of degree
$j$ put $m=e-j\ge0$. Then $\Lambda x=0$ and $Hx=-mx$, and the relation
\begin{equation}\label{eq:lambdapower}
  \Lambda L^{s}x=s(m-s+1)\,L^{s-1}x\qquad(s\ge1)
\end{equation}
follows by induction from $\Lambda L=L\Lambda-H$: the case $s=1$ is
\Cref{lem:LambdaL}, and
$\Lambda L^{s}x=L\Lambda L^{s-1}x-HL^{s-1}x
=\bigl((s-1)(m-s+2)-(2s-2-m)\bigr)L^{s-1}x=s(m-s+1)L^{s-1}x$.
So $\Omega x=(m^{2}-2m+4m)x=m(m+2)x$, and $\Omega$ commutes with $L$, since
$[L,H]=-2L$ and $[L,\Lambda L]=HL$ give $[L,\Omega]=-2(LH+HL)-4L+4HL=0$;
hence $\Omega$ acts on $L^{r}P^{j}$ by $m(m+2)$. Inside $H^{k}$ the summands
are $L^{r}P^{k-2r}$, with $m=e-k+2r$, and $m\mapsto m(m+2)$ is injective on
$m\ge0$, so Lagrange interpolation in $\Omega$ composed with $\pi_{k}$ gives
the projector $p_{k,r}$ onto $L^{r}P^{k-2r}$ as an algebraic operator.
Finally let $x\in P^{j}$ and $0\le r\le m$. If $2r\le m$, then
$L^{m-2r}(L^{r}x)=L^{m-r}x$. If $2r>m$, put $t=2r-m$, so $1\le t\le r$;
by \eqref{eq:lambdapower},
\[
  \Lambda^{t}L^{r}x=\prod_{i=0}^{t-1}(r-i)(m-r+i+1)\;L^{m-r}x,
\]
and every factor is at least $1$, because $r-i\ge r-t+1\ge1$ and
$m-r+i+1\ge m-r+1\ge1$. So on $L^{r}P^{j}$ the involution is $\pm L^{m-2r}$
or a nonzero rational multiple of $\Lambda^{2r-m}$, and
${*}=\sum_{k,r}c_{k,r}T_{k,r}\,p_{k,r}$ with rational $c_{k,r}$ and
$T_{k,r}$ a power of $L$ or of $\Lambda$. Every term is algebraic.
\end{proof}

\begin{proposition}[The conjecture is (L) and (M)]\label{prop:lefmot}
The Hodge conjecture for every variety is equivalent to \textup{(L)} and
\textup{(M)} together.
\end{proposition}

\begin{proof}
Assume the conjecture. Then \textup{(L)} holds by \Cref{thm:equivalence}. An
algebraic class $\xi$ on $X$ of dimension $d$ is motivated: take $Y$ a point,
$\alpha=\xi$, and $\beta=\varepsilon h^{d}$ with $h$ the ample class and
$\varepsilon=\pm1$ the sign for which ${*}(h^{d})=\varepsilon$, which exists
because the unit $1\in H^{0}$ is primitive of degree $0$; then
$\alpha\cup{*}\beta=\xi$. Every Hodge class is algebraic, so \textup{(M)}
holds.

Conversely assume \textup{(L)} and \textup{(M)}, and let $\xi$ be a Hodge
class on $X$. By \textup{(M)}, $\xi=\mathrm{pr}_{*}(\alpha\cup{*}\beta)$ with
$\alpha,\beta$ algebraic on $X\times Y$. By \textup{(L)}, $B(X\times Y)$
holds, so ${*}$ is induced by an algebraic cycle
(\Cref{lem:lefinvolution}); hence ${*}\beta$, the product
$\alpha\cup{*}\beta$ and its push-forward $\xi$ are algebraic. This is
Andr\'e's observation that motivated classes are algebraic as soon as the
Lefschetz involutions are \cite[\S2.1]{And96b}.
\end{proof}

\begin{proposition}[Two sources of (V)]\label{prop:lefvhc}
The Hodge conjecture implies \textup{(V)}, and so does \textup{(L)}.
\end{proposition}

\begin{proof}
Let $f\colon\cX\to S$, $\xi$ and $s$ be as in \textup{(V)}, with $\xi_{s}$
algebraic. As $S$ is smooth and connected it is irreducible, so two affine
open subsets meet, and it suffices to treat $S$ affine, joining $s$ to an
arbitrary $t$ through a point of the intersection of an affine neighbourhood
of each. Then $\cX$ is quasi-projective and, by resolution of singularities,
has a smooth projective compactification $\overline{\cX}$.

Assume the conjecture. Write $n=2p$ and
$\rho\colon H^{n}(\overline{\cX},\QQ)\to H^{0}(S^{\mathrm{an}},R^{n}f_{*}\QQ)$
for restriction; it is surjective by the theorem of the fixed part
\cite[4.1.1(ii)]{Del71}. For $t\in S$ the restriction
$r_{t}\colon H^{n}(\overline{\cX},\QQ)\to H^{n}(\cX_{t},\QQ)$ is $\rho$
followed by evaluation at $t$, and evaluation is injective because $S$ is
connected; so $K=\ker r_{t}=\ker\rho$ does not depend on $t$. It is a
sub-Hodge structure, and the Hodge structure $H^{n}(\overline{\cX},\QQ)$ is
polarisable, hence semisimple, so $H^{n}(\overline{\cX},\QQ)=K\oplus K'$
with $K'$ a sub-Hodge structure. Choose $\eta\in K'$ with $\rho(\eta)=\xi$.
Then $r_{s}|_{K'}$ is an injective morphism of Hodge structures and
$r_{s}(\eta)=\xi_{s}$ is of type $(p,p)$, so every component of $\eta$ of
type other than $(p,p)$ maps to zero and is zero: $\eta$ is a Hodge class.
So $\xi_{t}=r_{t}(\eta)$ is a Hodge class on $\cX_{t}$ for every $t$, and
algebraic by the conjecture.

Assume \textup{(L)}. The algebraic class $\xi_{s}$ is motivated (proof of
\Cref{prop:lefmot}), so $\xi_{t}$ is motivated for every $t$ by Andr\'e's
deformation theorem \cite[Th\'eor\`eme~0.5]{And96b}, and a motivated class is
algebraic under \textup{(L)} (proof of \Cref{prop:lefmot}).
\end{proof}

\begin{theorem}[Deligne, Andr\'e; Abdulali, Andr\'e]\label{thm:vhcab}
\textup{(V)} implies the Hodge conjecture for every abelian variety, and so
does \textup{(L)}.
\end{theorem}

\begin{proof}
Milne calls a class \emph{accessible} if it lies in the least family of
rational classes that contains the classes of algebraic cycles, is stable
under pull-back by morphisms of varieties, and contains $t_{t}$ for every $t$
whenever $(t_{s})_{s\in S}$ is a global section of $R^{r}f_{*}\QQ$ for a
smooth projective family over a smooth connected base and contains $t_{s}$
for one $s$; and he proves, following Deligne and Andr\'e, that every Hodge
class on an abelian variety is accessible \cite[\S5.1, Theorem~3]{Mil20}.
The algebraic classes contain the classes of algebraic cycles and are stable
under pull-back, and under \textup{(V)} they have the third property. So they
contain every accessible class, hence every Hodge class on an abelian
variety; this is \cite[Remark~2]{Mil20}. The statement for \textup{(L)}
follows by \Cref{prop:lefvhc}, and is the theorem of Abdulali and Andr\'e in
the form of \cite[Theorem~4]{Mil20}.
\end{proof}

\begin{remark}[What is known of the three]\label{rem:lmvknown}
Each of \textup{(L)}, \textup{(M)} and \textup{(V)} is open, and none is
proved here in general. $B(X)$ is known for curves, where it is immediate,
and for abelian varieties \cite[\S2]{Kle68}, and \Cref{cor:mumfordlefschetz}
below adds the total space of every Mumford family. \textup{(M)} holds for the Hodge
classes on abelian varieties \cite[Th\'eor\`eme~0.6.2]{And96b}.
\textup{(V)} holds in degree $2$, where by the argument of
\Cref{prop:lefvhc} the class $\xi_{t}$ is a Hodge class and the Lefschetz
theorem on $(1,1)$ classes applies; the one general mechanism for it in
higher degree is the semiregularity theorem \cite{Blo72,BF03,Pri24}, which is the
mechanism of \Cref{thm:semiregclosure}. \Cref{thm:vhcab} says what is at
stake for this paper: the single statement \textup{(V)}, which the conjecture
itself implies, gives the Weil classes of every CM field and the classes
beyond the Weil lines at once, by a route that uses none of the theorems
proved here.
\end{remark}

\subsection{A new case of the Lefschetz standard conjecture}
\label{ssec:newlefschetz}

\textup{(L)} is known for abelian varieties and open for the varieties fibred
in them. This subsection proves it for the total space of an abelian scheme
over a curve whose invariant cycles are algebraic, and so for the total space
of every Mumford family. The proof assembles the operator $\Lambda$ from a
relative part, which is a Pontryagin product and is algebraic on every abelian
scheme, and a part along the base, which is algebraic exactly when the
invariant cycles are.

\begin{proposition}[$\Lambda$ as a Pontryagin product]
\label{prop:pontryaginlambda}
Let $A$ be an abelian variety of dimension $g$, $\ell\in H^{2}(A,\QQ)$ the
class of an ample line bundle, $D=\int_{A}\ell^{g}/g!$ and
$\gamma=\ell^{g-1}/(g-1)!$. Let $\mu\colon A\times A\to A$ be the addition
and $x\star y=\mu_{*}(p_{1}^{*}x\cup p_{2}^{*}y)$ the Pontryagin product.
Then the operator $\Lambda$ of \Cref{def:Bconj} for the class $\ell$ is
\[
  \Lambda(x)=D^{-1}\,x\star\gamma .
\]
In particular $\Lambda$ is induced by the algebraic cycle
$D^{-1}(p_{1},\mu)_{*}(p_{2}^{*}Z)$ on $A\times A$, for any cycle $Z$ of
class $\gamma$.
\end{proposition}

\begin{proof}
Put $V=H_{1}(A,\QQ)$, so that $H^{*}(A,\QQ)=\bigwedge^{*}V^{\vee}$ under cup
product and $H_{*}(A,\QQ)=\bigwedge^{*}V$ under the Pontryagin product of
homology, and let $\mathrm{PD}(x)=x\cap[A]$. Cap product makes
$H_{*}$ a module over $H^{*}$, and for a torus it is contraction; so
$\mathrm{PD}(\ell\cup x)=\ell\cap\mathrm{PD}(x)$ is the contraction
$\iota_{\ell}$ applied to $\mathrm{PD}(x)$. Poincar\'e duality commutes
with $\mu_{*}$ and carries $p_{1}^{*}x\cup p_{2}^{*}y$ to the cross product
of $\mathrm{PD}(x)$ and $\mathrm{PD}(y)$, without sign when $y$ has even
degree, so $\mathrm{PD}(x\star\gamma)=\mathrm{PD}(x)\cdot\beta$ with
$\beta=\mathrm{PD}(\gamma)\in\bigwedge^{2}V$.

Choose a basis $x_{1},y_{1},\dots,x_{g},y_{g}$ of $V^{\vee}$ with
$\ell=\sum_{i}d_{i}\,x_{i}\wedge y_{i}$, $d_{i}>0$, oriented so that
$\int_{A}x_{1}\wedge y_{1}\wedge\dots\wedge x_{g}\wedge y_{g}=1$; then
$\int_{A}\ell^{g}=g!\,d_{1}\cdots d_{g}$, so $D=d_{1}\cdots d_{g}$. Let
$e_{i},f_{i}$ be the dual basis of $V$. Then
$\gamma=\sum_{j}\bigl(\prod_{i\ne j}d_{i}\bigr)\prod_{i\ne j}x_{i}\wedge y_{i}$
and $\beta=\varepsilon D\sum_{j}d_{j}^{-1}e_{j}\wedge f_{j}$ with
$\varepsilon=\pm1$ independent of $j$, since the permutations of the pairs
preserve the orientation and exchange the terms. The algebra $\bigwedge^{*}V$
is the graded tensor product of the algebras
$\bigwedge^{*}\langle e_{i},f_{i}\rangle$, and $\iota_{\ell}$ and
multiplication by $D^{-1}\beta$ are sums over $i$ of even operators on the
$i$-th factor, namely $d_{i}\,\iota_{x_{i}\wedge y_{i}}$ and
$\varepsilon d_{i}^{-1}e_{i}\wedge f_{i}$. Operators on different factors
commute, and on the $i$-th factor, with basis $1,e_{i},f_{i},e_{i}\wedge
f_{i}$ and $\iota_{x_{i}\wedge y_{i}}(e_{i}\wedge f_{i})=\sigma$ for a sign
$\sigma$ fixed by the convention for contraction, the commutator
$[\iota_{\ell},D^{-1}\beta]$ acts by $\varepsilon\sigma$ on $1$, by $0$ on
$e_{i}$ and $f_{i}$, and by $-\varepsilon\sigma$ on $e_{i}\wedge f_{i}$, that
is by $\varepsilon\sigma(1-k_{i})$ in degree $k_{i}$. Summing over $i$,
$[\iota_{\ell},D^{-1}\beta]=\varepsilon\sigma(g-k)$ on $\bigwedge^{k}V$. As
$\mathrm{PD}$ carries $H^{q}$ to $\bigwedge^{2g-q}V$, the operator
$\Lambda'(x)=D^{-1}x\star\gamma$ satisfies
$[L,\Lambda']=\varepsilon\sigma(q-g)=\varepsilon\sigma H$ on $H^{q}(A)$.

The sign is fixed by one value. For $x$ and $y$ of complementary degrees,
$x\star y\in H^{0}(A)$ is $\int_{A\times A}p_{1}^{*}x\cup p_{2}^{*}y\cup
\mu^{*}[\mathrm{pt}]$, and $\mu^{*}[\mathrm{pt}]$ is the class of the
antidiagonal $\{(a,-a)\}$, so $x\star y=\int_{A}x\cup[-1]^{*}y$; for
$y=\gamma$, of even degree, this is $\int_{A}x\cup\gamma$. Hence
$\Lambda'(\ell)=D^{-1}\int_{A}\ell\cup\gamma=g$, while
$\Lambda(\ell)=\Lambda(L\cdot1)=g$ by \Cref{lem:LambdaL}. If
$\varepsilon\sigma=-1$, then $-\Lambda'$ would be the operator of
\Cref{def:Bconj} and $\Lambda'(\ell)=-g$. So $\varepsilon\sigma=1$, and
$\Lambda'=\Lambda$ by the uniqueness in \Cref{def:Bconj}. Finally, the
correspondence $(p_{1},\mu)_{*}(p_{2}^{*}Z)$ acts by
$x\mapsto p_{2*}\bigl(p_{1}^{*}x\cup(p_{1},\mu)_{*}p_{2}^{*}[Z]\bigr)
=\mu_{*}(p_{1}^{*}x\cup p_{2}^{*}\gamma)$, by the projection formula.
\end{proof}

\begin{theorem}[The Lefschetz operator of an abelian scheme over a curve;
cf.\ Tankeev \cite{Tan03}]
\label{thm:lefschetzfamily}
Let $C$ be a smooth projective connected curve and $\pi\colon W\to C$ an
abelian scheme of relative dimension $g$ with $W$ projective. Suppose that
every class in $H^{*}(W_{c},\QQ)$ invariant under the monodromy of
$R^{*}\pi_{*}\QQ$ is the restriction of an algebraic class on $W$. Then there
is an ample class $L_{W}$ on $W$ whose operator $\Lambda$ is induced by an
algebraic cycle on $W\times W$: $B(W)$ holds.
\end{theorem}

\begin{proof}
\emph{Step 1: the decomposition.} Let $[n]\colon W\to W$ be multiplication
by $n$ over $C$. It acts on $R^{q}\pi_{*}\QQ=\bigwedge^{q}R^{1}\pi_{*}\QQ$ by
$n^{q}$. The Leray spectral sequence degenerates \cite[4.1.1(i)]{Del71}, and
$[2]^{*}$ preserves the Leray filtration and acts on its graded pieces
$H^{p}(C,R^{q}\pi_{*}\QQ)$ by $2^{q}$. So $H^{k}(W,\QQ)=\bigoplus_{q}
H^{k}_{(q)}$, where $H^{k}_{(q)}$ is the $2^{q}$-eigenspace of $[2]^{*}$ and
maps isomorphically onto $H^{k-q}(C,R^{q}\pi_{*}\QQ)$; only $k-q\in\{0,1,2\}$
occurs. The projectors
$P_{q}=\prod_{q'\ne q}([2]^{*}-2^{q'})/(2^{q}-2^{q'})$ are induced by
algebraic cycles, polynomials in the graph of $[2]$. As $[2]^{*}$ is a ring
homomorphism, $H_{(q)}\cup H_{(q')}\subset H_{(q+q')}$; as
$\int_{W}[2]^{*}z=2^{2g}\int_{W}z$, the integral vanishes on $H_{(q)}$ for
$q\ne2g$, so Poincar\'e duality pairs $H_{(q)}$ with $H_{(2g-q)}$ only; and
$[2]_{*}=2^{2g}([2]^{*})^{-1}$ acts on $H_{(q)}$ by $2^{2g-q}$.

\emph{Step 2: the ample class.} Let $\cL$ be a relatively ample line bundle,
replaced by $\cL\otimes[-1]^{*}\cL$ so that it is symmetric. Then $[-1]^{*}$
fixes $c_{1}(\cL)$ and acts on $H_{(q)}$ by $(-1)^{q}$, so $c_{1}(\cL)$ has no
component in $H^{2}_{(1)}$, while $H^{2}_{(0)}=H^{2}(C,\QQ)$ is spanned by the
class $f$ of a fibre. So $c_{1}(\cL)=\ell+af$ with $\ell\in H^{2}_{(2)}$
algebraic and $a\in\QQ$. Put $L_{W}=\ell+mf=c_{1}(\cL)+(m-a)f$ with
$m\in\QQ$ so large that $L_{W}$ is a positive rational multiple of the class
of an ample line bundle, which is possible because $\cL$ is relatively ample.
Write $L_{\mathrm{rel}}$ and $L_{C}$ for cup product with $\ell$ and with
$mf$, so that $L_{W}=L_{\mathrm{rel}}+L_{C}$; on
$H^{p}(C,R^{q}\pi_{*}\QQ)$ the first raises $q$ by two and the second raises
$p$ by two.

\emph{Step 3: the relative operator.} Put $\gamma=\ell^{g-1}/(g-1)!$, which is
algebraic and lies in $H_{(2g-2)}$, and $D=\int_{W_{c}}\ell^{g}/g!$. Let
$\mu,\mathrm{pr}_{1},\mathrm{pr}_{2}\colon W\times_{C}W\to W$ be the addition
and the projections, and $T(x)=\mu_{*}(\mathrm{pr}_{1}^{*}x\cup
\mathrm{pr}_{2}^{*}\gamma)$. As in \Cref{prop:pontryaginlambda}, $T$ is
induced by the algebraic cycle $(\mathrm{pr}_{1},\mu)_{*}(\mathrm{pr}_{2}^{*}Z)$
on $W\times W$. Since $\pi\circ\mathrm{pr}_{1}=\pi\circ\mu$, the projection
formula gives $T(x\cup\pi^{*}b)=T(x)\cup\pi^{*}b$, so $T$ commutes with
$L_{C}$. Since $[2]\circ\mu=\mu\circ([2]\times_{C}[2])$,
$[2]_{*}T(x)=\mu_{*}(\mathrm{pr}_{1}^{*}[2]_{*}x\cup\mathrm{pr}_{2}^{*}
[2]_{*}\gamma)$, and for $x\in H_{(q)}$ this is
$2^{2g-q}\,2^{2}\,T(x)$; so $T$ maps $H_{(q)}$ into $H_{(q-2)}$. A relative
correspondence acts on the Leray spectral sequence through the maps it
induces on the local systems, so on $H^{p}(C,R^{q}\pi_{*}\QQ)$ the operator
$T$ is $H^{p}(C,T_{\bullet})$, with $T_{c}$ the Pontryagin product with
$\gamma|_{W_{c}}$ on the fibre. By \Cref{prop:pontryaginlambda}, $D^{-1}T_{c}$
is the operator $\Lambda$ of the fibre, so $\Lambda_{\mathrm{rel}}=D^{-1}T$
satisfies $[L_{\mathrm{rel}},\Lambda_{\mathrm{rel}}]=H_{\mathrm{rel}}$, where
$H_{\mathrm{rel}}$ acts on $H_{(q)}$ by $q-g$.

\emph{Step 4: the invariant classes.} Restriction maps $H^{q}(W,\QQ)$ onto the
monodromy invariants $I_{q}\subset H^{q}(W_{c},\QQ)$ \cite[4.1.1(ii)]{Del71},
and it kills $H^{q}_{(q')}$ for $q'<q$, which lies in positive Leray degree;
so it maps $H^{q}_{(q)}$ isomorphically onto $I_{q}$. Every element of
$H^{q}_{(q)}$ is algebraic: if $v\in I_{q}$ is the restriction of an algebraic
class $z$, then $P_{q}z$ is algebraic, lies in $H^{q}_{(q)}$ and restricts to
$v$. The subspace $I_{q}$ is a sub-Hodge structure, as the image of a
morphism of Hodge structures, and it is stable under cup product with
$\ell_{c}$ and under the operator $\Lambda$ of the fibre, which are the
restrictions of $L_{\mathrm{rel}}$ and $\Lambda_{\mathrm{rel}}$. The
polarisation of $H^{q}(W_{c},\QQ)$ attached to $\ell_{c}$ has the form
$(v,w)\mapsto\int_{W_{c}}v\cup S(w)$, with $S$ a polynomial in cup product
with $\ell_{c}$, the operator $\Lambda$ and the projectors of the fibre, and
its restriction to the sub-Hodge structure $I_{q}$ is again a polarisation
\cite{Voi02}. Since $S(I_{q})\subset I_{2g-q}$, the pairing
$I_{q}\times I_{2g-q}\to\QQ$ is nondegenerate.

\emph{Step 5: the operator along the base.} For $y\in H^{q}_{(q)}$ and
$w\in H^{2g-q}_{(2g-q)}$ one has $\int_{W}(y\cup f)\cup w=\int_{W_{c}}
y|_{W_{c}}\cup w|_{W_{c}}$. By Step 4 and Step 1 this shows that
$L_{C}\colon H^{q}_{(q)}\to H^{q+2}_{(q)}$ is injective, and it is onto
because $\dim H^{q+2}_{(q)}=\dim H^{2g-q}_{(2g-q)}=\dim I_{2g-q}=\dim I_{q}$,
by Poincar\'e duality on $W$ and on the invariants. Choose bases $(a_{j})$ of
$H^{q}_{(q)}$ and $(b_{k})$ of $H^{2g-q}_{(2g-q)}$, all algebraic by Step 4,
let $M_{jk}=\int_{W}a_{j}\cup mf\cup b_{k}$, which is invertible, and put
\[
  \Lambda_{C}(x)=\sum_{j,k}(M^{-1})_{kj}\Bigl(\int_{W}x\cup b_{k}\Bigr)a_{j}
  \qquad(x\in H^{q+2}(W,\QQ)),
\]
for every $q$. It is induced by the algebraic cycle
$\sum_{j,k}(M^{-1})_{kj}\,b_{k}\times a_{j}$. By Step 1, $\int_{W}x\cup b_{k}$
vanishes for $x\in H^{q+2}_{(q')}$ with $q'\ne q$, so $\Lambda_{C}$ is zero
on $H^{p}(C,R^{q}\pi_{*}\QQ)$ for $p\ne2$, and $\Lambda_{C}(L_{C}a_{i})=a_{i}$.
Hence $[L_{C},\Lambda_{C}]$ is $-1$ on $H_{(q)}^{q}$, $0$ on
$H_{(q)}^{q+1}$ and $+1$ on $H_{(q)}^{q+2}$: it is the operator $H_{C}$ that
acts on $H^{p}(C,R^{q}\pi_{*}\QQ)$ by $p-1$.

\emph{Step 6: the cross relations and the conclusion.} $[L_{C},
\Lambda_{\mathrm{rel}}]=0$ by Step 3. Both $L_{\mathrm{rel}}\Lambda_{C}$ and
$\Lambda_{C}L_{\mathrm{rel}}$ vanish off the summands with $p=2$, and on
$L_{C}y$ with $y\in H^{q}_{(q)}$ both give $\ell\cup y$, because
$\ell\cup y\in H^{q+2}_{(q+2)}$ and $L_{C}$ commutes with $L_{\mathrm{rel}}$;
so $[L_{\mathrm{rel}},\Lambda_{C}]=0$. Put
$\Lambda=\Lambda_{\mathrm{rel}}+\Lambda_{C}$, of degree $-2$. Then
\[
  [L_{W},\Lambda]=[L_{\mathrm{rel}},\Lambda_{\mathrm{rel}}]+[L_{C},\Lambda_{C}]
  =H_{\mathrm{rel}}+H_{C},
\]
which acts on $H^{p}(C,R^{q}\pi_{*}\QQ)\subset H^{p+q}(W)$ by
$(q-g)+(p-1)=(p+q)-\dim W$. So $\Lambda$ is the operator of \Cref{def:Bconj}
for $L_{W}$, by its uniqueness, and it is induced by an algebraic cycle by
Steps 3 and 5.
\end{proof}

\begin{corollary}[Mumford families]\label{cor:mumfordlefschetz}
Let $\pi\colon W\to C$ be a non-isotrivial abelian scheme of relative dimension
four over a smooth projective connected curve, $W$ projective, whose very
general fibre $X$ is a fourfold as in \Cref{prop:mumfordproduct}, with Hodge
group the group $G$ of Mumford's construction \cite{Mum69}, the restriction of
scalars from a totally real cubic field $F$ of the norm one group of a
quaternion algebra over $F$. Then $B(W)$ holds. Such families exist, over the
compact Shimura curves of \cite{Mum69}.
\end{corollary}

\begin{proof}
By Andr\'e's theorem \cite[Theorem~1]{And92} the connected monodromy group is
a normal subgroup of the derived Mumford--Tate group of $X$, which is $G$. It
is not trivial, because the family is not isotrivial, and $G$ is $\QQ$-simple,
since over $\overline{\QQ}$ its three factors are indexed by the embeddings of
$F$ and permuted transitively by the Galois group; so it is $G$. The
monodromy invariants in $H^{*}(X,\QQ)$ are therefore contained in the
invariants of $G$, which are the Hodge classes of $X$, and these are the
powers of the polarisation by \Cref{prop:mumfordproduct}(i). Item (XLII) of
\Cref{sec:verification} and the Lean certificate confirm the count: the
invariants of $\mathfrak{sl}_{2}^{3}$ in $\bigwedge^{q}(V_{1}\otimes V_{2}
\otimes V_{3})$ have dimension $1$ for $q$ even and $0$ for $q$ odd. The powers
of $\ell$ restrict to the powers of the polarisation, so every invariant class
is the restriction of an algebraic class, and \Cref{thm:lefschetzfamily}
applies.
\end{proof}

\begin{figure}[tbp]
\centering
  \includegraphics[width=0.86\textwidth]{fig_lefschetzgrid}
\caption{The proof of \Cref{thm:lefschetzfamily} for relative dimension
four. The boxes are the summands $H^{p}(C,R^{q}\pi_{*}\QQ)$ of $H^{*}(W,\QQ)$,
split off by the eigenvalue $2^{q}$ of $[2]^{*}$ shown on the left; the dotted
diagonals join the summands of one $H^{k}(W,\QQ)$, $k=p+q$, and the dashed
magenta line is the axis $q=g$ of hard Lefschetz on the fibres. The relative
operators $L_{\mathrm{rel}}$ and $\Lambda_{\mathrm{rel}}$, in teal, move
within a column; $\Lambda_{\mathrm{rel}}$ is a Pontryagin product with an
algebraic class on every abelian scheme (\Cref{prop:pontryaginlambda}). The
base operators $L_{C}$ and $\Lambda_{C}$, in ochre, move along a row between
the outer columns and kill the middle one; $\Lambda_{C}$ is the displayed
cycle built from bases $(a_{j})$, $(b_{k})$ of the invariant summands, and is
algebraic exactly when the invariant cycles are. Below each column is the
weight of $H=[L_{W},\Lambda]$ on it, and the panel at the top is Step 6 of the
proof. The ochre boxes carry the invariant cycles and their images under
$L_{C}$; for a Mumford family they are the lines listed in magenta on the
right, spanned by $\ell^{i}f$ (and by $\ell^{i}$ in the left column), and the
other boxes of the outer columns are zero (\Cref{cor:mumfordlefschetz}).}
\label{fig:lefschetzgrid}
\end{figure}

\begin{remark}[What this does and does not reach]\label{rem:newlefschetzscope}
The total space $W$ is a fivefold fibred in abelian fourfolds whose Hodge
groups are not defined by endomorphisms; it is not a product and not an
abelian variety. For abelian schemes over curves the mechanism of
\Cref{thm:lefschetzfamily} goes back to Tankeev \cite{Tan03}, who proves $B$
for the total space when certain canonically defined Hodge cycles are
algebraic; we have not checked whether his hypotheses are visibly met by
Mumford families, and make no claim of priority for
\Cref{cor:mumfordlefschetz}. The
hypothesis of \Cref{thm:lefschetzfamily} is where the difficulty sits.
Applied to the fibre square $W\times_{C}W\to C$ the theorem gives
$B(W\times_{C}W)$ as soon as every invariant class of $X\times X$ is the
restriction of an algebraic class on $W\times_{C}W$; the invariant classes
include the two exceptional classes of \Cref{prop:mumfordproduct}, which are
not known to be algebraic, so this is the smallest case of (F2) in a global
form. \Cref{thm:lefschetzfamily} does not bear on (P2), on the Kuga--Satake
class of \Cref{thm:mumfordks}, or on \textup{(L)} for varieties that are not
fibred in abelian varieties over a curve.
\end{remark}

\begin{proposition}[Powers of a Weil variety; cf.\ \cite{Mil22}]\label{prop:powers}
Let $X$ be of $(K,-1,n)$-Weil type with $n\ge2$ and $\Hg(X)=\SU(V,H)$. For
every $k\ge1$ the $\QQ$-algebra $\Hdg^{*}(X^{k})$ is generated by the divisor
classes of $X^{k}$ and the pull-backs $f_{a}^{*}\omega$ of the Weil classes
$\omega$ of $X$ along the homomorphisms $f_{a}=\sum_{j}a_{j}\,pr_{j}\colon
X^{k}\to X$, $a\in\ZZ^{k}$. Consequently, if the Weil classes of $X$ are
algebraic, the Hodge conjecture holds for every power of $X$. By
\cite{Mar25a} this is so for every such $X$ with $n=2$, and with $n=3$ and
trivial discriminant.
\end{proposition}

\begin{proof}
Write $H^{1}(X,\CC)=V_{+}\oplus V_{-}$ for the eigenspaces of $K$ and
$W=V_{+}$. The polarisation makes $V_{+}$ and $V_{-}$ isotropic and pairs them
perfectly, so $V_{-}=W^{\vee}$, and restriction to $V_{+}$ identifies
$\SU(V,H)_{\CC}$ with $\SL(W)$. As $\Hg(X^{k})$ is $\Hg(X)$ acting diagonally
on $H^{1}(X^{k},\CC)=(W\otimes U)\oplus(W^{\vee}\otimes U)$, $U=\CC^{k}$, the
Hodge classes of $X^{k}$, tensored with $\CC$, are the $\SL(W)$-invariants in
$\bigwedge^{*}\bigl((W\otimes U)\oplus(W^{\vee}\otimes U)\bigr)$. That algebra
is an $\SL(W)$-equivariant quotient of a tensor algebra, with an equivariant
section given by antisymmetrisation, so by the first fundamental theorem of
invariant theory for $\SL(W)$ its invariants are generated by the images of
the contractions and the determinant tensors: the classes
$D(u,u')=\sum_{i}(e_{i}\otimes u)\wedge(e_{i}^{\vee}\otimes u')$ and the
polarised determinants, which by polarisation are combinations of
$\det_{u}=\bigwedge_{i}(e_{i}\otimes u)$ and
$\det^{\vee}_{u}=\bigwedge_{i}(e_{i}^{\vee}\otimes u)$, $u,u'\in U$. For
$2n\ge3$ the $D(u,u')$ span the invariants of degree two, which are the
divisor classes tensored with $\CC$. With $\alpha_{\pm}$ spanning
$\bigwedge^{2n}V_{\pm}$ (\Cref{prop:weilline}) one has
$f_{a}^{*}\alpha_{+}=\det_{a}$ and $f_{a}^{*}\alpha_{-}=\det^{\vee}_{a}$, and
since $u\mapsto\det_{u}$ is polynomial and $\ZZ^{k}$ is Zariski dense in
$U$, these span the same space as all $\det_{u}$. So the $\QQ$-algebra
generated by $\NS(X^{k})_{\QQ}$ and the $f_{a}^{*}\omega$ is contained in
$\Hdg^{*}(X^{k})$ and has the same complexification. The generators are
algebraic when the Weil classes of $X$ are, and they are for $n=2$ and for
$n=3$ at trivial discriminant by \cite{Mar25a}.
\end{proof}

The statement for powers is known (compare \cite{Mil22}, and \cite{Li26} for
the powers of CM fourfolds); we include the short proof because it is what
the next corollary uses.

\begin{corollary}[Abelian schemes of small relative dimension]
\label{cor:lefschetzsmall}
Let $\pi\colon W\to C$ be an abelian scheme of relative dimension $g$ over a
smooth projective connected curve, with $W$ projective, and suppose that
every monodromy invariant class in $H^{*}(W_{c},\QQ)$ is a Hodge class, as
happens when the identity component of the Zariski closure of the monodromy
group is the Hodge group of a very general fibre.
\begin{enumerate}[label=\textup{(\roman*)},leftmargin=2.6em,itemsep=2pt]
\item If $g\le4$, then $B(W)$ holds.
\item If the fibres are of $(K,-1,n)$-Weil type with connected monodromy
  group $\SU(V,H)$, and $n=2$, or $n=3$ with trivial discriminant, then
  $B$ holds for every fibre power $W\times_{C}\dots\times_{C}W$. In
  particular $B(W_{C})$ of \Cref{cor:weilequiv} holds in these cases.
\end{enumerate}
\end{corollary}

\begin{proof}
By the theorem of the fixed part the invariant classes form a sub-Hodge
structure of every fibre, with a Hodge structure independent of the point, so
they are Hodge classes on every fibre. (i) The Hodge conjecture holds for
abelian varieties of dimension at most four, by the Lefschetz theorem on
$(1,1)$ classes and hard Lefschetz up to dimension three and by \cite{Mar25a}
in dimension four, so each invariant
class is algebraic on every fibre; \Cref{lem:spreading} makes it the
restriction of an algebraic class on $W$, and \Cref{thm:lefschetzfamily}
applies. (ii) The fibre power is an abelian scheme with projective total
space and the monodromy acting diagonally. Its invariant classes are
$\SU(V,H)$-invariants, hence Hodge classes at every member with
$\Hg=\SU(V,H)$, an uncountable set of points, and there the Hodge conjecture
holds for the fibre by \Cref{prop:powers}; \Cref{lem:spreading} and
\Cref{thm:lefschetzfamily} conclude as before. For $W_{C}$ the connected
monodromy group is $\SU(V,H)$ when it is Zariski dense, which is the
hypothesis of \Cref{cor:weilequiv}(ii).
\end{proof}

\begin{proposition}[The total space of a Mumford family]
\label{prop:mumfordfivefold}
Let $W\to C$ be as in \Cref{cor:mumfordlefschetz}, with $C$ the compact
Shimura curve of \cite{Mum69} of torsion-free level, or a finite \'etale cover
of it. Then the Hodge conjecture holds for the fivefold $W$: its Hodge classes
of degree $2p$ are spanned by $\ell^{p}$ and $f\ell^{p-1}$, with $f$ the class
of a fibre.
\end{proposition}

\begin{proof}
By Step 1 of the proof of \Cref{thm:lefschetzfamily},
$H^{k}(W,\QQ)=\bigoplus_{q}H^{k}_{(q)}$ with $H^{k}_{(q)}\cong
H^{k-q}(C,R^{q}\pi_{*}\QQ)$, the summands being sub-Hodge structures cut out by
algebraic projectors. The summand with $k=q$ is the space of invariants, which
by the proof of \Cref{cor:mumfordlefschetz} is spanned by the powers of the
polarisation, restrictions of the powers of $\ell$. The summand with $k=q+2$
is $f\cup H^{q}_{(q)}$ by Step 5, spanned by the $f\ell^{q/2}$. It remains to
see that $H^{1}(C,R^{q}\pi_{*}\QQ)$ carries no Hodge class. For $q$ even its
weight $q+1$ is odd. For $q$ odd, over $\CC$ the local system splits as a sum
of $S^{k_{1}}V_{1}\otimes U_{\lambda}\otimes M_{\lambda}$, where $V_{1}$ is the
uniformising variation of the factor of $G$ split at the real place, $U_{\lambda}$ is
a unitary variation of type $(0,0)$, $M_{\lambda}$ is constant of type $(a,a)$
with $2a+k_{1}=q$, and $k_{1}\equiv q\pmod2$ because the central element
$(-1,1,1)$ acts on $V$ by $-1$. When $C$ is \'etale over the Shimura curve the
Kodaira--Spencer map of $V_{1}$ is an isomorphism, and then
$H^{1}(C,S^{k_{1}}V_{1}\otimes U_{\lambda})$ has only the types $(k_{1}+1,0)$
and $(0,k_{1}+1)$ \cite{Zuc79}; after the twist by $M_{\lambda}$ these are
never balanced, since $k_{1}\ge1$.
\end{proof}

We have not found \Cref{prop:mumfordfivefold} in the literature, but it is of a
kind studied for Kuga fibre varieties and we make no claim of priority. The
\'etale hypothesis cannot be dropped from the argument: over a ramified cover
the Kodaira--Spencer map has zeros and balanced types can occur in
$H^{1}(C,R^{q}\pi_{*}\QQ)$.

\subsection{The two remaining targets, each reduced to one case of (L)}
\label{ssec:targets}

The machinery of \Cref{ssec:newlefschetz} does more than add a case of
\textup{(L)}. For an abelian scheme over a curve it shows that \textup{(L)}
for the total space is exactly the statement that algebraicity propagates
along the curve, and the base points of this paper then turn each of the two
targets left open in \Cref{rem:bullseye}, one Weil family and the first
classes beyond the Weil lines, into \textup{(L)} for one explicit variety. The
reductions are equivalences, so they locate the difficulty rather than remove
it. Along the way two things are proved outright: the Hodge conjecture for the
self-product of a Mumford fourfold at every CM point, and the impossibility of
realising the two branches of \Cref{rem:twobranches} by two objects.

\begin{lemma}[Spreading along a curve]\label{lem:spreading}
Let $\pi\colon W\to C$ be a smooth projective morphism onto a smooth projective
connected curve, $W$ projective, and $\xi$ a global section of
$R^{2p}\pi_{*}\QQ$. If $\xi_{t}$ is algebraic for uncountably many $t\in C$,
there are an algebraic class $\tilde\xi$ on $W$ and an integer $d\ge1$ with
$\tilde\xi|_{W_{t}}=d\,\xi_{t}$ for every $t\in C$.
\end{lemma}

\begin{proof}
The proof of \Cref{thm:closedlocus} applies verbatim with $(\cA,\cS,\omega)$
replaced by $(W,C,\xi)$: the set of $t$ at which $\xi_{t}$ is algebraic is
the union, over countably many tuples of Hilbert polynomials and rational
coefficients, of the images of unions $T$ of connected components of fibre
products of relative Hilbert schemes, projective over $C$, on which
$\sum_{i}q_{i}[Z_{i}]=\xi$. Uncountably many $t$ force one image to be
infinite, hence equal to $C$. Choose an irreducible component of that $T$
surjecting onto $C$, two of its points over distinct points of $C$, and an
irreducible curve through them; its normalisation $C'$ maps finitely onto $C$,
of some degree $d$. Pulling the universal subschemes back to
$W'=W\times_{C}C'$ gives a cycle $z=\sum_{i}q_{i}[\cZ_{i}]$ whose restriction
to $W'_{c'}=W_{g(c')}$ is $\xi_{g(c')}$ for every $c'\in C'$. The projection
$h\colon W'\to W$ is finite and flat of degree $d$, and by base change
$h_{*}[z]$ restricts on $W_{t}$ to the sum of the restrictions over the
points of $C'$ above $t$, counted with multiplicity, which is $d\,\xi_{t}$.
\end{proof}

\begin{theorem}[\textup{(L)} for an abelian scheme over a curve is
propagation]\label{thm:lefschetzpropagation}
Let $\pi\colon W\to C$ be as in \Cref{thm:lefschetzfamily}, but with no
hypothesis on its invariant cycles, and keep the notation of its proof. Let
$\Lambda_{C}$ be the operator that is zero on $H^{p}(C,R^{q}\pi_{*}\QQ)$ for
$p\ne2$ and inverts $L_{C}\colon H^{q}_{(q)}\to H^{q+2}_{(q)}$.
\begin{enumerate}[label=\textup{(\roman*)},leftmargin=2.6em,itemsep=2pt]
\item The operator $\Lambda$ of \Cref{def:Bconj} for $L_{W}$ is
  $\Lambda_{\mathrm{rel}}+\Lambda_{C}$. So $B(W)$ holds for $L_{W}$ if and
  only if $\Lambda_{C}$ is induced by an algebraic cycle.
\item If $B(W)$ holds for $L_{W}$, then for every $c\in C$ and every algebraic
  class $y$ on $W_{c}$ the class $\Lambda_{C}(i_{c*}y)$ is algebraic; up to
  the factor $m$ it lies in $H^{0}(C,R^{q}\pi_{*}\QQ)$ and restricts on
  $W_{c}$ to the unique monodromy invariant class $v$ with
  $\int_{W_{c}}v\cup w=\int_{W_{c}}y\cup w$ for every invariant $w$ of
  complementary degree. In particular every global section of
  $R^{2p}\pi_{*}\QQ$ that is algebraic at one point is the restriction of an
  algebraic class on $W$, and is algebraic at every point.
\item If at one point $c_{0}$ every monodromy invariant class of
  $H^{*}(W_{c_{0}},\QQ)$ is algebraic, then $B(W)$ holds for $L_{W}$ if and
  only if every monodromy invariant class is the restriction of an algebraic
  class on $W$.
\end{enumerate}
\end{theorem}

\begin{proof}
(i) Steps 4 and 5 of the proof of \Cref{thm:lefschetzfamily} show that
$L_{C}\colon H^{q}_{(q)}\to H^{q+2}_{(q)}$ is an isomorphism without using
the hypothesis on invariant cycles: the nondegeneracy of
$I_{q}\times I_{2g-q}\to\QQ$ comes from the Hodge--Riemann relations alone.
So $\Lambda_{C}$ is defined, and Steps 5 and 6, which use only that
$\Lambda_{C}$ vanishes off $p=2$ and inverts $L_{C}$, give
$[L_{W},\Lambda_{\mathrm{rel}}+\Lambda_{C}]=H$. By the uniqueness in
\Cref{def:Bconj} this is $\Lambda$. As $\Lambda_{\mathrm{rel}}$ is induced by
an algebraic cycle (Step 3), $\Lambda$ is if and only if $\Lambda_{C}$ is.

(ii) The Gysin map $i_{c*}$ commutes with $[2]_{*}$, which acts on
$H^{q}(W_{c})$ by $2^{2g-q}$, so $i_{c*}y\in H^{q+2}_{(q)}$, the summand with
$p=2$. If $y$ is algebraic so is $i_{c*}y$, and by (i) so is
$\Lambda_{C}(i_{c*}y)$. Put $u=m\,\Lambda_{C}(i_{c*}y)\in H^{q}_{(q)}$, so
$f\cup u=i_{c*}y$. For $w\in H^{2g-q}_{(2g-q)}$ the projection formula gives
$\int_{W_{c}}u|_{W_{c}}\cup w|_{W_{c}}=\int_{W}(f\cup u)\cup w
=\int_{W}i_{c*}y\cup w=\int_{W_{c}}y\cup w|_{W_{c}}$, and $w|_{W_{c}}$ runs
over the invariants of complementary degree (Step 4); by the nondegeneracy of
Step 4 this determines $u|_{W_{c}}$. If $y=\xi_{c}$ for a global section
$\xi$, then $v=\xi_{c}$ has the property, so $u|_{W_{c}}=\xi_{c}$, and $u$,
being a global section of $R^{q}\pi_{*}\QQ$ equal to $\xi$ at $c$, restricts
to $\xi_{t}$ at every $t$.

(iii) If every invariant class is a restriction of an algebraic class, $B(W)$
holds by \Cref{thm:lefschetzfamily}. Conversely, under $B(W)$, (ii) applied at
$c_{0}$ to each invariant class makes it the restriction of an algebraic
class.
\end{proof}

\begin{figure}[tbp]
\centering
  \includegraphics[width=0.84\textwidth]{fig_propagation}
\caption{\Cref{thm:lefschetzpropagation}(ii) along a compact Shimura curve
$C$. An algebraic class $y$ on one fibre, at a base point of a Weil family or
at a CM point of a Mumford family (hollow dots: the CM points, countable and
dense), is pushed into $W$ by the Gysin map and brought back to the invariant
summand by the base operator $\Lambda_{C}$. When $B(W)$ holds the result $u$
is an algebraic class on $W$, pinned down by the displayed pairing against the
invariant classes, and its restriction to a very general fibre is the
invariant part of $y$. The dashed arcs are isogenies between CM members along
Hecke orbits: they carry algebraic classes to algebraic classes but reach
only countably many points, and \Cref{lem:spreading}, in the box, needs
uncountably many. This is how \Cref{cor:weilequiv} and
\Cref{cor:mumfordequiv} turn each target into one case of \textup{(L)}.}
\label{fig:propagation}
\end{figure}

\begin{theorem}[The Mumford square at a CM point]\label{thm:mumfordcm}
Let $W\to C$ be a family as in \Cref{cor:mumfordlefschetz} and $c\in C$ a CM
point. Then every Hodge class on $X_{c}\times X_{c}$ is algebraic. In
particular the two exceptional classes of \Cref{prop:mumfordproduct} are
algebraic on $X_{c}\times X_{c}$.
\end{theorem}

\begin{proof}
The Hodge group $T$ of $X_{c}$ is a $\QQ$-torus contained in $G$. It lies in a
maximal $\QQ$-torus of $G=\Res_{F/\QQ}\SL_{1}(D)$, which is the restriction of
scalars of the norm one torus of a quadratic extension of $F$ inside $D$, of
rank three, one factor for each embedding of $F$. The cocharacter of the Hodge
structure lies in the factor of the embedding at which $D$ splits, and the
Galois group permutes the embeddings of $F$ transitively, so the smallest
$\QQ$-subtorus containing it has rank three: $T_{\CC}$ is a maximal torus of
$G_{\CC}\cong\SL_{2}^{3}$. Its eight weights on
$V=H^{1}(X_{c},\CC)=V_{1}\otimes V_{2}\otimes V_{3}$ are the vectors
$(\pm1,\pm1,\pm1)$, pairwise distinct. Hence
$\End^{0}(X_{c})\otimes\CC=\End_{T}(V)$ is the algebra of diagonal matrices
in a weight basis, $\operatorname{Hom}(X_{c}\times X_{c},X_{c})\otimes\CC$ is
the space of pairs $(\varphi,\psi)$ of such matrices acting by
$e_{w}\mapsto\varphi_{w}e_{w}^{(1)}+\psi_{w}e_{w}^{(2)}$, and the Hodge
classes of $X_{c}$ and of $X_{c}\times X_{c}$, tensored with $\CC$, are
spanned by the monomials of weight zero.

Item (XLIII) of \Cref{sec:verification} shows that the ring generated by the
divisor classes of $X_{c}\times X_{c}$ and by the pull-backs $f^{*}(t)$, for
$f=(\varphi,\psi)$ and $t$ a Hodge class of $X_{c}$ of degree $4$ or $6$,
contains every Hodge class of $X_{c}\times X_{c}$, in every degree. The map
$f\mapsto f^{*}(t)$ is polynomial and the rational homomorphisms are Zariski
dense in their complexification, so the same span is reached with $f$ defined
over $\QQ$ and $t$ rational. Divisor classes are algebraic, and the Hodge
classes of the abelian fourfold $X_{c}$ are algebraic because the Hodge
conjecture holds for abelian fourfolds \cite{Mar25a}; pull-backs and products
of algebraic classes are algebraic.
\end{proof}

\begin{corollary}[The first classes beyond the Weil lines are one case of
\textup{(L)}]\label{cor:mumfordequiv}
Let $W\to C$ be as in \Cref{cor:mumfordlefschetz} and
$W^{(2)}=W\times_{C}W\to C$, an abelian scheme of relative dimension eight
with $W^{(2)}$ projective of dimension nine. The following are equivalent.
\begin{enumerate}[label=\textup{(\alph*)},leftmargin=2.6em,itemsep=2pt]
\item $B(W^{(2)})$ holds for the class $L_{W^{(2)}}$.
\item The two exceptional classes of $X_{t}\times X_{t}$ are algebraic for
  uncountably many $t\in C$.
\item The Hodge conjecture holds for $X_{t}\times X_{t}$ for every $t\in C$.
\item The Hodge conjecture holds for $W^{(2)}$; for this item $C$ is taken to
  be the compact Shimura curve of torsion-free level, or a finite \'etale
  cover of it.
\end{enumerate}
\end{corollary}

\begin{proof}
As in the proof of \Cref{cor:mumfordlefschetz}, the connected monodromy group
is $G$, so the monodromy invariants in $H^{*}(X_{t}\times X_{t},\QQ)$ are
contained in the $G$-invariants, which are the Hodge classes of
$X_{t}\times X_{t}$ for $t$ very general; conversely a $G$-invariant class is
flat, being fixed by the monodromy. Item (XLIII) shows that these invariants
form a ring generated by its elements of degree two, the divisor classes, and
of degree four, which are the six products of divisor classes and the two
exceptional classes.

(c) implies (b) trivially. (b) implies (a): the exceptional classes are
restrictions of algebraic classes on $W^{(2)}$ by \Cref{lem:spreading}; so are
the invariant divisor classes, which are algebraic on every fibre, by the
same lemma; hence so is every invariant class, a polynomial in these, and
\Cref{thm:lefschetzfamily} gives $B(W^{(2)})$. (a) implies (c): at a CM point
$c_{0}$ every Hodge class of $X_{c_{0}}\times X_{c_{0}}$ is algebraic by
\Cref{thm:mumfordcm}, so by \Cref{thm:lefschetzpropagation}(iii) every
invariant class is the restriction of an algebraic class, and is algebraic on
every fibre. At a CM point (c) is \Cref{thm:mumfordcm}. At any other point
$t$ the Hodge group of $X_{t}$ is $G$, so its Hodge classes are the
invariant ones. Indeed, if the Hodge group $H_{t}$ were smaller, the Hodge
classes of the powers of $X_{t}$ that are not $G$-invariant would stay Hodge
only on a discrete subset of the period domain, the upper half plane; every
point of the orbit of the Hodge structure at $t$ under $H_{t}(\RR)^{+}$
carries them, so that orbit is a point, the cocharacter of
the Hodge structure is central in $H_{t}$, and $H_{t}$, the smallest
$\QQ$-group containing it, is a torus: $t$ would be a CM point.

For (d), split $H^{*}(W^{(2)},\QQ)$ into the summands $H^{p}(C,R^{q})$,
$p\in\{0,1,2\}$, as in the proof of \Cref{prop:mumfordfivefold}. The summands
with $p=1$ carry no Hodge class, by the argument given there, applied to the
fibre square. Those with $p=2$ are $f\cup H^{q}_{(q)}$, and
$f\cup u=i_{c*}(u|_{W^{(2)}_{c}})$ for any $c$; taking $c$ a CM point, where
$u|_{W^{(2)}_{c}}$ is a Hodge class and so algebraic by \Cref{thm:mumfordcm},
they are algebraic. The summands with $p=0$ are the invariants, all of Hodge
type. So (d) holds exactly when every invariant class is the restriction of an
algebraic class, which by \Cref{thm:lefschetzpropagation}(iii), applied at a
CM point, is (a).
\end{proof}

\begin{corollary}[One Weil family is one case of \textup{(L)}]
\label{cor:weilequiv}
Keep \Cref{setup:shimura}, let $s_{0}\in\cS$ be the image of a base point of
\Cref{thm:everyfamily}, and let $t_{1}\in\cS$ lie outside every proper
Zariski closed set $\Sigma_{\underline{P},\underline{q}}$ of
\Cref{thm:closedlocus}. Let $\kappa\colon C\to\cS$ be a morphism from a smooth
projective connected curve whose image contains $s_{0}$ and $t_{1}$, and
$W_{C}=\cA\times_{\cS}C$.
\begin{enumerate}[label=\textup{(\roman*)},leftmargin=2.6em,itemsep=2pt]
\item If $B(W_{C})$ holds for $L_{W_{C}}$, then $\Sigma=\cD_{n,n}$, and the
  Hodge conjecture holds in every degree for every member of the family with
  $\Hg(A)=\SU(V,H)$.
\item Conversely, if $\Sigma=\cD_{n,n}$ and the monodromy of $W_{C}$ is Zariski
  dense in $\SU(V,H)$, then $B(W_{C})$ holds for $L_{W_{C}}$.
\end{enumerate}
\end{corollary}

\begin{proof}
(i) The class $\kappa^{*}\omega$ is a global section of
$R^{2n}\pi_{*}\QQ$ on $C$, algebraic at a point over $s_{0}$ by
\Cref{thm:everyfamily}. By \Cref{thm:lefschetzpropagation}(ii) it is algebraic
at a point over $t_{1}$, so $t_{1}$ lies in some closed set
$\Sigma_{\underline{P},\underline{q}}$, which is therefore all of $\cS$.
\Cref{thm:closedlocus} gives $\Sigma=\cD_{n,n}$, and \Cref{thm:residual}(a)
the rest. (ii) With dense monodromy the invariants are the
$\SU(V,H)$-invariants, which are the Hodge classes of a member with
$\Hg(A)=\SU(V,H)$: the powers of the polarisation and the Weil classes, by
\Cref{thm:hdgcount}. All are algebraic on every fibre when
$\Sigma=\cD_{n,n}$, so they are restrictions of algebraic classes by
\Cref{lem:spreading}, and \Cref{thm:lefschetzfamily} applies.
\end{proof}

Curves $\kappa$ as in \Cref{cor:weilequiv} exist: the Baily--Borel
compactification of $\cS$ is projective \cite{BB66}, its boundary has
codimension $2n-1\ge3$, the rational boundary components of the domain of type
$\mathrm{I}_{n,n}$ being of type $\mathrm{I}_{n-r,n-r}$ with $r\ge1$, and a
general complete intersection curve through $s_{0}$ and $t_{1}$ avoids it.

\begin{theorem}[Two branches cannot be two objects]\label{thm:nosum}
Let $A$ be of $(K,-1,n)$-Weil type with $n\ge2$, $\omega\ne0$ a rational class
on the Weil line, and $E$ a perfect complex with $\ch(E)=N\omega$, $N\ne0$, and
$\sigma_{E}$ injective, for instance with $\dim\Ext^{2}(E,E)=2n(2n-1)$. Let
$E\cong\bigoplus_{i}F_{i}[k_{i}]$ be any decomposition. Then every summand
$F_{i}$ with $\ch(F_{i})\ne0$ has a Yoneda product
$\bigwedge^{2}\Ext^{1}(F_{i},F_{i})\to\Ext^{2}(F_{i},F_{i})$ that is not
injective. In particular $E$ is not a direct sum of shifts of line bundles,
of skyscraper sheaves of points, or of their images under autoequivalences of
$D^{b}(A)$, since the Ext algebra of each of these is the exterior algebra on
its $\Ext^{1}$.
\end{theorem}

\begin{proof}
By \Cref{thm:p2numerical}(ii) the kernel of $\operatorname{ev}_{E}$ in degree
two is $P\wedge Q$. The characteristic morphism is compatible with the
decomposition, so $\operatorname{ev}_{F_{i}}(p\wedge q)=
\operatorname{ev}_{F_{i}}(p)\operatorname{ev}_{F_{i}}(q)=0$ for all
$p\in P$ and $q\in Q$. Suppose the Yoneda product of $F=F_{i}$ is injective on
$\bigwedge^{2}\Ext^{1}$. Then the images $\bar P$ and $\bar Q$ of $P$ and $Q$
in $\Ext^{1}(F,F)$ satisfy $\bar p\wedge\bar q=0$ for all $\bar p,\bar q$, so
either $\bar P=0$, or $\bar Q=0$, or $\bar P=\bar Q$ is a line. In the first
case $\operatorname{ev}_{F}$ kills $P\wedge HH^{1}$, in the second
$Q\wedge HH^{1}$, in the third all of $HH^{2}=HH^{1}\wedge HH^{1}$; by the
compatibility $\sigma_{F}\circ\operatorname{ev}_{F}=(\,\cdot\,)\lrcorner
\ch(F)$ of \Cref{setup:hh} the same subspace annihilates $x=\ch(F)$.

These annihilators are computed in the free module of rank one that
$H^{*}(A,\CC)$ is over $HH^{*}(A)=\bigwedge^{*}(P\oplus Q)$, with generator the
holomorphic volume form $\Omega$: an element $y\cdot\Omega$ is killed by
$P\wedge HH^{1}$ exactly when $p\wedge y$ has top degree for every $p\in P$,
that is when $y\in\det P\wedge\bigwedge^{*}Q+\bigwedge^{\ge4n-1}(P\oplus Q)$.
Since $\det P\cdot\Omega=\pm\alpha_{+}$, the first summand gives the classes
$\bigwedge^{j}Q\cdot\alpha_{+}$, of type $(n-b,n+a)$ with $a$ wedges by
$V_{-}^{0,1}$ and $b$ contractions by $T_{+}$, and the second gives classes of
type $(1,2n)$, $(0,2n-1)$ and $(0,2n)$. The components of $x$ are of type
$(k,k)$, so $x\in\CC\alpha_{+}$. Symmetrically the second case gives
$x\in\CC\alpha_{-}$ and the third $x=0$; item (XLIII) confirms the three
annihilators for $n\le3$, and so does the Lean certificate. Finally $x$ is
rational, complex conjugation carries $V_{+}$ to $V_{-}$ and $H^{1,0}$ to
$H^{0,1}$, so $\bar\alpha_{+}\in\CC\alpha_{-}$, and a rational multiple of
$\alpha_{+}$ or of $\alpha_{-}$ is zero. So $\ch(F)=0$. As
$\sum_{i}(-1)^{k_{i}}\ch(F_{i})=N\omega\ne0$, some summand has nonzero
Chern character, and the first statement follows; the second is the special
case in which every summand has an exterior Ext algebra.
\end{proof}

\begin{remark}[Where the two targets stand]\label{rem:targetsstatus}
\Cref{cor:weilequiv} and \Cref{cor:mumfordequiv} say that closing one Weil
family and closing the first classes beyond the Weil lines are each the
Lefschetz standard conjecture for one variety: the total space of the family
over a curve through a base point, and the ninefold $W\times_{C}W$. By
\Cref{thm:lefschetzpropagation} that case of \textup{(L)} is in turn the
algebraicity of the invariant cycles, so nothing is gained for free: the
reduction moves both targets onto the bullseye of \Cref{rem:bullseye} and
shows that they sit there, and no further. Some things are unconditional:
the Hodge conjecture for the total space of a Mumford family over its Shimura
curve (\Cref{prop:mumfordfivefold}), $B$ for every abelian scheme over a curve
of relative dimension at most four whose invariant classes are Hodge
(\Cref{cor:lefschetzsmall}), and the following two. The
Mumford classes are algebraic at every CM point (\Cref{thm:mumfordcm}), so the
first case of (F2) is a propagation problem along one compact Shimura curve,
exactly as a Weil family is a propagation problem from its base point. And the
object (P2) asks for cannot be assembled from two objects each of which sees
one branch (\Cref{thm:nosum}): the two branches of \Cref{rem:twobranches} must
live inside a single summand with a quadratic relation among its first order
deformations. Neither target is closed here. For the Mumford target
\Cref{ssec:mumfordrigid} goes further: it closes off degeneration and
deformation into a larger family, reduces semiregularity to one dimension,
$119$, and specialisation to a degree bound at CM points; and
\Cref{ssec:mumfordobject} closes off every construction from line bundles and
every twistor deformation of the square, and shows that modulo $p$ the
classes are represented at every closed point of every reduction, so that the
target there is a bound on those representatives (\Cref{thm:lefschetzclosure},
\Cref{prop:notwistor}, \Cref{thm:modpdensity}).

The logic runs one way only. $B(W_{C})$ and $B(W\times_{C}W)$ are two cases of
\textup{(L)}, not \textup{(L)}: \textup{(L)} for every variety implies both,
but both together close one Weil family and the first case of (F2) and
nothing more. The conjecture needs \textup{(L)} for every variety together
with \textup{(M)} (\Cref{prop:lefmot}), or (F3), which is the conjecture
itself (\Cref{prop:f3ishc}), and no case of \textup{(L)} proved or reduced in
this section bears on \textup{(M)} or on (F3).
\end{remark}

\subsection{What would close it, computed rather than listed}
\label{ssec:closure}

A list of open problems invites a reading it does not deserve, that each entry
is a step and that the steps together are a path. The implications proved
above, together with those of \Cref{ssec:bullseye}, decide the question, so
we settle it by computation instead of by description. The logical skeleton
of the paper and of the literature it quotes is finite: a set of statements,
each of them proved here, quoted from the literature, or open, and a set of
inference rules, each of them one theorem. Consequences can then be taken.

Such a computation is only as strong as its rule set. A set of statements
the rules do not carry to the conjecture is \emph{insufficient} in the sense
of the computation, and that means no more than that no rule derives the
conjecture from it. An implication that is true and left out makes the
computation report a set as necessary when it is not. One such implication
is elementary, and it decides the shape of the answer, so we prove it first.

\begin{proposition}[The conjecture for the varieties that are not abelian
is the conjecture]\label{prop:f3ishc}
Let $A$ be a complex abelian variety of dimension $g\ge1$ and let
$X=A\times\PP^{1}$. Then $X$ is a smooth complex projective variety that is
not an abelian variety, and if every Hodge class of degree $2p$ on $X$ is
algebraic, then so is every Hodge class of degree $2p$ on $A$. Consequently
the statement
\begin{quote}
the Hodge conjecture holds for every smooth complex projective variety that
is not an abelian variety
\end{quote}
implies the Hodge conjecture for every abelian variety, and it is equivalent
to $\mathbf{HC}$.
\end{proposition}

\begin{proof}
$X$ is smooth and projective of dimension $g+1$. By the K\"unneth formula
$h^{1,0}(X)=h^{1,0}(A)+h^{1,0}(\PP^{1})=g$, while an abelian variety of
dimension $g+1$ has $h^{1,0}=g+1$; so $X$ is not an abelian variety. (Equally,
$X$ contains the rational curves $\{a\}\times\PP^{1}$, and every morphism
from $\PP^{1}$ to a complex torus is constant, since it lifts to the
universal cover $\CC^{g+1}$ \cite[Chapter~1]{BL04}.)

Let $\mathrm{pr}\colon X\to A$ be the projection and $D=A\times\{0\}$, a
smooth divisor that $\mathrm{pr}$ maps isomorphically onto $A$, so that
$\mathrm{pr}_{*}[D]=1$ in $H^{0}(A,\QQ)$. Let $\alpha$ be a Hodge class of
degree $2p$ on $A$. Pull-back along a morphism is a morphism of Hodge
structures, so $\mathrm{pr}^{*}\alpha$ is a Hodge class of degree $2p$ on
$X$, and by hypothesis $\mathrm{pr}^{*}\alpha=[Z]$ for an algebraic cycle
$Z$ on $X$ with rational coefficients. The projection formula gives
\[
  \mathrm{pr}_{*}\bigl([Z]\cup[D]\bigr)
  =\mathrm{pr}_{*}\bigl(\mathrm{pr}^{*}\alpha\cup[D]\bigr)
  =\alpha\cup\mathrm{pr}_{*}[D]=\alpha .
\]
The cycle class map is compatible with intersection products and with proper
push-forward \cite[Chapter~19]{Ful98}, so the left side
is the class of the cycle $\mathrm{pr}_{*}(Z\cdot D)$ and $\alpha$ is
algebraic. A point, the abelian variety of dimension zero, satisfies the
conjecture trivially.

So the displayed statement gives the conjecture for every abelian variety,
and with it the conjecture for every smooth complex projective variety. The
converse is immediate: the displayed statement is a special case of
$\mathbf{HC}$.
\end{proof}

Since the conjecture for the varieties that are not abelian is the whole
conjecture, the statement that belongs on a frontier is weaker: it asks for
algebraicity only modulo what abelian varieties supply.

\begin{definition}[The conjecture modulo abelian varieties]
\label{def:f3prime}
For a smooth complex projective variety $X$ and $p\ge0$ let
$\mathrm{Ab}^{p}(X)\subseteq H^{2p}(X,\QQ)$ be the span of the classes
$\Gamma_{*}\gamma=\mathrm{pr}_{X*}\bigl([\Gamma]\cup\mathrm{pr}_{B}^{*}\gamma\bigr)$
of degree $2p$, where $B$ runs over the complex abelian varieties, the point
included, $\gamma$ over the Hodge classes on $B$, and $\Gamma$ over the
algebraic cycles on $B\times X$ with rational coefficients. The statement
\textup{(F3$'$)} is: for every $X$ and every $p$,
\[
  \Hdg^{p}(X)=\mathrm{Ab}^{p}(X).
\]
\end{definition}

Each $\Gamma_{*}\gamma$ is a Hodge class, since cup product with an
algebraic class and push-forward preserve Hodge classes, so
$\mathrm{Ab}^{p}(X)\subseteq\Hdg^{p}(X)$; and taking $B$ a point shows that
$\mathrm{Ab}^{p}(X)$ contains every algebraic class. So \textup{(F3$'$)} says
that every Hodge class is algebraic modulo the images of Hodge classes of
abelian varieties under algebraic correspondences, which is the form of (F3)
that \Cref{rem:f3collapse} asks for.

\begin{proposition}[What is known of \textup{(F3$'$)}]\label{prop:f3prime}
\leavevmode
\begin{enumerate}[label=\textup{(\roman*)},leftmargin=2.6em]
\item $\mathbf{HC}$ implies \textup{(F3$'$)}, and \textup{(F3$'$)} together
  with the conjecture for abelian varieties implies $\mathbf{HC}$.
\item \textup{(F3$'$)} holds on every abelian variety and on every product
  $A\times\PP^{1}$ with $A$ abelian; the argument of \Cref{prop:f3ishc} gives
  nothing for it.
\item Let $\mathcal{A}$ be the class of smooth complex projective varieties
  $X$ for which $H^{*}(X,\QQ)$ is spanned by classes $\Gamma_{*}y$ with $B$
  abelian, $\Gamma$ an algebraic cycle on $B\times X$ and
  $y\in H^{*}(B,\QQ)$. Every $X\in\mathcal{A}$ satisfies \textup{(F3$'$)}.
  The class $\mathcal{A}$ contains every abelian variety and every smooth
  projective curve, and it is closed under products and under images of
  surjective morphisms. So it contains every product of curves, every
  symmetric product of a curve, and every smooth projective variety onto
  which a product of curves or an abelian variety maps surjectively.
\item Every $X$ of dimension $n$ satisfies \textup{(F3$'$)} in the degrees
  $2p$ with $p\in\{0,1,n-1,n\}$, so every $X$ of dimension at most three
  satisfies it in every degree.
\end{enumerate}
\end{proposition}

\begin{proof}
(i) If $\mathbf{HC}$ holds, then $\Hdg^{p}(X)$ consists of algebraic classes,
which lie in $\mathrm{Ab}^{p}(X)$. Conversely, if the conjecture holds for
abelian varieties, every $\gamma$ in \Cref{def:f3prime} is algebraic, so
every $\Gamma_{*}\gamma$ is, since composition with a correspondence carries
algebraic classes to algebraic classes \cite[Chapter~16]{Ful98}; then
\textup{(F3$'$)} says that every Hodge class is algebraic.

(ii) On an abelian variety $A$ take $B=A$ and $\Gamma$ the diagonal. On
$X=A\times\PP^{1}$ the K\"unneth formula gives
$H^{*}(X,\QQ)=\mathrm{pr}^{*}H^{*}(A,\QQ)\oplus
\mathrm{pr}^{*}H^{*}(A,\QQ)\cup\eta$, with $\eta$ the class of a point of
$\PP^{1}$ pulled back to $X$, which is algebraic. The first summand is
$\Gamma_{*}H^{*}(A,\QQ)$ for $\Gamma$ the transpose of the graph of
$\mathrm{pr}$, and the second is $\Gamma'_{*}H^{*}(A,\QQ)$ for
$\Gamma'=\Gamma\cdot\mathrm{pr}_{X}^{*}\eta$. So $X\in\mathcal{A}$, and
(iii) applies.

(iii) Let $X\in\mathcal{A}$ and fix $p$. Since $H^{2p}(X,\QQ)$ is finite
dimensional, finitely many classes $\Gamma_{j*}y_{j}$ span it, and we may
take each $\Gamma_{j}$ of pure codimension $c_{j}$ on $B_{j}\times X$ and each
$y_{j}$ of pure degree $k_{j}$, with $k_{j}+2c_{j}-2\dim B_{j}=2p$. Then
$\Gamma_{j*}\colon H^{k_{j}}(B_{j},\QQ)\to H^{2p}(X,\QQ)$ is a morphism of
Hodge structures up to a Tate twist, and the sum of these maps is a
surjective morphism of polarisable Hodge structures onto $H^{2p}(X,\QQ)$.
Polarisable Hodge structures are semisimple, so its kernel has a complement
that is a sub-Hodge structure and maps isomorphically onto $H^{2p}(X,\QQ)$;
hence every Hodge class of degree $2p$ is the image of a Hodge class, that is
$\sum_{j}\Gamma_{j*}\gamma_{j}$ with each $\gamma_{j}$ a Hodge class on
$B_{j}$. So $\Hdg^{p}(X)=\mathrm{Ab}^{p}(X)$.

An abelian variety lies in $\mathcal{A}$ through its diagonal. A curve $C$
of genus $g\ge1$ lies in $\mathcal{A}$ through an Abel--Jacobi embedding
$\iota\colon C\to J(C)$: the transpose of its graph induces $\iota^{*}$,
which is an isomorphism on $H^{0}$ and $H^{1}$ and sends the class of the
theta divisor to $g$ times the class of a point, so it is surjective; the
projective line is the image of an elliptic curve under a double cover, and
is covered by the last closure property. For products, if $H^{*}(X_{i},\QQ)$
is spanned by the $\Gamma_{i*}y$, then by the K\"unneth formula
$H^{*}(X_{1}\times X_{2},\QQ)$ is spanned by the
$(\Gamma_{1}\times\Gamma_{2})_{*}(y_{1}\otimes y_{2})$, with
$\Gamma_{1}\times\Gamma_{2}$ an algebraic cycle on
$(B_{1}\times B_{2})\times(X_{1}\times X_{2})$. For images, let
$f\colon Y\to X$ be a surjective morphism with $Y\in\mathcal{A}$, let
$r=\dim Y-\dim X$ and let $h$ be an ample class on $Y$. The correspondence
$\Psi=[\Gamma_{f}]\cdot\mathrm{pr}_{Y}^{*}h^{r}$ on $Y\times X$ induces
$\Psi_{*}y=f_{*}(y\cup h^{r})$, and
$\Psi_{*}f^{*}x=x\cup f_{*}h^{r}=c\,x$ with $c>0$ the degree of $h^{r}$ on a
general fibre. So $\Psi_{*}$ is surjective, and $H^{*}(X,\QQ)$ is spanned by
the classes $(\Psi\circ\Gamma)_{*}y$, with $\Psi\circ\Gamma$ algebraic
\cite[Chapter~16]{Ful98}. A symmetric product of a curve is the image of a
power of it.

(iv) Every class of degree $0$ or $2n$ is algebraic, and every Hodge class of
degree $2$ is algebraic by the Lefschetz theorem on $(1,1)$ classes. For
$n\ge2$, the hard Lefschetz theorem makes cup product with $h^{n-2}$, for $h$
an ample class, an isomorphism $H^{2}(X,\QQ)\to H^{2n-2}(X,\QQ)$ of Hodge
structures up to a Tate twist, so every Hodge class of degree $2n-2$ is
$h^{n-2}$ times a Hodge class of degree $2$, and is algebraic. Algebraic
classes lie in $\mathrm{Ab}^{p}(X)$. When $n\le3$ every $p$ between $0$ and
$n$ is one of $0,1,n-1,n$.
\end{proof}

\begin{remark}[Why \textup{(F3$'$)} is not closed here]\label{rem:f3primeopen}
\textup{(F3$'$)} is the statement that the Hodge conjecture reduces to
abelian varieties along algebraic correspondences, and
\Cref{prop:f3prime} records what is known of such a reduction: the
varieties of $\mathcal{A}$, whose whole cohomology is reached from abelian
varieties by algebraic correspondences, and the degrees in which the
conjecture itself is known. The first degree in which it is not known in
general is degree four on fourfolds, for fourfolds not known to lie in
$\mathcal{A}$. The square $S\times S$ of a K3 surface is the standard
example. Its Hodge classes in $H^{2}(S)\otimes H^{2}(S)$ include the classes
that induce the Hodge endomorphisms of the transcendental lattice $T(S)$, and
$S\times S$ lies in $\mathcal{A}$ as soon as the Kuga--Satake embedding of
$T(S)$ is induced by an algebraic cycle, the Kuga--Satake Hodge conjecture for
$S$, which is not known in general: the transpose of such a cycle maps the
cohomology of the Kuga--Satake variety onto $T(S)$, and the rest of
$H^{*}(S,\QQ)$ is algebraic. The rational Hodge isometries are algebraic
\cite{Bus19}, but a real multiplication is not an isometry, and
\Cref{rem:mumfordks} meets exactly this obstacle for the surfaces
$S_{\lambda}$.

For varieties whose cohomology is not of abelian type the position is
starker. By \Cref{prop:noabeliantype} no correspondence moves the $H^{3}$ of a
very general quintic threefold into the cohomology of an abelian variety, so
such a variety does not lie in $\mathcal{A}$. That is no obstruction to
\textup{(F3$'$)} itself, since a Hodge class spans a Hodge structure of Tate
type, which is of abelian type, and nothing in Hodge theory forbids it from
lying in $\mathrm{Ab}^{p}(X)$. Nor is it a route: to put a Hodge class into
$\mathrm{Ab}^{p}(X)$ one must construct an algebraic cycle on $B\times X$
whose class carries it, which is the kind of construction the Hodge
conjecture asks for, now on $B\times X$ instead of on $X$. The methods of
this paper construct cycles on abelian varieties of Weil type and on nothing
else, so they give no case of \textup{(F3$'$)} outside $\mathcal{A}$. We
therefore leave it open, and \Cref{thm:closure} treats it as an open
statement.
\end{remark}

\begin{setupenv}[The rule set]\label{setup:rules}
Item (XXXIII) of \Cref{sec:verification} carries forty-two statements and
twenty-two rules; the propagation statement (P2) is carried family by
family, as (P2) for the split and for the other families of the imaginary
quadratic fields and of the fields of degree at least four. Each statement is marked proved in this paper, quoted from the
literature, or open. Each rule is a finite set of premises together with one
conclusion, and carries the label of the theorem of this paper that proves it;
\Cref{tab:rules} lists the rules whose conclusion is derived rather than
assumed. For a set $T$ of statements write $\Cn(T)$ for the least set that
contains $T$ and is closed under the rules, which exists because each rule has
a single conclusion and the consequence operator is therefore monotone. Write
$\mathbf{T}$ for the set of statements marked proved or quoted, and
$\mathbf{HC}$ for the Hodge conjecture for every smooth complex projective
variety. Call a set $S$ of open statements \emph{sufficient} if
$\mathbf{HC}\in\Cn(\mathbf{T}\cup S)$, and \emph{minimal} if no proper subset
of $S$ is sufficient.
\end{setupenv}

\begin{table}[htbp]
\centering
\small
\renewcommand{\arraystretch}{1.14}
\begin{tabular}{@{}>{\raggedright\arraybackslash}p{0.40\linewidth}>{\raggedright\arraybackslash}p{0.27\linewidth}l@{}}
\toprule
premises & conclusion & by \\
\midrule
the conjecture for abelian varieties, (F3)
  & $\mathbf{HC}$ & \Cref{ssec:closuremap} \\
(F3)
  & abelian varieties & \Cref{prop:f3ishc} \\
the conjecture for abelian varieties, (F3$'$)
  & $\mathbf{HC}$ & \Cref{prop:f3prime} \\
the Weil classes of every CM field, (F2)
  & abelian varieties & \Cref{ssec:closuremap} \\
$K$ imaginary quadratic, the fields of degree at least four
  & every CM field & \Cref{app:albert} \\
$\delta=\delta_{0}$, $\delta\ne\delta_{0}$
  & every discriminant & \Cref{prop:separate} \\
a base point, the reduction, the orbit, (P2) for the family
  & that family & \Cref{thm:semiregclosure} \\
the same three for a CM field of degree at least four, (P2) for the family
  & that family & \Cref{thm:cmpropagation} \\
its split families, its other families
  & the fields of degree at least four & \Cref{prop:cmweil} \\
$\delta=\delta_{0}$
  & $\delta\ne\delta_{0}$ & \Cref{prop:descent} \\
the split families of the fields of degree at least four
  & their other families; $\delta=\delta_{0}$ & \Cref{prop:descent} \\
a secant object in every dimension, Question 11.4
  & $\delta=\delta_{0}$ & \Cref{prop:splitgeom} \\
a Cohen--Macaulay support, smooth or singular
  & one secant object & \Cref{cor:disjointroutes} \\
one secant object, Question 11.4
  & $\mathbf{W}(K,4,\delta_{0})$ & \Cref{cor:markmancondition} \\
\midrule
(L), (M)
  & $\mathbf{HC}$ & \Cref{prop:lefmot} \\
(L), the deformation theorem for motivated classes
  & (V) & \Cref{prop:lefvhc} \\
(V), the accessibility of Hodge classes on abelian varieties
  & abelian varieties & \Cref{thm:vhcab} \\
\bottomrule
\end{tabular}
\caption{The rules of \Cref{setup:rules} whose conclusion is derived, with the
theorem that proves each; the rows for (P2) and for the Cohen--Macaulay
support each stand for two rules, the row for descent over the fields of
degree at least four for two, descent within a field and scalar extension
down to the imaginary quadratic fields, the second row is \Cref{prop:f3ishc},
which an earlier version of the set omitted, the third is
\Cref{prop:f3prime}(i) for the weaker statement (F3$'$) of
\Cref{def:f3prime}, the last three rows are the rules of
\Cref{ssec:bullseye}, one of whose conclusions, (V), is itself open, and the
one remaining rule records what the literature supplies on Weil classes. The
bounded criterion (P1) is absent: by
\Cref{thm:p1equivalent} it is equivalent to its own conclusion, so it is a
restatement and not a premise. Item (XXXIII) of
\Cref{sec:verification} carries the whole set as data and checks that every
label here is a label of this paper.}
\label{tab:rules}
\end{table}

\begin{theorem}[The closure of what is proved here]\label{thm:closure}
With the notation of \Cref{setup:rules}:
\begin{enumerate}[label=\textup{(\roman*)},leftmargin=2.6em]
\item $\mathbf{HC}\notin\Cn(\mathbf{T})$. Nothing in this paper, combined with the
  theorems it quotes, implies the Hodge conjecture.
\item There are exactly five minimal sufficient sets:
  \[
    \{\mathrm{F3}\},\quad\{\mathrm{L},\mathrm{M}\},\quad
    \{\mathrm{L},\mathrm{F3}'\},\quad\{\mathrm{V},\mathrm{F3}'\},\quad
    \{\mathrm{P2}_{\mathrm{split}},\mathrm{F2},\mathrm{F3}'\},
  \]
  with \textup{(L)}, \textup{(M)}, \textup{(V)} as in \Cref{def:lmv},
  \textup{(F2)}, \textup{(F3)} as in \Cref{cor:fourremain}, (F3$'$) as in
  \Cref{def:f3prime}, and $\mathrm{P2}_{\mathrm{split}}$ the propagation
  statement \textup{(P2)}, read in the corrected form of \Cref{rem:p2prime},
  for the split families of the CM fields of degree at least four. No
  minimal set contains \textup{(P2)} for an imaginary quadratic field or for
  a family of nontrivial discriminant: \Cref{prop:descent} derives those.
  The single open statement (F3) is sufficient, and it is equivalent to
  $\mathbf{HC}$ (\Cref{prop:f3ishc}); no other single open statement is.
  The sets $\{\mathrm{L},\mathrm{F3}\}$, $\{\mathrm{V},\mathrm{F3}\}$ and
  $\{\mathrm{P2}_{\mathrm{split}},\mathrm{F2},\mathrm{F3}\}$ are sufficient and
  not minimal. The last of the five is the only one with three elements and the
  only one whose derivation of $\mathbf{HC}$ uses a statement proved in this
  paper.
\item In each of the five, every element is necessary: removing any one
  leaves $\mathbf{HC}$ outside the closure.
\item The secant route lies in no minimal sufficient set. Granting both of its
  open demands, a secant object satisfying \eqref{eq:heisenberg} in every
  dimension and an affirmative answer to \cite[Question 11.4]{Mar25b}, puts
  $\mathbf{W}(K,n,\delta_{0})$ in the closure, and with it, by
  \Cref{prop:descent}, $\mathbf{W}(K,n,\delta)$ for every $\delta$; it leaves
  the Weil classes of the CM fields of degree at least four outside.
\item No statement proved in this paper rests, at any depth, on an open one,
  and the rule set is acyclic.
\item Every element of a minimal sufficient set other than
  $\mathrm{P2}_{\mathrm{split}}$ is a consequence of $\mathbf{HC}$, and $\{\mathrm{F3}\}$ and
  $\{\mathrm{L},\mathrm{M}\}$ are each equivalent to it; no implication from
  $\mathbf{HC}$ to \textup{(P2)} is known.
\item An earlier version of this theorem, over a rule set without
  \Cref{prop:descent}, stated in (iv) that the secant route stops at the
  trivial discriminant, and listed \textup{(P2)} for every family in the
  fifth set. Both came from the omission of the rules of
  \Cref{prop:descent}, and the five minimal sets are otherwise unchanged.
\end{enumerate}
\end{theorem}

\begin{proof}
The consequence operator is monotone, so $\Cn$ is computed by iterating the
one step map to a fixed point, and the set of subsets of the sixteen open
statements is finite. Item (XXXIII) of
\Cref{sec:verification} carries out the iteration and the exhaustive search
over all $2^{16}$ subsets of the sixteen open statements in exact arithmetic,
and Section 24 of the Lean certificate repeats (i), (ii), (iii) and (iv) in
the kernel, over the same rule set under the numbering recorded there,
checking in particular all sixteen single open statements and all one
hundred and twenty pairs. The content of the statement is the rule set itself, and every rule is
one of the theorems named in \Cref{tab:rules}. The sufficiency of (F3)
alone is the chain $\mathrm{(F3)}\Rightarrow$ abelian varieties by
\Cref{prop:f3ishc}, then abelian varieties $\wedge\,\mathrm{(F3)}
\Rightarrow\mathbf{HC}$; every sufficient set that contains (F3) and has
another element is therefore not minimal, and a sufficient set without
(F3) must reach $\mathbf{HC}$ through one of the two other rules with that
conclusion, \Cref{prop:lefmot} or \Cref{prop:f3prime}(i), the second of
which needs the conjecture for abelian varieties, reached by \textup{(V)},
by \textup{(L)} through \textup{(V)}, or by (P2) with (F2). For (vi):
\textup{(L)} follows from $\mathbf{HC}$ by \Cref{thm:equivalence},
\textup{(M)} by \Cref{prop:lefmot}, \textup{(V)} by \Cref{prop:lefvhc},
(F3$'$) by \Cref{prop:f3prime}(i), and (F2) and (F3) are special cases of
$\mathbf{HC}$.

In (iv), every base point the secant route produces is a member
$X\times\widehat{X}$, whose discriminant is trivial whatever $\beta$ is chosen
(\Cref{prop:splitgeom}(ii)); the deformation argument stays inside the family
of trivial discriminant, and it is \Cref{prop:descent}(i), not the
deformation, that carries the result to the other discriminants. It has no
base point over a field of degree at least four, so it reaches none of them.
\end{proof}

\begin{corollary}[The three statements, one of which is the
conjecture]\label{cor:fourremain}
The Hodge conjecture follows from the results of this paper together with the
following three statements. No rule of \Cref{setup:rules} derives it from
(F1) and (F2) together, and it follows from (F3) alone, which is equivalent
to it
(\Cref{prop:f3ishc}); so the list is not a reduction of the conjecture to
three smaller problems. With (F3) replaced by the weaker statement
\textup{(F3$'$)} of \Cref{def:f3prime}, the conjecture modulo abelian
varieties, the three statements (F1), (F2), \textup{(F3$'$)} still suffice,
and no rule of \Cref{setup:rules} derives the conjecture from a proper
subset of them.
\begin{enumerate}[label=\textup{(F\arabic*)},leftmargin=3.0em,itemsep=4pt]
\item \emph{One propagation statement, for every CM field.} \textup{(P2)},
  in the corrected form of \Cref{rem:p2prime}: at one base point of each
  Weil family of each CM field some perfect complex whose Chern character is
  a nonzero multiple of a rational Weil class plus a polynomial in the
  polarisation has injective semiregularity map. Its first form, with Chern
  character exactly a multiple of the Weil class, is false for every
  imaginary quadratic field and every $n\ge2$ (\Cref{thm:p2false}): such a
  complex would have to deform to non-algebraic Weil tori, whose coherent
  sheaves have no Chern character below the top degree. The corrected form
  closes
  $\mathbf{W}(K,n,\delta)$
  for every imaginary quadratic $K$, every $n$ and every $\delta$ at once, by
  \Cref{thm:semiregclosure}, and
  $\mathbf{W}(F,n,\delta)$ for every CM field $F$ of higher degree, by
  \Cref{thm:cmpropagation}, because those theorems put $\Sigma=\cD$, so the
  Weil class is algebraic at every member of the family and not only at the
  general one, and because \Cref{thm:everyfamily} and \Cref{thm:cmbasepoint}
  supply the base point in every one of those families. By
  \Cref{prop:descent} it is enough to have it for the split families of the
  CM fields of degree at least four, in infinitely many dimensions.
  \Cref{thm:p2forces}, \Cref{thm:p2support}, \Cref{thm:nosum},
  \Cref{thm:p2numerical} and \Cref{thm:p2endo} describe the object of the
  first form and are vacuous by \Cref{thm:p2false}; for the corrected form
  \Cref{thm:p2primenumber} computes the number that $\dim\Ext^{2}(E,E)$ has
  to reach, $7n^{2}-2n$ for a general shape. For one family its
  conclusion is equivalent to the Lefschetz standard conjecture for one
  variety (\Cref{cor:weilequiv}). The bounded criterion \textup{(P1)} is not a
  second form of the same demand: by
  \Cref{thm:p1equivalent} it is equivalent to the conclusion it was to imply,
  and \Cref{thm:p1forces} shows that the one cycle it could have been tested
  on, the transport of the base cycle, does not satisfy it.
\item \emph{The classes beyond the Weil lines.} Every Hodge class on an
  abelian variety that lies outside the subring generated by divisor classes
  and by Weil classes of quotients is algebraic. Through dimension five there
  is no such class \cite{MZ99}; in dimension eight there are
  (\Cref{prop:mumfordproduct}), so this is not a reduction of the conjecture
  for abelian varieties to Weil classes, which is false as a statement about
  Hodge rings: it is the conjecture for those classes. Its smallest case is
  brought down to one Kuga--Satake class in \Cref{thm:mumfordks}; it holds at
  every CM point (\Cref{thm:mumfordcm}), and on all members of a Mumford
  family it is equivalent to the Lefschetz standard conjecture for one
  ninefold (\Cref{cor:mumfordequiv}). Its classes are rigid along the base
  curve (\Cref{thm:mumfordrigid}), and one complex at one CM point with
  $\Ext^{2}$ of dimension $119$ would close it (\Cref{thm:mumfordnumerical}),
  as would a bound on the cycles that represent them at the closed points of
  the reductions of the curve modulo primes (\Cref{thm:modpdensity}).
\item \emph{The varieties that are not abelian.} The Hodge conjecture holds
  for every smooth complex projective variety that is not an abelian variety.
  This is the Hodge conjecture itself. The varieties it covers include
  $A\times\PP^{1}$ for every abelian variety $A$, and the conjecture for
  $A\times\PP^{1}$ gives it for $A$ (\Cref{prop:f3ishc}); so (F3) implies
  the conjecture for abelian varieties, hence (F1) and (F2), and hence
  $\mathbf{HC}$. Nor is it a reduction to abelian varieties in the other
  direction: no realisation of the cohomology of an arbitrary variety inside
  that of abelian varieties exists (\Cref{prop:noabeliantype}), and
  \Cref{thm:hyperplane} and \Cref{thm:returntrip} close the two routes
  through hyperplane sections that would be the natural candidates for one.
\end{enumerate}
\end{corollary}

\begin{proof}
The three together are sufficient by the derivation of \Cref{thm:closure},
and (F3) alone by \Cref{prop:f3ishc}. That (F1) and (F2) together are not
sufficient in the rule set is \Cref{thm:closure}(ii), since every minimal
sufficient set contains (F3), \textup{(F3$'$)} or \textup{(M)}; the
statement about (F1), (F2), \textup{(F3$'$)} is \Cref{thm:closure}(ii) and
(iii) for the last of the five minimal sets. The Weil classes of the CM
fields of degree at least four are not a fourth statement: they are derived
from (F1) by \Cref{thm:cmpropagation}. Nor is the bounded criterion a
premise, since by \Cref{thm:p1equivalent} it is a restatement of its own
conclusion.
\end{proof}

\begin{remark}[How the earlier count went wrong, and the wording of (F2)]
\label{rem:f3collapse}
An earlier version of this corollary said that the conjecture follows from
no proper subset of the three statements. That was a statement about a rule
set that used (F3) only through $\text{(abelian varieties)}\wedge
\text{(F3)}\Rightarrow\mathbf{HC}$ and did not contain the elementary rule
$\text{(F3)}\Rightarrow\text{(abelian varieties)}$ of \Cref{prop:f3ishc}.
The rule set can show that a set suffices; it cannot show that a set is
needed, except relative to the implications written into it.

The wording of (F2) has the same feature, one degree at a time. Let
$R(A)\subseteq H^{*}(A,\QQ)$ be the subring generated by divisor classes and
Weil classes of quotients. Suppose some Hodge class $\beta$ of degree $2p$ on
$A$ lies outside $R(A)$, and let $\alpha$ be any Hodge class of degree $2p$.
If $\alpha\notin R(A)$, (F2) makes it algebraic. If $\alpha\in R(A)$, then
$\alpha+\beta\notin R(A)$, because $R(A)$ is closed under subtraction, so
(F2) makes $\alpha+\beta$ and $\beta$ algebraic, and with them $\alpha$.
So in every degree in which (F2) says anything it is the whole conjecture for
that degree: on the square of a Mumford fourfold it gives every Hodge class of
$H^{4}$, the six products of divisor classes included. Whether (F2) alone
gives the conjecture for every abelian variety is not decided here, and
nothing depends on it: the rule set does not use such an implication, and
adding it would not change \Cref{thm:closure}(ii), since (F3) alone is
already sufficient.

The form of (F3) that is not the conjecture itself asks only for
algebraicity modulo classes coming from abelian varieties. It is
\textup{(F3$'$)} of \Cref{def:f3prime}; \Cref{prop:f3prime} says what is
known of it, and \Cref{rem:f3primeopen} why it is not closed.
\end{remark}

\begin{figure}[tbp]
\centering
  \includegraphics[width=0.92\textwidth]{fig_frontier}
\caption{The rule set of \Cref{setup:rules} with the rules of
\Cref{ssec:bullseye}, drawn as a hypergraph. A rule with several premises is
a junction marked $\wedge$, with arrows from its premises and one arrow to its
conclusion; each rule carries the result that proves it, and the junction
marked \S16.5 is the closure map of \Cref{ssec:closuremap}. Clay marks an
open statement, grass what this paper proves, indigo a derived statement,
and the tag $\Leftarrow\mathbf{HC}$ an open statement that the conjecture
implies: every clay node except (P2). The arrow from (F3) to the conjecture
is \Cref{prop:f3ishc}, and the junction above the conjecture for abelian
varieties is \Cref{prop:f3prime}(i). The five minimal sufficient sets of
\Cref{thm:closure}(ii) are drawn as nested ribbons under the arrows they
use, widest first, and listed in the cards below: (F3) alone; (L) with (M);
(L) through (V) with \textup{(F3$'$)}; (V) with \textup{(F3$'$)}; and, on
the left, (P2) with (F2) and \textup{(F3$'$)}. A shared arrow shows the
bands of every route through it. The secant route, which reaches the
imaginary quadratic families only, is drawn in \Cref{fig:closure}.}
\label{fig:frontier}
\end{figure}

\begin{remark}[The target, and what it would take]\label{rem:bullseye}
With the rules of \Cref{ssec:bullseye}, \Cref{prop:f3ishc} and
\Cref{prop:f3prime} in place, \Cref{thm:closure} changes the reading of
\Cref{cor:fourremain} in three ways. First, the route through this paper is
a route only with (F3) replaced by \textup{(F3$'$)}. As stated, its last
statement (F3) is the conjecture (\Cref{prop:f3ishc}), so (F1) and (F2), and
everything this paper proves about Weil classes, drop out of every minimal
derivation that uses (F3). With \textup{(F3$'$)} in its place the route
$\{\mathrm{P2},\mathrm{F2},\mathrm{F3}'\}$ is minimal, and it is not the
shortest: (F1) and (F2), which between them give the conjecture for abelian
varieties along this route, are replaced on the routes of the literature by
the single statement \textup{(V)}, or by \textup{(L)}, which implies it
(\Cref{thm:vhcab}). Second, every element of every minimal sufficient set
except (P2) is a consequence of the conjecture (\Cref{thm:closure}(vi)), and
two of the sets, $\{\mathrm{F3}\}$ and $\{\mathrm{L},\mathrm{M}\}$, are
each equivalent to it. (P2) is a statement about semiregularity maps, and a
proof of the conjecture need not prove it; in its first form, with a Chern
character that is exactly a Weil class, it is false (\Cref{thm:p2false}), and
the route runs through the corrected form of \Cref{rem:p2prime}. Third, every minimal sufficient set
contains (F3), \textup{(F3$'$)} or \textup{(M)}: the varieties that are not
abelian are reached only through the conjecture for them, which is the whole
conjecture, through the conjecture modulo abelian varieties, which is open
(\Cref{rem:f3primeopen}), or through motivated classes.

So the target at the centre is \textup{(L)}. It lies in two of the five
minimal sets, it gives \textup{(V)} and with it the whole abelian half, the
conjecture implies it, and with \textup{(M)} it is the conjecture. For each
$X$ it asks for one explicit Hodge class on $X\times X$ to be algebraic, the
class inducing $\Lambda$. It is open, as are \textup{(M)}, \textup{(V)}
and \textup{(F3$'$)}, and this paper proves none of them in general:
\Cref{cor:mumfordlefschetz} adds the total spaces of Mumford families to the
known cases of \textup{(L)}, \Cref{prop:f3prime} records the known cases of
\textup{(F3$'$)}, and the routes of the literature use no other theorem of
this paper.
\end{remark}

\subsubsection*{Two of the three are not reductions}

A reader might hope that (F2) and (F3) are reductions in the ordinary sense,
statements to the effect that every Hodge class is carried, by some
correspondence, to a class on a Weil type variety or on an abelian variety,
so that the whole conjecture would follow from the Weil classes alone. Both
hopes are refutable, and we refute them, because a frontier that contains a
false statement is not a frontier.

\begin{proposition}[The Hodge ring is not generated by divisor and Weil
classes]\label{prop:mumfordproduct}
Let $X$ be a complex abelian fourfold whose Hodge group is a $\QQ$-form of an
almost direct product of three copies of $\mathrm{SL}_{2}$, acting on
$V=H^{1}(X,\QQ)$ through the tensor product $V_{1}\otimes V_{2}\otimes V_{3}$
of the three standard representations, with $\End^{0}(X)=\QQ$; such
fourfolds exist and form families of dimension one, by Mumford
\cite{Mum69}, and by \cite[2.5(i)]{MZ99} they are exactly the fourfolds
with $\End^{0}(X)=\QQ$ and Hodge group smaller than $\Sp_{8}$. Then:
\begin{enumerate}[label=\textup{(\roman*)},nosep,leftmargin=2.6em]
\item the Hodge ring of $X$ is generated by the polarisation:
  $\dim_{\QQ}\Hdg^{1}(X)=\dim_{\QQ}\Hdg^{2}(X)=1$;
\item $\dim_{\QQ}\Hdg^{1}(X\times X)=3$ and $\dim_{\QQ}\Hdg^{2}(X\times X)=8$,
  while the products of divisor classes span a subspace of dimension $6$
  in $H^{4}(X\times X,\QQ)$;
\item the two remaining classes lie in $H^{4}(X\times X,\QQ)$ and are Weil
  classes of no CM field acting on $X\times X$, on a quotient of it, or on an
  abelian subvariety of it.
\end{enumerate}
In particular the statement that every Hodge class on an abelian variety
lies in the subring generated by divisor classes and Weil classes of
quotients is false in dimension eight, though true through dimension five
\cite{MZ99}.
\end{proposition}

\begin{proof}
Hodge classes are the invariants of the Hodge group, so
$\Hdg^{k}(Y)\otimes\CC$ is the space of $\Hg(Y)_{\CC}$-invariants in
$\bigwedge^{2k}H^{1}(Y,\CC)$, and $\Hg(X\times X)=\Hg(X)$ acting
diagonally on $V\oplus V$. Over $\CC$ the group is
$\mathrm{SL}_{2}^{3}$ and the weights of $V$ under its maximal torus are the
eight sign vectors $(\pm1,\pm1,\pm1)$. Item (XXXV) of \Cref{sec:verification}
decomposes every $\bigwedge^{k}V$ by peeling highest weights and finds
\begin{gather*}
  \textstyle\bigwedge^{2}V=V(2,2,0)\oplus V(2,0,2)\oplus V(0,2,2)\oplus\mathbf{1},\\
  \textstyle\bigwedge^{4}V=V(2,2,2)\oplus\bigoplus_{\text{cyc}}V(4,0,0)
  \oplus\bigoplus_{\text{cyc}}V(2,2,0)\oplus\mathbf{1},
\end{gather*}
each with one trivial summand, which is (i). For (ii),
$\bigwedge^{4}(V\oplus V)=\bigoplus_{a}\bigwedge^{a}V\otimes\bigwedge^{4-a}V$
and the invariants in $\bigwedge^{a}V\otimes\bigwedge^{4-a}V$ are
$\Hom_{\mathrm{SL}_{2}^{3}}(\bigwedge^{a}V,\bigwedge^{4-a}V)$, since every
summand is self-dual; the counts are $1,1,4,1,1$, total $8$, the $4$ coming
from the four distinct irreducible summands of $\bigwedge^{2}V$. Likewise the
invariants in $\bigwedge^{2}(V\oplus V)$ number $1+1+1=3$. The divisor
classes of $X\times X$ are the three symplectic forms $\psi_{11}$, $\psi_{12}$,
$\psi_{22}$ on $V\oplus V$ built from the polarisation $\psi$ of $V$, and the
same item multiplies them in the exterior algebra of $V\oplus V$ and finds
their six pairwise products linearly independent; it also counts the
$\Sp_{8}$-invariants of $\bigwedge^{4}(V\oplus V)$ by the Weyl character
formula and finds six, so the six products are all the classes a Hodge group
equal to $\Sp_{8}$ would allow. For (iii), the Weil classes of an abelian
variety $Y$ with respect to a CM field $F\subseteq\End^{0}(Y)$ are, by
definition, the classes in $\bigwedge^{2k}_{F}H^{1}(Y,\QQ)\subseteq
H^{2k}(Y,\QQ)$ with $2k=\dim_{F}H^{1}(Y,\QQ)$ (\Cref{setup:cmfield}). A CM
field acting on $X\times X$ lies in $\End^{0}(X\times X)=M_{2}(\QQ)$, so it is
imaginary quadratic, $2k=8$, and its Weil classes lie in $H^{8}$, while the two
classes of (ii) have degree four. Every proper
abelian subvariety and every proper quotient of $X\times X$ is isogenous to
$X$, because $X$ is simple, and $\End^{0}(X)=\QQ$ contains no CM field, so $X$
carries no Weil classes at all. This is the statement of \cite[0.1(iii)]{MZ99}
that the exceptional classes of $X\times X$ are not of Weil type. The last
sentence is (ii) and (iii) together with the theorem of \cite{MZ99} for
dimension at most five.
\end{proof}

\begin{proposition}[No Hodge-theoretic reduction to abelian varieties]
\label{prop:noabeliantype}
Let $Y$ be a very general smooth quintic threefold in $\PP^{4}$, and let
$V=H^{3}(Y,\QQ)$ with its polarised Hodge structure of weight three. Then $V$
is not isomorphic to a subquotient of any Hodge structure of the form
$\bigoplus_{i}T_{i}$, where each $T_{i}$ is obtained from the $H^{1}$ of
finitely many abelian varieties by tensor products, duals and Tate twists. In
particular the Hodge structure $H^{*}(Y,\QQ)$ is not of abelian type, and no
correspondence between $Y$ and an abelian variety, or a product of abelian
varieties, induces an injection of $V$ into the cohomology of the latter.
\end{proposition}

\begin{proof}
Write $\mathrm{MT}(W)$ for the Mumford--Tate group of a rational Hodge
structure $W$ and $\mu_{W}\colon\mathbb{G}_{m}\to\mathrm{MT}(W)_{\CC}$ for
its Hodge cocharacter, acting on $W^{p,q}$ by $z^{p}$. The adjoint action of
$\mu_{W}$ on $\operatorname{Lie}\mathrm{MT}(W)_{\CC}\subseteq\End(W_{\CC})$ has weights
among the differences $p-p'$ of two Hodge indices of $W$.

If $W$ has weight one, these differences lie in $\{-1,0,1\}$. If $W'$ is a
subquotient of a tensor construction $T$ on $W$, then $\mathrm{MT}(W')$ is the
image of $\mathrm{MT}(W)$ acting on $W'$, and $\mu_{W'}$ is the image of
$\mu_{W}$; so $\operatorname{Lie}\mathrm{MT}(W')$ is a quotient of $\operatorname{Lie}\mathrm{MT}(W)$
as a representation of $\mathbb{G}_{m}$ through $\mu$, and its weights are a
subset of $\{-1,0,1\}$. The same holds with $W=\bigoplus H^{1}(A_{i})$.
Hence for every Hodge structure of abelian type the adjoint weights of the
Hodge cocharacter lie in $\{-1,0,1\}$.

For $V=H^{3}(Y,\QQ)$ we show that the adjoint weight $3$ occurs. The
canonical bundle of a quintic threefold is trivial by adjunction, so
$h^{3,0}(Y)=1$; pick $0\ne e\in H^{3,0}$. The intersection form $\psi$ on $V$
is symplectic and pairs $H^{3,0}$ with $H^{0,3}$ perfectly, so the rank one
endomorphism $x\mapsto\psi(x,e)\,e$ of $V_{\CC}$ is nonzero, sends $H^{0,3}$
to $H^{3,0}$ and kills the rest, and lies in $\mathfrak{sp}(V_{\CC},\psi)$
because $\psi(\psi(x,e)e,y)+\psi(x,\psi(y,e)e)
=\psi(x,e)\bigl(\psi(e,y)+\psi(y,e)\bigr)=0$; its adjoint weight is
$3-0=3$. It remains to see that $\mathfrak{sp}(V_{\CC},\psi)\subseteq
\operatorname{Lie}\mathrm{MT}(V)_{\CC}$. Let $S$ be the parameter space of smooth quintic
threefolds, connected and nonsingular, and $\mathbb{V}$ the polarisable
variation of Hodge structure $R^{3}$ of the universal family. The connected
algebraic monodromy group $G^{0}$ of $\mathbb{V}$ contains the connected
monodromy group of any Lefschetz pencil of quintics, that is of hyperplane
sections of the fifth Veronese image of $\PP^{4}$; for a Lefschetz pencil of
odd fibre dimension the image of the monodromy representation on the
$\ell$-adic vanishing cohomology $E\otimes\QQ_{\ell}$ is open in
$\Sp(E\otimes\QQ_{\ell},\psi)$ \cite[Theorem 5.10]{Del74}, so the Zariski
closure of the monodromy group over $\QQ$ is $\Sp(E,\psi)$; here
$E=H^{3}(Y_{t},\QQ)$, the vanishing cohomology being the orthogonal
complement of the image of $H^{3}(\PP^{4},\QQ)=0$. Hence $G^{0}\supseteq
\Sp(V,\psi)$, and $G^{0}\subseteq\Sp(V,\psi)$ since the monodromy preserves
$\psi$. At a very general
point $Y$ the group $G^{0}$ is a normal subgroup of the derived group of
$\mathrm{MT}(V)$ \cite[Theorem 1]{And92}, so $\Sp(V,\psi)\subseteq
\mathrm{MT}(V)^{\mathrm{der}}\subseteq\Sp(V,\psi)$, and the adjoint weight
$3$ occurs in $\operatorname{Lie}\mathrm{MT}(V)$. So $V$ is not of abelian type, and
neither is any Hodge structure containing it as a subquotient, in particular
$H^{*}(Y,\QQ)$. A correspondence inducing an injection of $V$ into the
cohomology of an abelian variety would exhibit $V$ as a sub-Hodge structure
of a Tate twist of that cohomology, which is of abelian type. Item (XXXV)
records the Hodge numbers $(1,101,101,1)$ from the Jacobian ring of the
Fermat quintic, checks the symplectic identity and the weight of the
displayed endomorphism in a four-dimensional model with the same weights, and
lists the adjoint weights that occur.
\end{proof}

\begin{remark}[What the two propositions do to the frontier]
\label{rem:notreductions}
Neither proposition changes \Cref{thm:closure}: the rule set uses (F2) and
(F3) as implications, and an implication whose hypothesis is the conjecture
for one class of Hodge classes and whose conclusion is the conjecture for a
larger class is true whatever the classes are. What does change it is
\Cref{prop:f3ishc}: the varieties of (F3) are not a class disjoint from the
abelian ones but contain them, through products with $\PP^{1}$, so (F3) is
the whole conjecture and not its non-abelian half. What they change is how the
two statements may be read. Before them one could hope that (F2) and (F3) were
structural facts awaiting proof, so that the conjecture would be a theorem
about Weil classes and nothing else; \Cref{prop:mumfordproduct} shows that the
Hodge classes on abelian varieties are not exhausted by divisor and Weil
classes once the dimension reaches eight, and \Cref{prop:noabeliantype} that
the cohomology of a general variety cannot be moved into that of abelian
varieties at all. Each of the two is therefore the Hodge conjecture for a class
of varieties on which the methods of this paper, which are methods about Weil
lines, have no purchase. The product $X\times X$ is itself of
Weil type for every imaginary quadratic $K$, through an embedding
$K\subset M_{2}(\QQ)=\End^{0}(X\times X)$ whose Rosati involution for the
product polarisation is conjugation, with signature $(4,4)$; its Weil classes
lie in $H^{8}$, and its Hodge group is far smaller than $\SU(4,4)$. The two
classes in $H^{4}$ are of a different kind, and they show that even on a
variety of Weil type the Weil line need not be the whole of the exceptional
cohomology once the Hodge group is small. They are the smallest instance of
(F2), and \Cref{ssec:mumford} identifies them and says what one K3 surface
would have to supply for them to be algebraic.
\end{remark}

\begin{figure}[tbp]
\centering
  \includegraphics[width=\textwidth]{fig_closure}
\caption{The route of this paper in \Cref{thm:closure}, drawn as a stack of
plates, one per level of the derivation; the other routes of the literature
are drawn in \Cref{fig:frontier}. The spine in indigo runs from the Weil families at the bottom to
the conjecture at the top. The two statements (F2) and (F3) enter in clay from
the right, each at the level where it is needed; (F2) is the conjecture for a
class of Hodge classes that no Weil line reaches, not a reduction, and (F3) is
the whole conjecture (\Cref{prop:f3ishc}), so the route drawn is minimal
only with (F3) replaced by \textup{(F3$'$)} of \Cref{def:f3prime}. From
below,
the propagation statement (P2) in grass reaches the two discriminant cases of
the quadratic fields and the fields of higher degree alike, because
\Cref{thm:everyfamily} and \Cref{thm:cmbasepoint} put a base point in every
family; it is read in the corrected form of \Cref{rem:p2prime}. The two
edges inside the levels are \Cref{prop:descent}: descent from the trivial
discriminant to the others, and scalar extension from the fields of degree at
least four to the quadratic fields; with them (P2) is needed only for the
split families of the fields of higher degree. The three
pillars under the grass node are the results of \Cref{thm:p2numerical},
\Cref{thm:p2support} and \Cref{cor:hhfactor} on the first form, in which the
Chern character is exactly a Weil class; by \Cref{thm:p2false} that form is
never met, so the number on the first pillar is attained by no complex. The
secant route in clay reaches the trivial discriminant, since the member it is
built on is split, and through descent the rest of the quadratic case; it has
no edge to the fields of higher degree. The bounded criterion (P1) is drawn apart, in slate, closed on
itself: it is equivalent to its own conclusion.}
\label{fig:closure}
\end{figure}

\begin{remark}[Where this leaves the secant route]\label{rem:secantplace}
\Cref{thm:closure}(iv) is a statement about the shape of the argument, not
about its difficulty, and we say plainly what it does and does not
mean.

It does not mean the secant route is worthless. It is the only route on which
anything above dimension four has ever been proved, it is the route that
settles $n=2$ at every discriminant and $n=3$ at trivial discriminant, and the
reduction of
\Cref{thm:factor} turns its remaining demand into one identity in the Yoneda
algebra of a sheaf on a fourfold, which is the form in which a demand can be
attacked. What \Cref{cor:cmdichotomy}, \Cref{thm:noci} and
\Cref{thm:nosplit} do to that identity is narrow it to two shapes of
support, neither resolved by sums of line bundles, and, in the smooth shape,
to four discriminants.

What it does mean is that finishing the secant route would not finish the
problem, and would not reduce it either. Through \Cref{prop:descent} it would
give every imaginary quadratic family, the trivial discriminant directly and
the others by descent, which is a great deal; an earlier version of this
remark said that it stops at the trivial discriminant, and that was an
omission of the rule set. It would give nothing over the CM fields of degree
at least four, whose split families need (P2) for themselves, and by
\Cref{prop:descent}(ii) that statement gives the imaginary quadratic fields as
well. So the statement that would have to be proved anyway is the
propagation statement for those split families, and once proved it subsumes
the secant route entirely.

It is tempting to read (C1) and (C2) as what remains for the imaginary
quadratic case, and the computation shows why that reading is wrong: it puts
the demand for a construction where a theorem about propagation belongs.
(C1) and (C2) are demands whose satisfaction would be a considerable
achievement, would settle the Weil classes of every imaginary quadratic field,
and would close nothing that (P2) does not close.
\end{remark}

\subsection{The first classes beyond the Weil lines}\label{ssec:mumford}

Statement (F2) is the Hodge conjecture for the Hodge classes on abelian
varieties that no Weil line reaches, and \Cref{prop:mumfordproduct} exhibits
the first of them: two classes in $H^{4}(X\times X,\QQ)$ for an abelian
fourfold $X$ of Mumford type. None of the constructions of the earlier
sections acts on them. In this subsection we identify them as a real
multiplication on the primitive part of $H^{2}(X)$, and we prove that they
are algebraic as soon as the Kuga--Satake class of one K3 surface of Picard
number thirteen attached to $X$ is algebraic. We do not prove that the
Kuga--Satake class is algebraic, and we know of no proof. What changes is the
kind of question: two classes on an eightfold are traded for one
correspondence of a well-studied type, between a surface and its Kuga--Satake
variety. The field enters through the Clifford relations, which force the
Kuga--Satake map to see all of it.

\begin{setupenv}[Mumford type]\label{setup:mumford}
Let $X$ be as in \Cref{prop:mumfordproduct}, with $V=H^{1}(X,\QQ)$,
polarisation $\psi$ and polarisation class $\theta\in H^{2}(X,\QQ)$, and let
$G=\Hg(X)$. Over $\bar{\QQ}$ the Hodge cocharacter is nontrivial on exactly
one of the three factors of $G$: on $V_{i}$ it has weights $\pm a_{i}/2$ with
$a_{i}\ge0$, and $V$ has only the two weights $\pm\frac12$, so exactly one
$a_{i}$ is nonzero. Because
$\End^{0}(X)=\QQ$ the group $G$ is therefore $\QQ$-simple: the Lie algebra of
a $\QQ$-factor not containing that geometric factor would consist of Hodge
classes in $\End(V)$, that is of endomorphisms of $X$. So $G$ is the image in
$\GL(V)$ of $\Res_{F/\QQ}D^{1}$, where $F$ is a totally real cubic field, $D$
is a quaternion algebra over $F$ split at one real place $\sigma_{3}$ and
definite at the other two, $\sigma_{1}$ and $\sigma_{2}$, with
$\operatorname{Cor}_{F/\QQ}D\cong M_{8}(\QQ)$, and $D^{1}$ is the group of
elements of reduced norm one \cite{Mum69,Gal00}; the kernel is the central
subgroup of triples of signs with product one. Over $\bar{\QQ}$,
$V=V_{1}\otimes V_{2}\otimes V_{3}$, where $V_{i}$ is the standard
representation of $D\otimes_{F,\sigma_{i}}\bar{\QQ}=M_{2}(\bar{\QQ})$; the
Galois group permutes the three factors as it permutes the three embeddings,
and $H^{1,0}(X)=V_{1}\otimes V_{2}\otimes V_{3}^{1,0}$. Put
\[
  T=\operatorname{Lie}G=D^{0},
  \qquad
  q_{\lambda}(x,y)=\Tr_{F/\QQ}\bigl(\lambda\operatorname{trd}(xy)\bigr)
  \quad(\lambda\in F^{\times}),
\]
where $D^{0}$ is the space of quaternions of reduced trace zero, regarded as
a $\QQ$-vector space of dimension nine, and $\operatorname{trd}$ is the
reduced trace. Over $\bar{\QQ}$, $T=T_{1}\oplus T_{2}\oplus T_{3}$ with
$T_{i}=\mathfrak{sl}(V_{i})$, and $a\in F$ acts on $T_{i}$ by $\sigma_{i}(a)$.
Write $m_{a}$ for multiplication by $a$. The space $T\subset\End(V)$ is a
sub-Hodge structure of weight zero.
\end{setupenv}

\begin{proposition}[The two classes are a real multiplication]%
\label{prop:mumfordrm}
In \Cref{setup:mumford}:
\begin{enumerate}[label=\textup{(\roman*)},nosep,leftmargin=2.6em]
\item $H^{2}(X,\QQ)=\QQ\theta\oplus U$, where $U$ is the primitive part, of
  dimension $27$, and over $\bar{\QQ}$
  \[
    U=U_{12}\oplus U_{13}\oplus U_{23},\qquad
    U_{ij}=S^{2}V_{i}\otimes S^{2}V_{j}\otimes{\textstyle\bigwedge^{2}}V_{k},
    \quad \{i,j,k\}=\{1,2,3\},
  \]
  each of dimension nine. The summand $U_{12}$ is of type $(1,1)$, and
  $U_{13}$ and $U_{23}$ have Hodge numbers $(3,3,3)$.
\item The algebra of Hodge endomorphisms of $H^{2}(X,\QQ)$ is $\QQ\times F$,
  acting by scalars on $\QQ\theta$, with $a\in F$ acting on $U_{ij}$ by
  $\sigma_{k}(a)$. Through Poincar\'e duality and the Lefschetz isomorphism
  $\theta^{2}\colon H^{2}\to H^{6}$ this algebra is the space of Hodge classes
  in the K\"unneth component $H^{2}(X)\otimes H^{2}(X)$ of
  $H^{4}(X\times X,\QQ)$; the subalgebra $\QQ\times\QQ$ is spanned by the
  products of divisor classes that lie there, and the two classes of
  \Cref{prop:mumfordproduct}(ii) may be taken to be the elements of $F$ of
  trace zero.
\item $T(-1)$ is a Hodge structure of K3 type with Hodge numbers $(1,7,1)$,
  and its Hodge endomorphisms are the $m_{a}$, $a\in F$. The $G$-invariant
  symmetric forms on $T$ are the $q_{\lambda}$, and $q_{\lambda}$ polarises
  $T(-1)$ if and only if $\lambda$ is totally positive.
\end{enumerate}
\end{proposition}

\begin{proof}
(i) is the decomposition of $\bigwedge^{2}(V_{1}\otimes V_{2}\otimes V_{3})$
in the proof of \Cref{prop:mumfordproduct}, where $U_{ij}$ appears as the
summand of highest weight $2$ on the $i$-th and $j$-th factors. The Hodge
cocharacter acts trivially on $V_{1}$ and $V_{2}$ and splits
$V_{3}=V_{3}^{1,0}\oplus V_{3}^{0,1}$, so $\bigwedge^{2}V_{3}$ has type
$(1,1)$ and $S^{2}V_{3}$ has one dimension of each of the types $(2,0)$,
$(1,1)$, $(0,2)$.

(ii) A Hodge endomorphism commutes with the Mumford--Tate group $G\cdot\mathbb{G}_{m}$, and
$\mathbb{G}_{m}$ acts on $H^{2}$ by scalars, so the Hodge endomorphisms are
the $G$-endomorphisms. Over $\bar{\QQ}$ the four summands $\QQ\theta$,
$U_{12}$, $U_{13}$, $U_{23}$ are irreducible and pairwise non-isomorphic, so
the $G$-endomorphisms are the combinations of the four projections. The
Galois group fixes the projection onto $\QQ\theta$ and permutes the other
three as it permutes the indices $k$, that is as it permutes the embeddings of
$F$. An endomorphism acting on $U_{ij}$ by $\sigma_{k}(a)$ is therefore
Galois invariant, hence rational, and these exhaust the rational points,
which are $\QQ\times F$. Poincar\'e duality identifies $\End H^{2}(X,\QQ)$
with $H^{6}(X)\otimes H^{2}(X)$, and $\theta^{2}\otimes1$ carries
$H^{2}\otimes H^{2}$ onto it; both are given by algebraic classes, and on an
abelian variety so is the inverse of the Lefschetz isomorphism
\cite[\S2]{Kle68}. The products of divisor classes in $H^{2}\otimes H^{2}$
are $\psi_{11}\psi_{22}$ and $\psi_{12}^{2}$ in the notation of the proof of
\Cref{prop:mumfordproduct}. They are invariant under the whole of
$\Sp(V,\psi)$, whose endomorphisms of $H^{2}=\QQ\theta\oplus U$ are spanned
by the identity and the projection onto $\QQ\theta$, and they are linearly
independent by item (XXXV); so they span $\QQ\times\QQ$, and any complement of
$\QQ\times\QQ$ in $\QQ\times F$ may serve as the two remaining classes.

(iii) The adjoint action of the Hodge cocharacter is trivial on $T_{1}$ and
$T_{2}$ and has weights $-1$, $0$, $1$ on $T_{3}$, so $T$ has Hodge numbers
$(1,7,1)$ in the types $(-1,1)$, $(0,0)$, $(1,-1)$, and $T(-1)$ is of K3
type. Its Hodge endomorphisms are again the $G$-endomorphisms, and the
argument of (ii), applied to the three pairwise non-isomorphic summands
$T_{i}$, gives $F$ acting by multiplication. On each $T_{i}$ the invariant
symmetric forms are the multiples of $\operatorname{tr}(xy)$, and
$\operatorname{tr}(x_{i}y_{i})=\sigma_{i}(\operatorname{trd}(xy))$ for
$x,y\in D^{0}$, so the Galois-invariant ones are the $q_{\lambda}$. At
$\sigma_{1}$ and $\sigma_{2}$ the algebra $D\otimes\RR$ is Hamilton's
quaternions, $T_{i}$ is of type $(1,1)$, and $\operatorname{trd}(x^{2})=-2\operatorname{nrd}(x)$ is
negative on nonzero pure quaternions. At $\sigma_{3}$ the form
$\operatorname{tr}(xy)$ on $\mathfrak{sl}_{2}(\RR)$ is negative on the line of
type $(1,1)$, the Lie algebra of the rotations fixed by the Hodge
cocharacter, and positive on its orthogonal complement, the real points of
the types $(2,0)$ and $(0,2)$. The Hodge--Riemann relations for a structure
of K3 type ask for exactly these signs after multiplication by
$\sigma_{i}(\lambda)$, so $q_{\lambda}$ polarises $T(-1)$ precisely when every
$\sigma_{i}(\lambda)$ is positive. Item (XLI) of \Cref{sec:verification}
computes the four projections from the Casimir operators of the three factors
on an explicit model and checks (i) and the Hodge numbers.
\end{proof}

The two classes are thus the endomorphisms of $U$ given by the elements of
$F$ of trace zero. A Hodge structure of K3 type with the same field sits
beside them, and the next lemma connects the two.

\begin{lemma}[The product map]\label{lem:mumfordmu}
For $x,y\in T$ let $\mu(x\otimes y)\in H^{2}(X,\QQ)$ be the class that
corresponds under $\psi$ to the alternating form
$(v,w)\mapsto\psi(xv,yw)+\psi(yv,xw)$ on $V$. Then $\mu$ is a $G$-equivariant
map $T\otimes T\to H^{2}(X,\QQ)$, and:
\begin{enumerate}[label=\textup{(\roman*)},nosep,leftmargin=2.6em]
\item for $i\ne j$ it restricts to an isomorphism $T_{i}\otimes T_{j}\to
  U_{ij}$, and it carries $T_{i}\otimes T_{i}$ onto $\QQ\theta$;
\item for $e,\lambda\in F^{\times}$ put $\mu_{e}=\mu\circ(m_{e}\otimes m_{e})$
  and let $\mu_{e}^{\dagger}$ be its adjoint for the form
  $q_{\lambda}\otimes q_{\lambda}$ on $T\otimes T$ and the Lefschetz pairing
  $B(\alpha,\beta)=\int_{X}\alpha\beta\,\theta^{2}$ on $H^{2}(X,\QQ)$. Then
  $\mu_{e}\mu_{e}^{\dagger}$ preserves $\QQ\theta$ and acts on $U$ as the
  element
  \[
    \nu\,\Nm_{F/\QQ}(g)\,g^{-1}\in F,\qquad g=e^{2}/\lambda,
  \]
  where $\nu$ is a nonzero rational number independent of $e$ and $\lambda$.
\end{enumerate}
\end{lemma}

\begin{proof}
Equivariance is clear, since $\mu$ is built from $\psi$ and the action of
$\End(V)$ on $V$. An element $x\in T_{i}$ acts on $V$ through the $i$-th factor
and lies in the symplectic Lie algebra of $\psi=\epsilon_{1}\otimes
\epsilon_{2}\otimes\epsilon_{3}$, so $\psi(xv,yw)=-\psi(v,xyw)$ and
$\mu(x\otimes y)$ is the form $-\psi(v,(xy+yx)w)$. If $x\in T_{i}$ and
$y\in T_{j}$ with $i\ne j$, then $xy=yx$ acts as $x\otimes y$ on
$V_{i}\otimes V_{j}$, and the form is $-2$ times the product of
$\epsilon_{i}(\cdot,x\,\cdot)$, $\epsilon_{j}(\cdot,y\,\cdot)$ and
$\epsilon_{k}$; the first two are symmetric because $x$ and $y$ lie in
$\mathfrak{sp}(\epsilon)=\mathfrak{sl}_{2}$, and $x\mapsto\epsilon(\cdot,x\,\cdot)$
is an isomorphism $\mathfrak{sl}_{2}\to S^{2}V^{\vee}$, so $\mu$ maps
$T_{i}\otimes T_{j}$ isomorphically onto $U_{ij}$. If $i=j$, then
$xy+yx=\operatorname{tr}(xy)$ on $V_{i}$, because $x$ and $y$ are trace-free
$2\times2$ matrices, and $\mu(x\otimes y)=-\operatorname{tr}(xy)\,\theta$.

For (ii), $m_{e}$ is self-adjoint for $q_{\lambda}$, and the adjoint of a
map out of $(T\otimes T,q_{\lambda}\otimes q_{\lambda})$ is
$(m_{\lambda}\otimes m_{\lambda})^{-1}$ times its adjoint for
$q_{1}\otimes q_{1}$, so
$\mu_{e}\mu_{e}^{\dagger}=\mu\,(m_{g}\otimes m_{g})\,\mu^{\dagger_{1}}$ with
$\mu^{\dagger_{1}}$ the adjoint for $q_{1}\otimes q_{1}$. This map is
$G$-equivariant, so it preserves $\QQ\theta$ and each $U_{ij}$, and
$\mu^{\dagger_{1}}$ carries $U_{ij}$ into
$(T_{i}\otimes T_{j})\oplus(T_{j}\otimes T_{i})$, on which $m_{g}\otimes m_{g}$
is multiplication by $\sigma_{i}(g)\sigma_{j}(g)=\Nm(g)/\sigma_{k}(g)$. Hence
$\mu_{e}\mu_{e}^{\dagger}$ is $\Nm(g)\sigma_{k}(g)^{-1}$ times
$\mu\mu^{\dagger_{1}}$ on $U_{ij}$. Item (XLI) computes $\mu\mu^{\dagger_{1}}$
on the split model and finds one and the same nonzero rational number $\nu$
on the three summands $U_{ij}$; it also checks the formula for
$\mu_{e}\mu_{e}^{\dagger}$ directly for several values of the conjugates of
$e$ and $\lambda$. By \Cref{prop:mumfordrm}(ii) an endomorphism of $U$ acting
on $U_{ij}$ by $\nu\Nm(g)/\sigma_{k}(g)$ is the element $\nu\Nm(g)g^{-1}$ of
$F$.
\end{proof}

By the Kuga--Satake class of a projective K3 surface $S$ we mean the Hodge
class on $S\times\mathrm{KS}(S)\times\mathrm{KS}(S)$ that induces the embedding
of \Cref{prop:ksembed} for $T(S)_{\QQ}$; its algebraicity is the Kuga--Satake
Hodge conjecture for $S$ of \Cref{rem:kstransport}.

\begin{theorem}[One K3 surface for the two classes]\label{thm:mumfordks}
In \Cref{setup:mumford}, let $\lambda\in F$ be totally positive.
\begin{enumerate}[label=\textup{(\roman*)},nosep,leftmargin=2.6em]
\item There is a projective K3 surface $S_{\lambda}$ of Picard number
  thirteen and an isomorphism of Hodge structures
  $T(S_{\lambda})_{\QQ}\cong T(-1)$ carrying the intersection form to
  $q_{\lambda}$. Its Kuga--Satake variety is isogenous to $X^{32}$.
\item If the Kuga--Satake class of $S_{\lambda}$ is algebraic, then the real
  multiplication by $F$ on $T(S_{\lambda})$ is induced by algebraic cycles on
  $S_{\lambda}\times S_{\lambda}$, and every Hodge class in
  $H^{4}(X\times X,\QQ)$ is algebraic; in particular the two classes of
  \Cref{prop:mumfordproduct}(ii) are algebraic.
\end{enumerate}
\end{theorem}

\begin{proof}
(i) The quadratic space $(T,q_{\lambda})$ has dimension nine and signature
$(2,7)$ by \Cref{prop:mumfordrm}(iii), and $H^{2}$ of a K3 surface with its
intersection form has dimension $22$ and signature $(3,19)$. Over $\QQ_{p}$
a nondegenerate quadratic space embeds isometrically in every nondegenerate
space of dimension at least three larger, and over $\RR$ the signatures
allow it, so by the Hasse--Minkowski theorem for representations of forms by
forms \cite[\S\S63, 66]{OMe00} there is an isometric embedding of
$(T,q_{\lambda})$ into $\Lambda_{\QQ}$, where $\Lambda$ is the K3 lattice.
The line $\ell=T(-1)^{2,0}\subset T_{\CC}\subset\Lambda_{\CC}$ satisfies
$(\ell,\ell)=0$ and $(\sigma,\bar{\sigma})>0$ for $0\ne\sigma\in\ell$, so by the
surjectivity of the period map there is a K3 surface $S_{\lambda}$ with a
marking $H^{2}(S_{\lambda},\ZZ)\cong\Lambda$ that carries
$H^{2,0}(S_{\lambda})$ to $\ell$ \cite[Chs.~6--7]{Huy16}. A rational class
is of type $(1,1)$ exactly when it is orthogonal to $\ell$, and its component
in $T$ is then a rational class of type $(1,1)$ in $T(-1)$, that is a
$G$-invariant vector of $T=\operatorname{Lie}G$; there is none, since $G$ is
semisimple. So $\NS(S_{\lambda})_{\QQ}=T^{\perp}$, of dimension $13$ and
signature $(1,12)$, and $T(S_{\lambda})_{\QQ}=T$ with the form $q_{\lambda}$.
A class of positive square in $\NS(S_{\lambda})$ makes $S_{\lambda}$
projective \cite[Ch.~8]{Huy16}.

For the Kuga--Satake variety, note first that over $\bar{\QQ}$ the space $T$
is the orthogonal sum of the three ternary spaces $T_{i}$, so $C(T)$ is the
graded tensor product of the $C(T_{i})$. Scaling a ternary form does not
change its even Clifford algebra, so $D^{1}$ is the spin group of
$(D^{0},\lambda\operatorname{trd})$ for every $\lambda$, and
$\Res_{F/\QQ}D^{1}$ maps to $\mathrm{Spin}(T,q_{\lambda})$ by multiplication
in the Clifford algebra, with the same kernel as its map to $\GL(V)$, so that
$G$ is a subgroup of $\mathrm{Spin}(T,q_{\lambda})$; as a representation of the $i$-th factor,
$C(T_{i})$ is four copies of $V_{i}$. Hence $C^{+}(T)$ is $V^{\oplus32}$ as a
representation of $G$ over $\bar{\QQ}$, and over $\QQ$ as well, because $V$ is
absolutely irreducible. The Hodge structures on $C^{+}(T)=H^{1}(\mathrm{KS}(S_{\lambda}),\QQ)$
and on $V^{\oplus32}=H^{1}(X^{32},\QQ)$ are both given by lifts of weight one of
the same homomorphism $\mathbb{S}\to\MT(T)_{\RR}$ through $G\cdot\mathbb{G}_{m}$, and
two such lifts coincide.\footnote{They differ by a homomorphism from the
anisotropic torus $\mathbb{S}/\mathbb{G}_{m}$ to the centre of
$G\cdot\mathbb{G}_{m}$, whose identity component is $\mathbb{G}_{m}$, and
such a homomorphism is trivial.} So
$H^{1}(\mathrm{KS}(S_{\lambda}),\QQ)\cong H^{1}(X^{32},\QQ)$ as Hodge
structures, and $\mathrm{KS}(S_{\lambda})$ is isogenous to $X^{32}$. For the
form Galluzzi uses this is her theorem \cite{Gal00}, and van Geemen proves the
corresponding splitting for every K3 surface with real multiplication
\cite{vG08}.

(ii) Write $S=S_{\lambda}$ and $A=\mathrm{KS}(S)$, and let $Z_{S}$ be an
algebraic cycle on $S\times A\times A$ inducing the Kuga--Satake embedding
$T(S)\to\End H^{1}(A)$. Composing $Z_{S}$ with the graph of an isogeny
$A\to X^{32}$ and with the graphs of the inclusions and projections of the
factors of $X^{32}$ gives algebraic cycles on $S\times X\times X$ inducing the
$32^{2}$ block components of the embedding, which are morphisms of Hodge
structures $T(S)\to\End(V)$, up to a Tate twist; both lie in the tensor
category generated by $H^{1}(X)$, so the blocks are equivariant for
$\MT(X)=G\cdot\mathbb{G}_{m}$. Over $\bar{\QQ}$,
$\End(V)=S^{2}V\oplus\bigwedge^{2}V$, and $T_{i}\cong S^{2}V_{i}$ occurs once
in $S^{2}V$, as $S^{2}V_{i}\otimes\bigwedge^{2}V_{j}\otimes\bigwedge^{2}V_{k}$,
and not at all in $\bigwedge^{2}V$. So each block takes values in $T$, and
through the identification of (i) it is a $G$-endomorphism of $T$, that is
$m_{e_{ab}}$ for some $e_{ab}\in F$ by \Cref{prop:mumfordrm}(iii). Over
$\bar{\QQ}$ the embedding is therefore $v\mapsto v\otimes N_{i}$ on $T_{i}$,
acting on $H^{1}(A)=V\otimes\QQ^{32}$, with
$N_{i}=(\sigma_{i}(e_{ab}))_{a,b}$. The embedding $\Psi$ of
\Cref{prop:ksembed} satisfies the Clifford relation
$\Psi(v)\Psi(w)+\Psi(w)\Psi(v)=2q_{\lambda}(v_{0},v_{0})\,q_{\lambda}(v,w)$,
because
$\Psi(v)\Psi(w)x=vwx\,v_{0}^{2}$. Take nonzero $v\in T_{i}$ and
$w\in T_{j}$ with $i\ne j$. As endomorphisms of $V$ they commute and $vw\ne0$,
so the left side is $vw\otimes(N_{i}N_{j}+N_{j}N_{i})$, while the right side
is zero because $T_{i}$ and $T_{j}$ are orthogonal; hence $N_{i}$ and $N_{j}$
anticommute. For $v\in T_{i}$ with $q_{\lambda}(v,v)\ne0$ the relation reads
$v^{2}\otimes N_{i}^{2}=q_{\lambda}(v_{0},v_{0})\,q_{\lambda}(v,v)$, and $v^{2}$
is a nonzero scalar on $V_{i}$, so $N_{i}^{2}$ is a nonzero scalar. Three pairwise
anticommuting invertible matrices are linearly independent: multiplying a
relation $\sum_{i}c_{i}N_{i}=0$ by $N_{j}$ on either side and adding gives
$2c_{j}N_{j}^{2}=0$. If the $e_{ab}$ spanned a subspace of $F$ of dimension at
most two, the $N_{i}$ would lie in the span of at most two rational matrices;
so the $e_{ab}$ span $F$, and two of them, $e$ and $e'$, are not rational
multiples of each other. Then $(e/e')^{2}$ is irrational, since a cubic field
contains no irrational element with rational square, so $e^{2}/\lambda$ and
$e'^{2}/\lambda$ are not both rational. Fix a block $\kappa=m_{e}$ with
$e^{2}/\lambda\notin\QQ$, induced by a cycle $\Gamma$. Item (XLI) builds the
Clifford algebra of the split model and checks that $\Psi$ has this shape,
with three anticommuting $N_{i}$.

The real multiplication first. The trace form $(a,b)\mapsto\operatorname{tr}(ab)$ on
$\End(V)$ is induced by an algebraic class: for correspondences $a,b$ on $X$
the Lefschetz trace formula gives
$\operatorname{tr}(a_{*}b_{*}\mid H^{1})=-\deg\bigl((a\circ b\circ\pi_{1})\cdot\Delta_{X}\bigr)$,
where $\pi_{1}$ is the K\"unneth projector onto $H^{1}$, algebraic on an
abelian variety \cite[\S2]{Kle68}. So the adjoint $\kappa^{\dagger}$ of
$\kappa$, for this form and the intersection form on $S$, is induced by a
cycle built from the transpose of $\Gamma$, $\pi_{1}$ and $\Delta_{X}$, and
$\kappa^{\dagger}\kappa$ is an algebraic endomorphism of $T(S)$. An element $x$ of $T_{i}$ acts on the
eight-dimensional space $V$ through a two-dimensional factor, with
multiplicity four, so $\operatorname{tr}_{V}(xy)=4\Tr_{F/\QQ}\operatorname{trd}(xy)$, and
\[
  q_{\lambda}(\kappa^{\dagger}\kappa x,y)
  =\operatorname{tr}_{V}(ex\cdot ey)
  =4\Tr_{F/\QQ}\bigl(e^{2}\operatorname{trd}(xy)\bigr)
  =q_{\lambda}\bigl(4e^{2}\lambda^{-1}x,y\bigr),
\]
that is $\kappa^{\dagger}\kappa=m_{4e^{2}/\lambda}$. By the choice of $e$ this
element is irrational, so it generates the cubic field $F$ and every $m_{a}$
is a polynomial in $\kappa^{\dagger}\kappa$. Composing with the projector of
$H^{2}(S)$ onto $T(S)$, which is algebraic because it is the diagonal minus
$[\mathrm{pt}\times S]$, $[S\times\mathrm{pt}]$ and $\sum_{r}d_{r}\times d_{r}^{\vee}$
over dual bases of $\NS(S)_{\QQ}$, gives algebraic cycles on $S\times S$
inducing the real multiplication.

The classes on $X\times X$ next. The map $\mu_{\kappa}=\mu\circ(\kappa\otimes\kappa)$
from $T(S)\otimes T(S)$ to $H^{2}(X,\QQ)$ is induced by an algebraic cycle on
$S\times S\times X$, obtained from $\Gamma\times\Gamma$, the divisor
$m^{*}\theta-p_{1}^{*}\theta-p_{2}^{*}\theta$ on $X\times X$, whose class is
$\psi$ in $H^{1}\otimes H^{1}$, and the diagonal of $X$ by pull-back,
push-forward and intersection, because $\mu$ uses nothing but $\psi$ and the
composition of endomorphisms of $V$.
The form $q_{\lambda}\otimes q_{\lambda}$ is the intersection form of
$S\times S$ on $T(S)\otimes T(S)$ and $B$ is the intersection form of $X$
twisted by $\theta^{2}$, so the adjoint $\mu_{\kappa}^{\dagger}$ is induced by
cup product with $\theta^{2}$ followed by the transposed cycle, and
$E=\mu_{\kappa}\mu_{\kappa}^{\dagger}$ is an algebraic endomorphism of
$H^{2}(X,\QQ)$. Since $\kappa=m_{e}$, \Cref{lem:mumfordmu}(ii) says that $E$
preserves $\QQ\theta$ and acts on $U$ as $\nu\Nm(g)g^{-1}$ with
$g=e^{2}/\lambda$, and this element of $F$ is not rational, since $g$ is
not. The projection $\pi_{\theta}$ onto $\QQ\theta$ is algebraic, being
given by the class $\theta\times\theta$ suitably normalised, so the algebraic
Hodge endomorphisms of $H^{2}(X,\QQ)$ contain $1$, $\pi_{\theta}$ and
$E(1-\pi_{\theta})$, which generate $\QQ\times F$, that is all of them by
\Cref{prop:mumfordrm}(ii). By the same part every Hodge class in
$H^{2}(X)\otimes H^{2}(X)$ is then algebraic. The Hodge classes of
$H^{4}(X\times X,\QQ)$ in the other K\"unneth components number one in each
of $H^{a}\otimes H^{4-a}$, $a\ne2$, by the count in the proof of
\Cref{prop:mumfordproduct}, and are products of divisor classes.
\end{proof}

\begin{figure}[tbp]
\centering
  \includegraphics[width=\textwidth]{fig_mumford}
\caption{\Cref{prop:mumfordrm} and \Cref{lem:mumfordmu}, drawn over the
triangle of the three factors of the Hodge group; every height is a
dimension. At the corners stand the summands $T_{i}=\mathfrak{sl}(V_{i})$ of
$T=\operatorname{Lie}G$, banded by Hodge type after the Tate twist $T(-1)$: blue for
$(1,1)$, clay for $(2,0)$ and amber for $(0,2)$, so that only $T_{3}$ carries
the K3 type. On the edges stand the three summands $U_{ij}$ of the primitive
part of $H^{2}(X)$, of dimension nine, $U_{12}$ of pure type $(1,1)$ and the
other two in three bands of three; in the middle is the line of $\theta$. The
solid arcs are the product map $\mu$, which carries $T_{i}\otimes T_{j}$ onto
$U_{ij}$, and the dashed ones carry $T_{i}\otimes T_{i}$ to $\QQ\theta$. The
violet arcs are the Galois group permuting the $U_{ij}$ as it permutes the
embeddings of $F$, which is why the two exceptional classes may be taken to be
the trace-zero part of a cubic field. Above, the K3 surface $S_{\lambda}$ of
\Cref{thm:mumfordks}, whose Kuga--Satake class, if algebraic, comes down to
$T$ as the map $\kappa=m_{e}$.}
\label{fig:mumford}
\end{figure}

\begin{remark}[What is known about the Kuga--Satake class]
\label{rem:mumfordks}
The Kuga--Satake class of $S_{\lambda}$ is not known to be algebraic, and
\Cref{thm:mumfordks} does not make it so. The known cases of the Kuga--Satake
Hodge conjecture lie elsewhere. They include the K3 surfaces of Picard number
at least seventeen, which Morrison relates to abelian surfaces \cite{Mor84},
and the double planes branched along six lines, of Picard number
at least sixteen, settled by Paranjape's cycle; see \cite{vG00}. Floccari
proves the conjecture for the K3 surfaces attached to projective sixfolds of
generalised Kummer type \cite{Flo24}, whose transcendental lattices embed,
after scaling, in the rational second cohomology of such a sixfold, of rank
seven and signature $(3,4)$; their transcendental rank is therefore at most
six, and the surfaces $S_{\lambda}$ have transcendental rank nine.
Schlickewei made the real multiplication algebraic on those double planes
that have it, necessarily by a field of degree at most two, from Paranjape's
cycle and van Geemen's description of their Kuga--Satake varieties
\cite{Sch10,vG08}, and Varesco showed that when the Kuga--Satake Hodge
conjecture holds for two hyperk\"ahler varieties, every Hodge similarity
between their transcendental lattices is algebraic up to the Lefschetz
isomorphism \cite{Var23}, which reaches a quadratic field through its square
root. A cubic field contains no irrational element with rational square, so
none of its elements acts as a similarity; the element
$\kappa^{\dagger}\kappa=m_{4e^{2}/\lambda}$ in the proof takes the place that
the similarities occupy there, and the Clifford relations are what make it
irrational. Buskin's theorem that rational Hodge isometries between K3
surfaces are algebraic \cite{Bus19} does not give the real multiplication
either. A Hodge isometry of $T(S_{\lambda})$ is a Hodge endomorphism, hence
some $m_{a}$ by \Cref{prop:mumfordrm}(iii), and it is an isometry only when
$a^{2}=1$. For $b\in F^{\times}$ the theorem does give algebraic cycles between
$S_{\lambda}$ and $S_{\lambda b^{-2}}$ that realise $m_{b}$, but every composite
of them that returns to $S_{\lambda}$ is again a Hodge isometry of
$T(S_{\lambda})$, so it multiplies by an element of square one.

The surfaces $S_{\lambda}$ include those of Zhu, who computes the integral
structure that the fourfold induces on $T$ and studies the K3 surfaces of
Picard number thirteen that carry it \cite{Zhu18}; its form is the
corestriction of the reduced norm of $D^{0}$, which up to sign is
$\frac12q_{1}$.
\Cref{thm:mumfordks} applies to them as it stands, because the irrationality
of $e^{2}/\lambda$ comes from the anticommuting blocks and not from $\lambda$.
The closest explicit families are those of van Geemen and Sch\"utt, with
real multiplication by the cyclic cubic fields $\QQ(\zeta_{7}+\zeta_{7}^{-1})$
and $\QQ(\zeta_{9}+\zeta_{9}^{-1})$ and Picard numbers four and ten, on which
the real multiplication is induced by explicit algebraic cycles
\cite[Theorem~1.2 and Section~4.8]{vGS25}; their Picard numbers are not
thirteen, and the Kuga--Satake class, which \Cref{thm:mumfordks} needs and
algebraic real multiplication does not supply, is not treated there.
The first case of (F2) is therefore implied by one instance of the
Kuga--Satake Hodge conjecture, for one K3 surface of Picard number thirteen,
and the argument stops there.
\end{remark}

\subsection{The Mumford target: rigidity, a numerical criterion, and the
bypasses}\label{ssec:mumfordrigid}

The first case of (F2) is now a propagation problem along one compact
Shimura curve: the exceptional classes are algebraic at every CM point
(\Cref{thm:mumfordcm}), and on every member exactly when $B(W\times_{C}W)$
holds (\Cref{cor:mumfordequiv}). This subsection asks whether the
propagation can be obtained without passing through the Lefschetz standard
conjecture, by one of the mechanisms that have closed other cases of the
Hodge conjecture: a degeneration, a deformation into a larger family on which
the class stays of Hodge type, a specialisation to special points, or
semiregularity. It proves that the first two cannot work, that the third
works exactly when the representing cycles have bounded degree, and it
reduces the fourth to one dimension count at one CM point. It does not prove
the classes algebraic at a very general point.

\begin{setupenv}[The classes $\omega_{a}$]\label{setup:omegaa}
Keep \Cref{setup:mumford} and let $W\to C$ be a family as in
\Cref{cor:mumfordlefschetz}, with fibres $X_{t}$. Over $\bar{\QQ}$ let
$\pi_{0},\pi_{12},\pi_{13},\pi_{23}\in H^{2}(X)\otimes H^{2}(X)\subset
H^{4}(X\times X)$ be the classes that correspond, under the isomorphism
$H^{2}\otimes H^{2}\cong\End H^{2}$ given by the form ${\textstyle\bigwedge^{2}}\psi$
on $H^{2}(X)={\textstyle\bigwedge^{2}}V^{\vee}$, to the projections onto $\QQ\theta$,
$U_{12}$, $U_{13}$ and $U_{23}$ of \Cref{prop:mumfordrm}(i), and for
$a=(a_{0},a_{1},a_{2},a_{3})$ put
\[
  \omega_{a}=a_{0}\pi_{0}+a_{1}\pi_{12}+a_{2}\pi_{13}+a_{3}\pi_{23},
  \qquad L_{0}(a)=a_{0}+a_{1}+a_{2}+a_{3}.
\]
The four tensors are invariant under $G$ and under the exchange $\tau$ of the
two factors of $X\times X$, because ${\textstyle\bigwedge^{2}}\psi$ is symmetric; so each
$\omega_{a}$ is a flat section over $C$ of Hodge type at every point. The
argument of \Cref{prop:mumfordrm}(ii) shows that the rational ones are those
with $a_{0}\in\QQ$ and $(a_{1},a_{2},a_{3})=(\sigma_{3}(\alpha),
\sigma_{2}(\alpha),\sigma_{1}(\alpha))$ for some $\alpha\in F$; we call such a
class \emph{exceptional} when $\alpha\notin\QQ$. The products of divisor
classes in $H^{2}\otimes H^{2}$ are the rational $\omega_{a}$ with
$\alpha\in\QQ$, and $\pi_{0}$ is a multiple of
$\mathrm{pr}_{1}^{*}\theta\cup\mathrm{pr}_{2}^{*}\theta$; so adding a rational
multiple of $\pi_{0}$ to an exceptional class changes $L_{0}$ and not the
question of algebraicity. For $c\in C$ write $Y=X_{c}\times X_{c}$, let
$HT^{2}(Y)=H^{0,2}(Y)\oplus H^{1}(T_{Y})\oplus H^{0}({\textstyle\bigwedge^{2}}T_{Y})$ act on
$H^{*}(Y,\CC)$ by contraction as in \Cref{setup:criterion}, of dimension
$28+64+28=120$, and let $\kappa\in H^{1}(T_{Y})$ be the Kodaira--Spencer class
at $c$ of the family $W\times_{C}W\to C$. Put
\[
  e(\omega)=\dim HT^{2}(Y)-\dim\Ann_{HT^{2}(Y)}(\omega)
  =120-\dim\Ann_{HT^{2}(Y)}(\omega).
\]
\end{setupenv}

\begin{lemma}[Linear relations among conjugates]\label{lem:conjugates}
Let $F$ be a cubic number field with embeddings $\sigma_{1},\sigma_{2},
\sigma_{3}$, let $\alpha\in F\setminus\QQ$, and let
$c_{0},c_{1},c_{2},c_{3},a_{0}\in\QQ$. The three conjugates
$\sigma_{i}(\alpha)$ are pairwise distinct and nonzero, and if
$c_{0}a_{0}+\sum_{i}c_{i}\sigma_{i}(\alpha)=0$ then $c_{1}=c_{2}=c_{3}$.
\end{lemma}

\begin{proof}
The minimal polynomial of $\alpha$ has degree three, since $[F:\QQ]$ is prime,
and is separable, so the conjugates are distinct; they are nonzero because
$\alpha\ne0$. Let $\Gamma$ be the Galois group of a normal closure; it
permutes the $\sigma_{i}$ transitively by $\gamma\circ\sigma_{i}=
\sigma_{\gamma(i)}$. Put $\beta=\alpha-\Tr_{F/\QQ}(\alpha)/3$, so that
$\beta\ne0$ has trace zero, and $r=-c_{0}a_{0}-(c_{1}+c_{2}+c_{3})
\Tr_{F/\QQ}(\alpha)/3\in\QQ$; the hypothesis reads $\sum_{i}c_{i}\sigma_{i}
(\beta)=r$. Applying every $\gamma\in\Gamma$ and averaging, each $\sigma_{j}$
occurs $|\Gamma|/3$ times in the place of each $\sigma_{i}$, so
$r=\frac13(\sum_{i}c_{i})\sum_{j}\sigma_{j}(\beta)=0$. The relations
$c\in\QQ^{3}$ with $\sum_{i}c_{i}\sigma_{i}(\beta)=0$ form a $\Gamma$-stable
subspace $N$ containing $(1,1,1)$. The plane $\{\sum c_{i}=0\}$ is an
irreducible representation of $\Gamma$ over $\QQ$, since $\Gamma$ contains
the $3$-cycles, which have no rational eigenvector in that plane; so if $N$
were larger than the line of $(1,1,1)$ it would contain the plane and be all
of $\QQ^{3}$, and $\sigma_{1}(\beta)=0$. Hence $c\in\QQ(1,1,1)$.\footnote{For a
cyclic cubic field $\Gamma$ is the group of $3$-cycles; otherwise it is the
full symmetric group. The argument uses only that $\Gamma$ contains the
$3$-cycles, which holds for every transitive subgroup of the symmetric group
on three letters.}
\end{proof}

\begin{lemma}[One exceptional class gives the others]\label{lem:oneexceptional}
Let $t\in C$. If one rational exceptional class $\omega$ is algebraic on
$X_{t}\times X_{t}$, then every Hodge class in $H^{2}(X_{t})\otimes
H^{2}(X_{t})$ is algebraic; in particular so are the two exceptional classes
of \Cref{prop:mumfordproduct}.
\end{lemma}

\begin{proof}
At a CM point every Hodge class of $X_{t}\times X_{t}$ is algebraic
(\Cref{thm:mumfordcm}); so let $t$ be another point, where the Hodge group is
$G$, and write $X=X_{t}$. For $x\in H^{2}(X)\otimes H^{2}(X)$ let $E_{x}$ be
the endomorphism $y\mapsto x_{*}(\theta^{2}\cup y)$ of $H^{2}(X,\QQ)$, where
$x_{*}$ is the action of $x$ as a correspondence. By Poincar\'e duality and
the hard Lefschetz isomorphism $\theta^{2}\colon H^{2}\to H^{6}$, the map
$x\mapsto E_{x}$ is a $G$-equivariant isomorphism onto $\End H^{2}(X,\QQ)$,
so it carries the Hodge classes onto the Hodge endomorphisms, which are
$\QQ\times F$ by \Cref{prop:mumfordrm}(ii). It carries the products of
divisor classes, which are invariant under $\Sp(V,\psi)$, onto the
$\Sp(V,\psi)$-equivariant endomorphisms, which are $\QQ\times\QQ$, so it
carries $\omega$ to an element $(r,\beta)$ with $\beta\notin\QQ$. The
algebraic Hodge classes form a subspace $A$ containing the products of divisor
classes, and $A$ is closed under composition: if $x_{1},x_{2}$ are algebraic,
then the composite correspondence $x_{1}\circ\Delta_{*}(\theta^{2})\circ
x_{2}$ is algebraic, and so is its K\"unneth component in
$H^{2}\otimes H^{2}$, since the K\"unneth projectors of $X\times X$ are
polynomials in the graphs of $[2]\times1$ and $1\times[2]$, which act on
$H^{i}(X)\otimes H^{j}(X)$ by $2^{i}$ and $2^{j}$ (as in Step 1 of the proof
of \Cref{thm:lefschetzfamily}); that component $x$ has
$E_{x}=E_{x_{1}}E_{x_{2}}$. So the image of $A$ is a subalgebra of
$\QQ\times F$ containing $\QQ\times\QQ$ and $(r,\beta)$, hence
$\QQ\times\QQ[\beta]=\QQ\times F$, because $\beta$ generates the cubic
field $F$.
\end{proof}

\begin{theorem}[Rigidity of the exceptional classes]\label{thm:mumfordrigid}
Let $c\in C$, $Y=X_{c}\times X_{c}$, and let $\omega=\omega_{a}$ be a
rational exceptional class.
\begin{enumerate}[label=\textup{(\roman*)},leftmargin=2.6em,itemsep=2pt]
\item $\Ann_{H^{1}(T_{Y})}(\omega)=\CC\kappa$ and
  $\Ann_{H^{0,2}(Y)}(\omega)=0$.
\item $\Ann_{H^{0}(\bigwedge^{2}T_{Y})}(\omega)=0$ if and only if
  $L_{0}(a)=a_{0}+\Tr_{F/\QQ}(\alpha)\ne0$. So
  $\Ann_{HT^{2}(Y)}(\omega)=\CC\kappa$ and $e(\omega)=119$ whenever
  $L_{0}(a)\ne0$, which can always be arranged by adding a rational multiple
  of $\pi_{0}$.
\item Let $\mathfrak{S}$ be the Siegel space of weight one Hodge structures on
  $H^{1}(Y,\QQ)$ polarised by $\psi\oplus\psi$, of dimension $36$. The germ at
  $Y$ of the locus in $\mathfrak{S}$ where $\omega$ stays of Hodge type is the
  germ of the curve $\{X_{t}\times X_{t}\}$ through $Y$.
\end{enumerate}
\end{theorem}

\begin{proof}
Over $\CC$ the Hodge decomposition at $c$ is
$H^{1,0}(X_{c})=V_{1}\otimes V_{2}\otimes V_{3}^{1,0}$, as at every point of
$C$, since the Hodge cocharacter lies in the third factor. A maximal torus
$S$ of $\SL(V_{1})\times\SL(V_{2})$ therefore preserves the Hodge
decomposition of $H^{1}(Y,\CC)=V\oplus V$ and acts on $HT^{2}(Y)$ and on
$H^{*}(Y,\CC)$, and so does $\tau$. Contraction is equivariant for both, and
$\omega$ is fixed by both, so $x\mapsto x\lrcorner\,\omega$ preserves the
decomposition of $HT^{2}(Y)$ by type ($H^{0,2}$, $H^{1}(T)$ or
$H^{0}(\bigwedge^{2}T)$, which land in the distinct summands $H^{2,4}$,
$H^{1,3}$ and $H^{0,2}$ of $H^{*}(Y)$), by weight of $S$, and by parity under
$\tau$. The annihilator is the sum of the kernels on the $46$ nonzero
blocks. Item (XLIV) of \Cref{sec:verification} writes each block as a matrix
whose entries are linear in $a$ and computes the Gram determinant of its
columns, a polynomial in $a$ that, for real $a$, vanishes exactly where the
columns become dependent. Every Gram determinant is a positive constant times a product of
the following factors, and the one block with a kernel at a generic point is
the $\tau$-even block of weight zero in $H^{1}(T)$, whose kernel is spanned,
for every $a$, by a vector independent of $a$, which is dropped before the
determinant is taken.
\begin{enumerate}[label=\textup{(\alph*)},leftmargin=2.6em,itemsep=1pt]
\item Eleven quadratic forms $P_{1},P_{2},Q_{1},\dots,Q_{4},R,S_{1},\dots,
  S_{4}$, each a positive combination of squares of rational linear forms
  whose common zeros have two of $a_{1},a_{2},a_{3}$ equal or one of them
  zero, for instance
  $2P_{1}=\sum_{i<j}(a_{i}-a_{j})^{2}$,
  $Q_{4}=(3a_{1}-a_{0}-a_{2}-a_{3})^{2}+16(a_{2}-a_{3})^{2}$ and
  $S_{2}=(a_{0}+a_{2})^{2}+3(a_{1}+a_{3})^{2}+8a_{2}^{2}$.
\item Six positive definite forms $Z_{1},\dots,Z_{6}$, each a positive
  combination of the squares of four independent linear forms.
\item The linear forms $a_{2}$, $a_{3}$, $a_{0}-3a_{1}+a_{2}-3a_{3}$,
  $a_{0}-3a_{1}-3a_{2}+a_{3}$, $a_{0}+9a_{1}-3a_{2}-3a_{3}$, whose
  coefficients on $a_{1},a_{2},a_{3}$ are not all equal.
\item The linear form $L_{0}$, which divides the Gram determinants of the
  nine $\tau$-odd blocks of $H^{0}(\bigwedge^{2}T)$ and no other.
\end{enumerate}
The sums of squares are certified in exact arithmetic by item (XLIV) and by
the Lean theorem \texttt{mumford\_rigidity\_sos}. For a rational exceptional
class the coefficients are real, and by \Cref{lem:conjugates} the factors of
(a), (b) and (c) do not vanish: those of (a) would force two conjugates of
$\alpha$ to coincide or one to vanish, those of (b) vanish only at $a=0$,
and those of (c) would give a rational relation among the conjugates with
unequal coefficients. So every block is injective except the one named, and
$L_{0}$ decides the bivector blocks: if $L_{0}(a)\ne0$ they are injective, and
if $L_{0}(a)=0$ the four blocks of dimension one, whose Gram determinants are
$L_{0}^{2}/16$, lie in the annihilator.

The vector of the exceptional block is $\kappa$. The family of Hodge
structures along $C$ is the orbit of the Hodge structure at $c$ under the
real points of the third factor of $G$, so the Kodaira--Spencer class
$\kappa$ is the component of type $(-1,1)$ of an element of
$\mathfrak{sl}(V_{3})$, the lowering operator acting on both copies of $V$.
It kills $\omega$ because $\omega$ is invariant under $G$, it is fixed by
$\tau$ and has weight zero for $S$, and item (XLIV) checks that it spans the
kernel. This proves (i) and (ii).

For (iii), a point of $\mathfrak{S}$ is a Hodge filtration $F^{1}$ on
$H^{1}(Y,\CC)$, and $\omega$ is of Hodge type at it when $\omega\in F^{2}
H^{4}$, a holomorphic condition. For weight one every tangent direction of
$\mathfrak{S}$ at $Y$ is an element $\xi$ of
$\mathfrak{sp}(\psi\oplus\psi)^{-1,1}\subset H^{1}(T_{Y})$, and it moves
$F^{2}H^{4}$ by the derivation $\xi$ of ${\textstyle\bigwedge^{*}}H^{1}$, which maps
$H^{2,2}$ into $H^{1,3}$; so the Zariski tangent space of the Hodge locus is
$\{\xi:\xi\lrcorner\,\omega=0\}$, which is $\CC\kappa$ by (i).\footnote{This is
Griffiths transversality in its simplest case: for a variation of weight one
there is no horizontality condition, so the tangent space of the Hodge locus
of a class of type $(p,p)$ is the kernel of the contraction into
$H^{p-1,p+1}$, see \cite[Chapter~5]{Voi03}. For a class of higher weight on a
general variety only the horizontal directions enter.} The Hodge locus is an
analytic germ containing the smooth curve germ $\{X_{t}\times X_{t}\}$, whose
tangent line is $\CC\kappa$, and its Zariski tangent space is that line; so
the two germs are equal.
\end{proof}

\begin{remark}[What the divisor classes do instead]\label{rem:rigidcompare}
The same computation at classes with $\alpha\in\QQ$ shows what rigidity
excludes. At $a=(3,5,5,5)$, a combination of products of divisor classes, the
tangent space of the Hodge locus in $\mathfrak{S}$ has dimension $10$, and it
contains the tangent space of the diagonal Siegel space
$\{X'\times X'\}$ of dimension $10$, because that class is invariant under
the whole of $\Sp(V,\psi)$; at $a=(1,1,1,1)$ it has dimension $20$. So the
products of divisor classes move with $X'\times X'$ for every $X'$, the
exceptional classes move with $X_{t}\times X_{t}$ for $t\in C$ only, and every
deformation of $Y$ off the curve, however small, makes them cease to be of
Hodge type. \Cref{fig:rigidity} draws the three loci.
\end{remark}

\begin{figure}[tbp]
\centering
  \includegraphics[width=\textwidth]{fig_rigidity}
\caption{\Cref{thm:mumfordrigid}. Left: the Siegel space $\mathfrak{S}$ of
dimension $36$ near $Y=X_{c}\times X_{c}$ and three Hodge loci through it,
labelled by the tangent dimension $T$ at $Y$: the class with equal
coefficients ($T=20$), the products of divisor classes with $a_{1}=a_{2}=a_{3}$
($T=10$, the diagonal Siegel space), and a rational exceptional class ($T=1$),
whose locus is the compact Mumford curve with its CM points (hollow). The
dashed rim is the boundary of $\mathfrak{S}$, which the curve does not reach.
Right: the nine weight blocks of $\mathfrak{sp}^{-1,1}$ under the torus of
the first two factors, with their dimensions and the forms whose product is
the Gram determinant of each; only the block of weight $(0,0)$ has a kernel,
the direction $\kappa$ of the curve. The same holds in the $46$ blocks of
$HT^{2}(Y)$ off the hyperplane $L_{0}=0$. Below: the sums of squares that
make the determinants nonzero at the conjugates of an irrational element of
$F$.}
\label{fig:rigidity}
\end{figure}

\begin{corollary}[No degeneration and no larger family]
\label{cor:nodegeneration}
Let $\omega$ be a rational exceptional class on $X_{c}\times X_{c}$. Every
connected complex analytic family of polarised abelian eightfolds through
$X_{c}\times X_{c}$ on which $\omega$ stays of Hodge type lies, near each
point of the curve, inside the curve $\{X_{t}\times X_{t}:t\in C\}$, which is
compact. In particular no such family degenerates, and on no member of it
other than the CM points of the curve is $\omega$ a combination of products
of divisor classes or of Weil classes of CM fields acting on the member or on
its quotients.
\end{corollary}

\begin{proof}
By \Cref{thm:mumfordrigid}(iii) the Hodge locus of $\omega$ coincides with the
curve in a neighbourhood of every point of the curve, so a connected family
through $Y$ on which $\omega$ stays of Hodge type maps into the curve near
each of its points over the curve; by connectedness and the identity theorem
its image is the curve. The curve is compact because $D$ is a division
algebra, so the arithmetic group of the Shimura curve is cocompact
\cite{Mum69}.\footnote{By \Cref{setup:mumford} the algebra $D$ is definite at
$\sigma_{1}$ and $\sigma_{2}$, so $D\otimes_{F,\sigma_{1}}\RR$ is Hamilton's
quaternions and $D$ is not $M_{2}(F)$: it is a division algebra, split at the
one real place $\sigma_{3}$. The norm one group of an order in it, embedded in
$\SL_{2}(\RR)$ through $\sigma_{3}$, is a Fuchsian group of finite covolume
with no parabolic element: a parabolic $\gamma\ne\pm1$ would make
$\gamma\mp1$ a nonzero nilpotent element of $D$. A Fuchsian group of finite
covolume without parabolic elements has compact quotient.} For the last sentence: a member of the curve other than a CM
point has Hodge group $G$ (proof of \Cref{cor:mumfordequiv}), and on it
$\omega$ lies outside the span of the products of divisor classes and of the
Weil classes of quotients, by \Cref{prop:mumfordproduct}(ii) and (iii). At a
CM point the class is algebraic already (\Cref{thm:mumfordcm}), and the
countably many CM points do not carry algebraicity to the other members
(\Cref{lem:spreading}).
\end{proof}

The same computation controls deformations that are not algebraic at all,
and so decides whether the twistor lines along which Markman moves sheaves
for Weil classes can be used here.

\begin{proposition}[The Hodge locus among all complex tori, and twistor
lines]\label{prop:notwistor}
Let $\omega$ be a rational exceptional class, $Y=X_{c}\times X_{c}$ with
$c\in C$, and $\Lambda=H_{1}(Y,\ZZ)$. Let $\mathcal{T}$ be the space of
complex structures on the real torus $\Lambda_{\RR}/\Lambda$, a complex
manifold of dimension $64$ with tangent space $H^{1}(T_{Y})$ at $Y$, and
$\mathcal{T}_{\omega}\subset\mathcal{T}$ the closed analytic subset on which
$\omega$ is of type $(2,2)$.
\begin{enumerate}[label=\textup{(\roman*)},leftmargin=2.6em,itemsep=2pt]
\item The connected component of $\mathcal{T}_{\omega}$ that contains $Y$ is
  the universal cover $\mathfrak{H}$ of the Mumford curve, the set of
  squares of members of the family, biholomorphic to the upper half plane.
  Every connected complex analytic family of complex tori, algebraic or not,
  through a member of the Mumford family on which $\omega$ stays of Hodge
  type lies in the family.
\item No twistor line through a member of the Mumford family, and no compact
  connected complex curve in $\mathcal{T}$ through it, keeps $\omega$ of
  Hodge type.
\item $\mathcal{T}_{\omega}$ contains twistor lines nevertheless. The algebra
  $D\otimes_{F,\sigma_{1}}\RR$ is Hamilton's quaternions and acts on
  $V_{\RR}$ through the first factor of $G(\RR)$; for a unit quaternion $j$
  with $j^{2}=-1$, the complex structure $j\oplus j$ on $V_{\RR}\oplus
  V_{\RR}$ makes $\omega$ of type $(2,2)$, and these structures form a twistor
  line. On each factor $\psi(x,jy)$ has signature $(4,4)$, and a member whose
  Mumford--Tate group is $G$ is not an abelian variety. These twistor lines
  lie in components of $\mathcal{T}_{\omega}$ other than $\mathfrak{H}$.
\end{enumerate}
\end{proposition}

\begin{proof}
(i) The complex structures on $\Lambda_{\RR}$ for which all the tensors
invariant under $G$ are of Hodge type are those $J$ with $J^{2}=-1$ in the
real points of $G\cdot\mathbb{G}_{m}$, since a reductive group containing the
homotheties is the stabiliser of its invariant tensors. Let $\mathfrak{H}$ be
the connected component of this closed set that contains $Y$; it is closed in
$\mathcal{T}$ and consists of the $J=1\otimes1\otimes j_{3}$ with
$j_{3}\in\SL_{2}(\RR)$, $j_{3}^{2}=-1$, in the half plane that contains the
structure of $Y$, that is of the squares of the members of the Mumford
family. At each of its points the
first order condition for $\omega$ to stay of type $(2,2)$ along
$\xi\in H^{1}(T_{Y})$ is $\xi\lrcorner\,\omega=0$, as in the proof of
\Cref{thm:mumfordrigid}(iii), where no polarisation was used, and by
\Cref{thm:mumfordrigid}(i) the solutions form the line $\CC\kappa$ tangent to
$\mathfrak{H}$. So near each of its points $\mathcal{T}_{\omega}$ is the smooth
curve germ of $\mathfrak{H}$, and $\mathfrak{H}$ is open in
$\mathcal{T}_{\omega}$. Being also closed and connected, it is a connected
component. A connected analytic family on which $\omega$ stays of Hodge type
maps into $\mathcal{T}_{\omega}$, hence into $\mathfrak{H}$ once one of its
points does.

(ii) A twistor line is a nonconstant holomorphic image of $\PP^{1}$ in
$\mathcal{T}$. If $\omega$ stayed of Hodge type along one through a point of
$\mathfrak{H}$, or along any compact connected curve through it, the curve
would lie in $\mathfrak{H}$ by (i), and a holomorphic map from a compact
connected curve to the upper half plane is constant.

(iii) At $\sigma_{1}$ the quaternion algebra $D$ is definite
(\Cref{setup:mumford}), so the first factor of $G(\RR)$ is the group of unit
quaternions, and $V_{\RR}$ restricted to it is a module over
$D\otimes_{F,\sigma_{1}}\RR\cong\mathbb{H}$. The unit quaternions $j$ with
$j^{2}=-1$ form a two-sphere, and the complex structures given by left
multiplication by them on a quaternionic vector space form a twistor line.
Each $j$ lies in $G(\RR)$, and so does the circle $\cos\vartheta+j\sin\vartheta$,
which fixes $\omega$; so $\omega$ is of type $(2,2)$. Item (XLVII) of
\Cref{sec:verification} realises $V_{\RR}$ inside $V_{1}\otimes V_{2}
\otimes V_{3}$ as the fixed space of the real structure that is quaternionic on
the first two factors and real on the third, and finds that $\psi(x,Jy)$ is
definite for the structures of the Mumford family and of signature $(4,4)$ for
$J=j\otimes1\otimes1$. At a member whose Mumford--Tate group is $G$, the Hodge
classes of degree two on $V_{\RR}\oplus V_{\RR}$ are the forms $\psi\otimes s$
with $s$ a symmetric $2\times2$ matrix, by \Cref{prop:mumfordproduct}(ii), and
the signature of $\psi(\cdot,J\cdot)\otimes s$ is $(4(p+q),4(p+q))$ when $s$ has
signature $(p,q)$, so none is definite, none is a polarisation, and the member
is not an abelian variety. The members of the Mumford family are not on these
lines, so by (i) the lines lie in other components.
\end{proof}

Markman's proof for Weil classes on abelian fourfolds moves a sheaf along
twistor lines on which its Chern character stays of Hodge type. By
\Cref{prop:notwistor} no such line passes through the square of a member of
the Mumford family: the twistor lines on which $\omega$ is of Hodge type exist,
but they carry non-algebraic tori whose Hodge cocharacter lies in a definite
factor of $D$, in components of the Hodge locus that the Mumford family never
meets. The Kuga--Satake class of \Cref{thm:mumfordks} behaves differently. The
Kuga--Satake construction applies to every Hodge structure of K3 type on
$(T,q_{\lambda})$, so that class stays of Hodge type over the whole period
domain of $(T,q_{\lambda})$, of dimension seven, and not only along the curve;
it is the one class attached to the target that deformation methods of this
kind are not excluded from.

The same space $T$ also makes the square a holomorphic symplectic variety,
and this decides what the mechanism of van Geemen and Rapagnetta would have
to supply here. They prove the Hodge conjecture for the general abelian
fourfold of Weil type with trivial discriminant by pulling back the second
Chern class of a hyperk\"ahler sixfold along a map from the fourfold, which
gives an algebraic class of codimension two that is not a product of divisor
classes \cite{vGR26}. For a projective hyperk\"ahler manifold $M$ write
$\sigma_{M}$ for its symplectic form, $q_{M}$ for its Beauville--Bogomolov
form and $T(M)$ for the orthogonal complement of $\NS(M)_{\QQ}$ in
$H^{2}(M,\QQ)$ for $q_{M}$.

\begin{proposition}[The square as a holomorphic symplectic
variety]\label{prop:hksquare}
Let $c\in C$ be a point at which the Mumford--Tate group of $X=X_{c}$ is
$G\cdot\mathbb{G}_{m}$, let $Y=X\times X$, and let $\Psi\in\bigwedge^{2}V$ be
the $G$-invariant bivector dual to $\psi$. For $x\in T$ put
$\iota(x)=(x\otimes1)\Psi$ in $V\otimes V=H^{1}(X,\QQ)\otimes H^{1}(X,\QQ)$, a
K\"unneth component of $H^{2}(Y,\QQ)$, and let $\iota_{2}\colon S^{2}T\to
H^{4}(Y,\QQ)$ be the map $xy\mapsto\iota(x)\cup\iota(y)$. Let $C_{i}\in
S^{2}T_{i}$ be the Casimir element of the trace form of
$\mathfrak{sl}(V_{i})$, and for $\nu\in F$ let
$c_{\nu}=\sum_{i}\sigma_{i}(\nu)C_{i}$, a rational element of $S^{2}T$; for
$\nu\ne0$ it is the element dual to the form $q_{1/\nu}$.
\begin{enumerate}[label=\textup{(\roman*)},leftmargin=2.6em,itemsep=2pt]
\item $\iota$ is an embedding of Hodge structures $T(-1)\to H^{2}(Y,\QQ)$, its
  image is the only irreducible sub-Hodge structure of $H^{2}(Y,\QQ)$ with
  $h^{2,0}=1$, and the holomorphic $2$-form $\sigma_{0}$ spanning its part of
  type $(2,0)$ is symplectic.
\item $\iota_{2}(c_{\nu})=\omega_{a}$ with $a_{0}=3\Tr_{F/\QQ}(\nu)$ and
  $\alpha=4\nu-\Tr_{F/\QQ}(\nu)$, in the notation of \Cref{setup:omegaa}. So
  $\iota_{2}(c_{\nu})$ is exceptional if and only if $\nu\notin\QQ$, and every
  rational exceptional class is some $\iota_{2}(c_{\nu})$ plus a rational
  multiple of $\pi_{0}$.
\item Let $M$ be a projective hyperk\"ahler manifold and
  $f\colon Y\dashrightarrow M$ a rational map with $f^{*}\sigma_{M}\ne0$. Then
  $f^{*}$ maps $T(M)$ isomorphically onto $\iota(T)$ and is compatible with
  products of classes in $T(M)$; $f$ is generically finite onto a subvariety
  of dimension eight on which $\sigma_{M}$ is generically nondegenerate;
  $b_{2}(M)\ge10$; and $\dim M\ge10$.
\item Suppose in (iii) that $M=S^{[n]}$ is the Hilbert scheme of $n$ points
  on a projective K3 surface $S$, and
  let $q_{\lambda}$ be the form on $T$ that corresponds to the intersection
  form of $S$ on $T(S)=T(M)$ under $\iota^{-1}f^{*}$
  (\Cref{prop:mumfordrm}(iii)). If $\lambda\notin\QQ$, or if the real
  multiplication of $T(S)$ is induced by algebraic cycles on $S\times S$,
  then the Hodge conjecture holds for $Y$.
\end{enumerate}
\end{proposition}

\begin{proof}
(i) The map $\iota$ is $G$-equivariant because $\Psi$ is invariant, and
injective. Its image lies in $S^{2}V$: for $x$ in the symplectic Lie algebra
of $\psi$ one has $(x\otimes1)\Psi=-(1\otimes x)\Psi$, and the exchange of the
two factors carries $(x\otimes1)\Psi$ to $-(1\otimes x)\Psi$ because $\Psi$
is alternating. An element $g$ of $G\cdot\mathbb{G}_{m}$ multiplies $\Psi$ by
the multiplier $\chi(g)$ of $\psi$, so $(g\otimes g)\iota(x)=
\chi(g)\,\iota(gxg^{-1})$: the map $\iota$ is equivariant for the adjoint
action on $T$ twisted by $\chi$, which is the action defining $T(-1)$. Both
Hodge structures come from the Hodge cocharacter through
$G\cdot\mathbb{G}_{m}$, so $\iota$ is a morphism of Hodge structures from
$T(-1)$. The sub-Hodge structures
of $H^{2}(Y,\QQ)$ are its $G$-stable subspaces, since $\mathbb{G}_{m}$ acts
by scalars. Over $\bar{\QQ}$, item (XLVIII) of \Cref{sec:verification} splits
$H^{2}(Y)=\bigwedge^{2}(V\oplus V)$ by the three Casimir operators into the
isotypic pieces of highest weights $(0,0,0)$, $(2,2,0)$, $(2,0,2)$,
$(0,2,2)$, $(2,2,2)$, $(2,0,0)$, $(0,2,0)$, $(0,0,2)$, of dimensions
$3,27,27,27,27,3,3,3$, whose parts of type $(2,0)$ have dimensions
$0,0,9,9,9,0,0,1$; the last three pieces are $\iota(T_{1})$, $\iota(T_{2})$,
$\iota(T_{3})$. So an irreducible summand of highest weight $(2,0,2)$ or
$(0,2,2)$ contributes $3$ to $h^{2,0}$, one of weight $(2,2,2)$ contributes
$9$, one of weight $(0,0,2)$ contributes $1$, and the others contribute
nothing. A rational sub-Hodge structure $W$ is stable under the Galois group,
which permutes the highest weights as it permutes the factors, transitively.
So if $W\otimes\bar{\QQ}$ has a summand of weight $(2,2,0)$, $(2,0,2)$ or
$(0,2,2)$, it has summands of all three and $h^{2,0}(W)\ge6$; if it has one of
weight $(2,2,2)$, then $h^{2,0}(W)\ge9$. If $h^{2,0}(W)=1$, therefore,
$W\otimes\bar{\QQ}$ contains the piece $\iota(T_{3})$ of weight $(0,0,2)$,
which has multiplicity one, and with it $\iota(T_{1})$ and $\iota(T_{2})$; so
$W\supseteq\iota(T)$, and $W=\iota(T)$ if $W$ is irreducible. The
part of type $(2,0)$ of $\iota(T)$ is spanned by $\sigma_{0}=\iota(E_{3})$,
where $E_{3}\in\mathfrak{sl}(V_{3})\otimes\CC$ maps $V_{3}^{0,1}$ onto
$V_{3}^{1,0}$, and item (XLVIII) finds that $\sigma_{0}$ pairs $H^{1,0}$ of the
first factor with $H^{1,0}$ of the second through an invertible matrix. So
$\sigma_{0}$ has rank eight: it is a holomorphic symplectic form on $Y$.

(ii) Item (XLVIII) computes
$\iota_{2}(C_{1})=3\pi_{0}-\pi_{12}-\pi_{13}+3\pi_{23}$ and the two cyclic
permutations of this identity, in the split model; both sides are built
from $\psi$ and the actions of the three factors, so the identities hold over
$\bar{\QQ}$ in \Cref{setup:mumford}. Summing with the coefficients
$\sigma_{i}(\nu)$ gives $\iota_{2}(c_{\nu})=3\Tr(\nu)\pi_{0}+
\sum\sigma_{k}(4\nu-\Tr(\nu))\pi_{ij}$, the sum over $\{i,j,k\}=\{1,2,3\}$ with
$i<j$, which is $\omega_{a}$ with the stated $a_{0}$ and $\alpha$. Since
$\Tr(\nu)\in\QQ$, the element $\alpha$ is rational exactly when $\nu$ is; and
$\Tr(\alpha)=\Tr(\nu)$, so $\nu=(\alpha+\Tr(\alpha))/4$ recovers $\nu$ from
any $\alpha\in F$. Two rational classes with the same $\alpha$ differ by a
multiple of $\pi_{0}$. The Casimir element $C_{i}$ is dual to the trace form
on $T_{i}$, so $c_{\nu}$ is dual to $\sum_{i}\sigma_{i}(\nu)^{-1}
\operatorname{tr}(x_{i}y_{i})=q_{1/\nu}(x,y)$.

(iii) Resolve the indeterminacy of $f$ by a sequence of blow-ups along
smooth centres, $\pi\colon\tilde{Y}\to Y$ and $\tilde{f}\colon\tilde{Y}\to M$,
and put $f^{*}=\pi_{*}\tilde{f}^{*}$, a morphism of Hodge structures. The
kernel of $\pi_{*}$ on $H^{2}(\tilde{Y},\QQ)$ is spanned by the classes of
the exceptional divisors, so $\alpha\mapsto\tilde{f}^{*}\alpha-
\pi^{*}f^{*}\alpha$ is a morphism of Hodge structures from $H^{2}(M,\QQ)$ to a
space of Hodge classes. The Hodge structure $T(M)$ is irreducible. A
sub-Hodge structure $W\subseteq T(M)$ with $W^{2,0}=0$ consists of rational
classes of type $(1,1)$, which lie in $\NS(M)_{\QQ}\cap T(M)=0$; if
$W^{2,0}\ne0$, then $\sigma_{M}\in W_{\CC}$, the orthogonal complement of $W$
in $H^{2}(M,\QQ)$ is a sub-Hodge structure orthogonal to $\sigma_{M}$ and
$\bar{\sigma}_{M}$, hence of type $(1,1)$ and contained in $\NS(M)_{\QQ}$,
and $W\supseteq T(M)$. It is not spanned by Hodge classes, so that
morphism vanishes on it, and $\tilde{f}^{*}\alpha=\pi^{*}f^{*}\alpha$ for
$\alpha\in T(M)$. Since $\tilde{f}^{*}$ is a ring homomorphism and
$\pi_{*}\pi^{*}$ is the identity, $f^{*}(\alpha\beta)=f^{*}\alpha\cup
f^{*}\beta$ for $\alpha,\beta\in T(M)$. The kernel of $f^{*}$ on $T(M)$ is a
sub-Hodge structure not containing $\sigma_{M}$, hence zero, and the image is
irreducible with $h^{2,0}=1$, hence $\iota(T)$ by (i). So $f^{*}\sigma_{M}$ is
a nonzero multiple of $\sigma_{0}$ and has rank eight at every point of $Y$
where $f$ is defined; the rank of a pulled-back form is at most the dimension
of the image, so $f$ is generically finite onto its image $Z$, of dimension
eight, and $\sigma_{M}$ is generically nondegenerate on $Z$. As $M$ is
projective, $b_{2}(M)=\dim T(M)+\dim\NS(M)_{\QQ}\ge9+1$; this excludes the
generalised Kummer varieties, whose second Betti number is seven
\cite{Bea83}. Finally $\dim M\ge\dim Z=8$, and suppose $\dim M=8$. Then $f$
is dominant, and for $\alpha\in T(M)_{\CC}$ the relation of Fujiki and
Beauville \cite{Bea83,Fuj87} gives
\[
  \int_{Y}(f^{*}\alpha)^{8}=\int_{\tilde{Y}}\tilde{f}^{*}(\alpha^{8})
  =\deg(f)\int_{M}\alpha^{8}=\deg(f)\,c_{M}\,q_{M}(\alpha)^{4},
  \qquad c_{M}>0 .
\]
Write $f^{*}\alpha=\iota(x)$ with $x=x_{1}+x_{2}+x_{3}\in T_{\CC}$ and
$N_{i}=-\det x_{i}$. The form $x\mapsto q_{M}((f^{*})^{-1}\iota(x))$ is
invariant under $G$, so it is $q_{\lambda}(x,x)=
\sum_{i}\sigma_{i}(\lambda)\operatorname{tr}(x_{i}^{2})=
2\sum_{i}\sigma_{i}(\lambda)N_{i}$ for some $\lambda\in F^{\times}$ by
\Cref{prop:mumfordrm}(iii). On the other side, item (XLVIII) finds
$\iota(x)^{8}=8!\det(x)$ times a volume form, and
$\det(x_{1}\otimes1\otimes1+1\otimes x_{2}\otimes1+1\otimes1\otimes
x_{3})=\Delta(N_{1},N_{2},N_{3})^{2}$ with
$\Delta=N_{1}^{2}+N_{2}^{2}+N_{3}^{2}-2N_{1}N_{2}-2N_{1}N_{3}-2N_{2}N_{3}$;
this is also the product of the eight numbers $\pm r_{1}\pm r_{2}\pm r_{3}$,
where $\pm r_{i}$ are the eigenvalues of $x_{i}$ and $r_{i}^{2}=N_{i}$. The
$N_{i}$ are algebraically independent functions on $T_{\CC}$, so
$\Delta^{2}$ would be a constant multiple of the fourth power of the linear
form $\sum_{i}\sigma_{i}(\lambda)N_{i}$. But $\Delta$ is a nondegenerate
ternary quadratic form, hence irreducible, and a contradiction results. So
$\dim M\ge10$.

(iv) By \cite{Bea83} the map $H^{2}(S,\QQ)\to H^{2}(S^{[n]},\QQ)$ induced by
the universal subscheme $Z_{n}\subset S\times S^{[n]}$ is injective, carries
the intersection form to the Beauville--Bogomolov form, and has
$H^{2}(S^{[n]},\QQ)=H^{2}(S,\QQ)\oplus\QQ\delta$ with $\delta$ algebraic, so
$T(M)=T(S)$. Let $\Delta_{T}$ be the K\"unneth component in $T(S)\otimes
T(S)$ of the class of the diagonal of $S$; it is algebraic, being the class
of the diagonal minus $[\mathrm{pt}\times S]$, $[S\times\mathrm{pt}]$ and
$\sum_{r}d_{r}\times d_{r}^{\vee}$ over dual bases of $\NS(S)_{\QQ}$. Its image
under the correspondence $Z_{n}\times Z_{n}$, pulled back along the diagonal
of $M$, is the class $\sum_{a}e_{a}\cup e_{a}^{\vee}$ over dual bases of
$T(S)$ for the intersection form, which is therefore algebraic on $M$, and by
(iii) its image under $f^{*}$ is $\iota_{2}$ of the element of $S^{2}T$ dual
to $q_{\lambda}$, that is $\iota_{2}(c_{1/\lambda})$, an algebraic class on
$Y$. If $\lambda\notin\QQ$ it is exceptional by (ii). If instead the real
multiplication is algebraic, apply it to the first factor of $\Delta_{T}$:
for $b\in F$ the class $\sum_{a}(m_{b}e_{a})\cup e_{a}^{\vee}$ is algebraic,
and its image is $\iota_{2}(c_{b/\lambda})$, which is exceptional for
$b\notin\QQ\lambda$. In both cases \Cref{lem:oneexceptional} makes every
Hodge class in $H^{2}(X)\otimes H^{2}(X)$ algebraic; the Hodge classes in the
other K\"unneth components of $H^{4}(Y,\QQ)$ are products of divisor classes,
as at the end of the proof of \Cref{thm:mumfordks}; and the Hodge classes of
$Y$ form a ring generated in degrees two and four (proof of
\Cref{cor:mumfordequiv}).
\end{proof}

\begin{remark}[What the mechanism of van Geemen and Rapagnetta would
need]\label{rem:hkroute}
By \Cref{prop:hksquare}(ii) the exceptional classes are the classes
$\iota_{2}(c_{\nu})$ with $\nu\notin\QQ$, the images of the elements dual to
the forms $q_{\lambda}$, $\lambda=1/\nu$, on the part $\iota(T)$ of $H^{2}(Y)$
generated by the symplectic class; the form with $\lambda$ rational gives a
product of divisor classes. They are the analogues
for $(Y,\sigma_{0})$ of the dual Beauville--Bogomolov class, whose pull-back is
what \cite{vGR26} uses. Part (iii) says what a transfer of that argument
needs: a projective hyperk\"ahler manifold of dimension at least ten with
transcendental part $T(-1)$, such as $S_{\mu}^{[n]}$ with $n\ge5$ for a
surface $S_{\mu}$ of \Cref{thm:mumfordks}, and a rational map from $Y$ onto a
symplectic subvariety of dimension eight. For $M=S_{\mu}^{[n]}$ the form of
(iv) is $q_{\mu/e^{2}}$ for the $e\in F^{\times}$ with $\iota^{-1}f^{*}=m_{e}$
on $T(S_{\mu})=T$, so $\lambda\notin\QQ$ for every such map as soon as
$\mu\notin\QQ^{\times}F^{\times2}$. No such map is known. It would give an
algebraic cycle on $S_{\mu}\times X\times X$ inducing an isomorphism of
$T(S_{\mu})$ onto $\iota(T)\subset H^{1}(X)\otimes H^{1}(X)$, a correspondence
of the kind that the Kuga--Satake class supplies in the proof of
\Cref{thm:mumfordks}; so this is the Kuga--Satake route in geometric form,
not an independent one. Maps to K3 surfaces and to hyperk\"ahler eightfolds
cannot serve, by (iii).
\end{remark}

\begin{theorem}[The Mumford criterion is a dimension]
\label{thm:mumfordnumerical}
Let $c\in C$, $Y=X_{c}\times X_{c}$, $\omega=\omega_{a}$ a rational
exceptional class with $L_{0}(a)\ne0$, and $E$ a perfect complex on $Y$ with
$\ch(E)=N\omega$, $N\ne0$, and no other component.
\begin{enumerate}[label=\textup{(\roman*)},leftmargin=2.6em,itemsep=2pt]
\item $\dim\Ext^{2}(E,E)\ge119$.
\item If $\dim\Ext^{2}(E,E)=119$, then $\operatorname{ev}_{E}\colon
  HT^{2}(Y)\to\Ext^{2}(E,E)$ is surjective with kernel $\CC\kappa$, and
  $\sigma_{E}$ is injective. Conversely, if $\sigma_{E}$ is injective, then
  $\operatorname{ev}_{E}$ has rank $119$ and kernel $\CC\kappa$, and $E$
  extends to first order along the curve.
\item If some such $E$ at one point $c$ has $\dim\Ext^{2}(E,E)=119$, then
  $\omega_{t}$ is algebraic for every $t\in C$, the Hodge conjecture holds for
  $X_{t}\times X_{t}$ for every $t$, and $B(W\times_{C}W)$ holds.
\item At every CM point $c$ some perfect complex $E$ on $Y$ has
  $\ch(E)=N\omega$ for some $N\ne0$ and no other component.
\end{enumerate}
\end{theorem}

\begin{proof}
(i) and (ii). By the square $\sigma_{E}\circ\operatorname{ev}_{E}=
(\,\cdot\,)\lrcorner\,\ch(E)$ of \Cref{setup:criterion}, $\ker
\operatorname{ev}_{E}\subseteq\Ann_{HT^{2}}(\omega)=\CC\kappa$ by
\Cref{thm:mumfordrigid}(ii), so $\operatorname{ev}_{E}$ has rank at least
$119$, which is (i). If $\dim\Ext^{2}(E,E)=119$, then
$\operatorname{ev}_{E}$ is onto, its kernel has dimension one and lies in
$\CC\kappa$, so equals it; and $\sigma_{E}(\operatorname{ev}_{E}(x))=
N\,x\lrcorner\,\omega$ vanishes only for $x\in\CC\kappa=\ker
\operatorname{ev}_{E}$, so $\sigma_{E}$ is injective on
$\image\operatorname{ev}_{E}=\Ext^{2}(E,E)$. Conversely, if $\sigma_{E}$ is
injective, then $\sigma_{E}(\operatorname{ev}_{E}(\kappa))=N\kappa\lrcorner
\,\omega=0$ forces $\operatorname{ev}_{E}(\kappa)=0$, so the kernel is
$\CC\kappa$ and the rank is $119$; and $\operatorname{ev}_{E}$ restricted to
$H^{1}(T_{Y})$ is the obstruction map \cite{Tod09}, which vanishes at
$\kappa$.

(iii) By (ii) $\sigma_{E}$ is injective. The Chern character of $E$ is
$N\omega$, a flat section of Hodge type over the whole of $C$, so
\Cref{thm:BF} applies to every infinitesimal deformation of $Y$ inside
$W\times_{C}W\to C$. The proof of \Cref{thm:semiregclosure} now applies word
for word, with the stack of perfect complexes on the fibres of $W\times_{C}
W\to C$ \cite{Lie06} in place of that of $\cA\to\cS$, or the analytic
Kuranishi germ if $\Ext^{<0}(E,E)\ne0$: it is smooth over $C$ at $[E]$, its
image contains a nonempty open subset $U\subseteq C$, and for $t\in U$ a
perfect complex $E_{t}$ on $X_{t}\times
X_{t}$ in the same connected component has $\ch(E_{t})=N\omega_{t}$, so
$\omega_{t}$ is algebraic, and by \Cref{lem:oneexceptional} so are both
exceptional classes. $U$ is uncountable, and \Cref{cor:mumfordequiv} gives
the rest.

(iv) By \Cref{thm:mumfordcm} there is $\xi\in\CH^{2}(Y)_{\QQ}$ with class
$\omega$, and the argument of \Cref{prop:ktheory} gives $E$: the Chern
character is an isomorphism $K_{0}(Y)_{\QQ}\to\CH^{*}(Y)_{\QQ}$
\cite[Ex.~15.2.16(b) and Cor.~18.3.2]{Ful98}, so $\xi=\ch(u)$ for a unique
$u$, and $Nu$ is the class of a perfect complex for a suitable $N$.
\end{proof}

Part (iii) is the Mumford analogue of \Cref{thm:p2numerical}(iii): one object,
at one point, and one number, which is $119$ here and $2n(2n-1)$ for a Weil
family. For a Weil family that number is attained by no object
(\Cref{thm:p2false}), because an object attaining it would have to deform to
non-algebraic tori; the classes here are rigid, so that argument does not
apply (\Cref{rem:p2prime}). The objects of (iv) exist, and the only open question is whether one
of them, or another representative of the same class in $K_{0}$, attains the
bound of (i). We have not found one, and we do not claim that one exists. The
rigidity of \Cref{thm:mumfordrigid} is not an obstacle to it: semiregularity
moves an object along the Hodge locus of its Chern character, and the curve
is exactly that locus, so a semiregular object would propagate along the
curve and nowhere else, which is all that is needed.

\begin{corollary}[Bounded data at CM points]\label{cor:mumfordbounded}
Fix a relatively ample line bundle on $W\times_{C}W\to C$, and let $\omega$ be
a rational exceptional class. The following are equivalent.
\begin{enumerate}[label=\textup{(\alph*)},leftmargin=2.6em,itemsep=2pt]
\item $\omega_{t}$ is algebraic for every $t\in C$; equivalently, by
  \Cref{lem:oneexceptional} and \Cref{cor:mumfordequiv}, the Hodge conjecture
  holds for $X_{t}\times X_{t}$ for every $t\in C$.
\item There are a finite set $\Phi$ of pairs $(\underline{P},\underline{q})$,
  each a tuple of Hilbert polynomials with a tuple of rational numbers of the
  same length, and infinitely many points $c\in C$ at each of which
  $\omega_{c}=\sum_{i}q_{i}[Z_{i}]$ for subschemes $Z_{i}\subset X_{c}\times
  X_{c}$ with Hilbert polynomials $P_{i}$, for some $(\underline{P},
  \underline{q})\in\Phi$.
\end{enumerate}
In particular the classes are algebraic on every member as soon as they are
represented at infinitely many CM points by cycles of bounded degree.
\end{corollary}

\begin{proof}
For a pair $(\underline{P},\underline{q})$ let $\Sigma_{\underline{P},
\underline{q}}\subseteq C$ be the set of points at which $\omega$ is
represented with these data. The proof of \Cref{thm:closedlocus} applies
verbatim to $(W\times_{C}W,C,\omega)$, as in \Cref{lem:spreading}, and shows
that each $\Sigma_{\underline{P},\underline{q}}$ is Zariski closed. If (b)
holds, one pair of $\Phi$ serves at infinitely many points, so its
$\Sigma$ is an infinite closed subset of the irreducible curve $C$, hence
$C$, and (a) holds. If (a) holds, $C$ is the union of the countably many
$\Sigma_{\underline{P},\underline{q}}$, and $C$ is uncountable, so one of
them is infinite, hence equal to $C$; then (b) holds with $\Phi$ that one
pair and every point of $C$.
\end{proof}

\Cref{thm:mumfordcm} gives a representation at every CM point, through
Markman's theorem for abelian fourfolds \cite{Mar25a} and the ring of
\Cref{thm:mumfordcm}, with no control of degree, and isogenies between CM
members, which move CM points along their Hecke orbits, do not bound it
either. \Cref{cor:mumfordbounded} says that such a bound is not a convenience
but the whole remaining content: with it the target falls, and without it the
CM points, being countable, say nothing about a very general member
(\Cref{lem:spreading} needs uncountably many points).

\begin{remark}[The bypass mechanisms, one by one]\label{rem:bypassaudit}
\Cref{fig:bypass} collects what is known about each route to the Mumford
target that avoids the Lefschetz standard conjecture.
\begin{enumerate}[label=\textup{(\arabic*)},leftmargin=2.6em,itemsep=2pt]
\item \emph{Degeneration.} Blocked: the base is compact and the Hodge locus is
  the base (\Cref{cor:nodegeneration}), so no family on which $\omega$ stays
  of Hodge type approaches the boundary, and there is no limit mixed Hodge
  structure to specialise to.
\item \emph{Deformation into a larger Hodge locus}, on which $\omega$ would
  become a product of divisor classes or a Weil class and be algebraic by
  an earlier theorem, then be carried back. Blocked by
  \Cref{thm:mumfordrigid}(iii), already to first order.
\item \emph{Hochschild deformations}: twisting $D^{b}(Y)$ by a $B$-field or a
  Poisson bivector, as in the gerby and noncommutative deformations used for
  Weil classes. Blocked: an object with $\ch(E)=N\omega$ extends along no
  direction of $HH^{2}(Y)$ except that of the curve
  (\Cref{cor:mumfordsplit}(i)).
\item \emph{Constructions from line bundles}: cones, tensor products, duals,
  pull-backs and push-forwards along homomorphisms, and Fourier--Mukai
  functors whose kernels are built in the same way. Blocked at every point of
  the curve: everything so built has a Chern character invariant under the
  Lefschetz group, and the exceptional classes are not
  (\Cref{thm:lefschetzclosure}).
\item \emph{Twistor and hyperholomorphic deformation}, the mechanism of
  Markman's proof for Weil classes: move a sheaf along twistor lines on which
  its Chern character stays of Hodge type. Blocked on the square: no twistor
  line through a member of the Mumford family keeps $\omega$ of Hodge type,
  and those that do carry non-algebraic tori, in other components of the
  Hodge locus (\Cref{prop:notwistor}).
\item \emph{Specialisation.} The class is algebraic at every CM point
  (\Cref{thm:mumfordcm}), on every power of $X_{c}$ \cite[Theorem~1.3]{Li26};
  this suffices exactly when the degrees are bounded
  (\Cref{cor:mumfordbounded}), and no bound is known.
\item \emph{Reduction modulo $p$.} At every closed point of every reduction
  of the curve the class is a $\QQ_{\ell}$-combination of algebraic cycles
  (\Cref{prop:modp}), by Li's theorem \cite[Theorem~1.1]{Li26}; this
  suffices, with no lifting step, exactly when the Hilbert data of the cycles
  are bounded on a Zariski dense set of closed points of the arithmetic model
  (\Cref{thm:modpdensity}), and no bound is known.
\item \emph{Semiregularity.} Reduced to one dimension count at one CM point
  (\Cref{thm:mumfordnumerical}); an object attaining it has large odd
  self-extensions (\Cref{thm:mumfordprofile}), is not built from line
  bundles (\Cref{prop:mumforddivisor}) and splits only in a constrained way
  (\Cref{cor:mumfordsplit}); no such object is known.
\item \emph{Kuga--Satake.} Reduced to the Kuga--Satake class of one K3 surface
  of Picard number thirteen (\Cref{thm:mumfordks}), which the known cases do
  not reach (\Cref{rem:mumfordks}). Its geometric form is the mechanism of
  van Geemen and Rapagnetta \cite{vGR26}, pull-back of the dual
  Beauville--Bogomolov class of a hyperk\"ahler manifold: the exceptional
  classes are exactly such twisted dual forms of the part of $H^{2}$ of the
  square generated by its symplectic class, and a rational map onto a
  symplectic subvariety of some $S_{\mu}^{[n]}$, $n\ge5$, would suffice;
  K3 surfaces and hyperk\"ahler eightfolds cannot serve
  (\Cref{prop:hksquare}, \Cref{rem:hkroute}).
\item \emph{Motivated classes.} Every Hodge class on an abelian variety is
  motivated \cite[Th\'eor\`eme~0.6.2]{And96b}, so $\omega_{t}$ is motivated for
  every $t$; a motivated class is algebraic once the Lefschetz involutions of
  the auxiliary varieties are (\Cref{prop:lefmot}), and for the auxiliary
  varieties of Andr\'e's proof that is not known. Along the family the auxiliary variety
  that suffices is $W\times_{C}W$ (\Cref{thm:lefschetzpropagation}), and its
  Lefschetz standard conjecture is equivalent to the target
  (\Cref{cor:mumfordequiv}): the route moves the problem and does not shorten
  it.
\end{enumerate}
Five of the ten are closed off by theorems; two work at special points
only, the CM points and the closed points of the reductions modulo $p$, and
each of them suffices exactly when a degree bound holds; two are reduced to a
single object, one Kuga--Satake class or one perfect complex with
$\dim\Ext^{2}=119$ at one CM point; and one is equivalent to the target. None
of them is carried out here, and the target remains open. The five that are
not closed off are one statement and two stronger inputs, not five
obligations: any one of them would suffice (\Cref{cor:mumfordone},
\Cref{rem:mumfordneeded}).
\end{remark}

\begin{figure}[tbp]
\centering
  \includegraphics[width=0.96\textwidth]{fig_bypass}
\caption{\Cref{rem:bypassaudit}: the ten routes to the Mumford target that
avoid the Lefschetz standard conjecture. Clay, with a crossed dashed spoke: a
route closed off by a theorem of \Cref{ssec:mumfordrigid} or
\Cref{ssec:mumfordobject}. Grass: a route that works at special points only,
at the CM points or at the closed points of the reductions modulo $p$, and
suffices exactly when a degree bound holds (\Cref{cor:mumfordbounded},
\Cref{thm:modpdensity}). Indigo: a route reduced to one open statement, of
the same kind as the target or equivalent to it.}
\label{fig:bypass}
\end{figure}

\subsection{What an object meeting the Mumford criterion must look like, and
  the target modulo $p$}
\label{ssec:mumfordobject}

\Cref{thm:mumfordnumerical} reduces the semiregularity route to one object at
one CM point. This subsection derives what any object with the right Chern
character must satisfy, whether or not it is semiregular: no sequence of
standard operations produces it from line bundles, it extends along no first
order deformation of $D^{b}(Y)$ other than the direction of the curve, a
semiregular one splits only into summands with no $\Ext^{2}$ between them,
and its self-extensions are large in every degree, with an odd part that the
Euler characteristic forces beyond what Hochschild cohomology supplies. It
then turns to the reduction modulo $p$, after the recent proof of the Tate
conjecture for powers of abelian fourfolds over finite fields: the classes
are represented by cycles at every closed point of every reduction of the
curve, and the target is equivalent to a bound on those cycles on a Zariski
dense set of closed points. None of this constructs the object or closes the
target.

\begin{lemma}[The contraction ranks are symmetric]\label{lem:rankduality}
Let $A$ be an abelian variety of dimension $g$ and $c\in H^{*}(A,\CC)$ a class
each of whose components has Hodge type $(p,p)$ for some $p$. For $k\ge0$ put
$r_{k}(c)=\dim\bigl(HT^{k}(A)\lrcorner\,c\bigr)$, with the action of
\Cref{setup:hh}. Then $r_{k}(c)=r_{g-k}(c)$ for $0\le k\le g$, and
$r_{k}(c)=0$ for $k>g$.
\end{lemma}

\begin{proof}
Write $H^{1,0}=H^{1,0}(A)$ and $A'=H^{0,1}(A)$, both of dimension $g$, so
that $H^{*}(A,\CC)=\bigwedge^{*}H^{1,0}\otimes\bigwedge^{*}A'$ and
$HT^{*}(A)=\bigwedge^{*}(A'\oplus T)$ with $T=(H^{1,0})^{\vee}$. Fix a
generator of $\bigwedge^{g}A'$ and let
$\star\colon\bigwedge^{q}A'\to\bigwedge^{g-q}A'^{\vee}$ send $\eta$ to
$\zeta\mapsto\eta\wedge\zeta$. Under $\Phi=\mathrm{id}\otimes\star$, wedging
with $z\in A'$ becomes, up to a sign fixed on each bigraded piece, contraction
with $z$ on the second factor, while $T$ still contracts the first. So $\Phi$
identifies the action of $HT^{*}(A)$ with the action of
$\bigwedge^{*}U^{\vee}$ by interior products on $\bigwedge^{*}U$, where
$U=H^{1,0}\oplus A'^{\vee}$ has dimension $2g$, up to signs that do not change
ranks. A class of type $(p,p)$ goes to
$\bigwedge^{p}H^{1,0}\otimes\bigwedge^{g-p}A'^{\vee}\subset\bigwedge^{g}U$,
so $\eta=\Phi(c)$ lies in the middle degree whatever the $p$ of its
components. Hence $r_{k}(c)$ is the rank of
$\bigwedge^{k}U^{\vee}\to\bigwedge^{g-k}U$, $x\mapsto\iota_{x}\eta$, which is
zero for $k>g$. Its transpose for the natural pairings is
$\bigwedge^{g-k}U^{\vee}\to\bigwedge^{k}U$, $y\mapsto\pm\iota_{y}\eta$,
because $\langle\iota_{x}\eta,y\rangle=\langle\eta,x\wedge y\rangle
=\pm\langle\iota_{y}\eta,x\rangle$. A map and its transpose have the same
rank.
\end{proof}

\begin{theorem}[The self-extensions of an object with an exceptional Chern
character]\label{thm:mumfordprofile}
Let $c\in C$ and $Y=X_{c}\times X_{c}$, and write
$e_{k}=\dim\Ext^{k}(E,E)$ for a perfect complex $E$ on $Y$.
\begin{enumerate}[label=\textup{(\roman*)},leftmargin=2.6em,itemsep=2pt]
\item For every $E$ and every $k$, $e_{k}\ge r_{k}(\ch E)$ and
  $e_{k}=e_{8-k}$.
\item For a rational exceptional class $\omega=\omega_{a}$ with
  $L_{0}(a)\ne0$: $r_{0}=r_{8}=1$, $r_{1}=r_{7}=16$ and $r_{2}=r_{6}=119$.
\item There is a set of rational exceptional classes, Zariski dense in the
  space of rational Hodge classes in $H^{2}\otimes H^{2}$, on which
  \[
    \bigl(r_{0},\dots,r_{8}\bigr)(\omega)\;\ge\;
    (1,16,119,328,560,328,119,16,1)
  \]
  term by term. When $F$ is the cubic subfield of $\QQ(\zeta_{7})$ it contains
  the class with $a_{0}=2$ and $(a_{1},a_{2},a_{3})$ the conjugates of
  $\zeta_{7}+\zeta_{7}^{-1}$.
\item Let $\omega$ be in the set of (iii) and $\ch(E)=N\omega$ with $N\ne0$.
  Then
  \[
    e_{1}+e_{3}=e_{0}+e_{2}+\tfrac12e_{4}
      +\sum_{j\ge1}(-1)^{j}e_{-j}.
  \]
  If $\Ext^{<0}(E,E)=0$, for instance if $E$ is a sheaf, then
  $\sum_{k\ \mathrm{odd}}e_{k}=\sum_{k\ \mathrm{even}}e_{k}\ge800$,
  $e_{1}+e_{3}\ge400$, and $\sum_{k}e_{k}\ge1600$, while the bounds of
  (iii) add up to $1488$: the odd degrees carry at least $112$ dimensions
  beyond the Hochschild bound $2(16+328)=688$.
\end{enumerate}
\end{theorem}

\begin{proof}
(i) The square $\sigma_{E}\circ\operatorname{ev}_{E}=(\,\cdot\,)\lrcorner
\ch(E)$ of \Cref{setup:hh} holds in every degree, so the kernel of
$\operatorname{ev}_{E}$ on $HT^{k}(Y)$ lies in the annihilator of $\ch(E)$
and its image has dimension at least $r_{k}(\ch E)$. Serre duality with the
trivial canonical bundle gives $e_{k}=e_{8-k}$.

(ii) $r_{0}=1$ because $\omega\ne0$, and $r_{2}=119$ is
\Cref{thm:mumfordrigid}(ii). If $w\in HT^{1}(Y)$ annihilates $\omega$, so does
$w'\wedge w$ for every $w'\in HT^{1}(Y)$, since the annihilator is a left
ideal; for $w\ne0$ these products span a space of dimension $15$, which cannot
lie in $\Ann_{HT^{2}}(\omega)=\CC\kappa$. So $r_{1}=16$, and
\Cref{lem:rankduality} gives $r_{8},r_{7},r_{6}$.

(iii) For each $k$ the set of $a\in\CC^{4}$ at which the contraction on
$HT^{k}(Y)$ into $\omega_{a}$ has rank at least the displayed number is Zariski
open, being the complement of the common zeros of minors, and item (XLV) of
\Cref{sec:verification} shows that it contains $a=(3,5,-7,11)$, by an exact
computation over $\QQ$ block by block. The rational Hodge classes in
$H^{2}\otimes H^{2}$ form a $\QQ$-vector space of dimension four whose complex
span is $\CC^{4}$ (\Cref{setup:omegaa}), hence Zariski dense, and so they
meet the intersection of these open sets in a Zariski dense set; the classes
with $\alpha\in\QQ$ form a proper linear subspace and are removed. Item (XLV)
also computes the ranks modulo a prime at which the minimal polynomial of
$\zeta_{7}+\zeta_{7}^{-1}$ splits, for $a_{0}=2$ and the three conjugates in
each of the six orders, and finds the displayed numbers; a rank modulo a prime
of good reduction is a lower bound for the rank in characteristic zero.

(iv) Since $\ch(E)$ has only a component in $H^{4}$, $\ch(E^{\vee})\ch(E)$
lies in $H^{8}(Y)$ and $\chi(E,E)=\int_{Y}\ch(E^{\vee})\ch(E)=0$, the Todd
class being trivial. By (i) and Serre duality, which also pairs
$\Ext^{-j}$ with $\Ext^{8+j}$,
\[
  0=\chi(E,E)=2(e_{0}-e_{1}+e_{2}-e_{3})+e_{4}
  +2\sum_{j\ge1}(-1)^{j}e_{-j},
\]
which is the identity. With no negative extensions the odd and even parts are
equal, and the even part is at least $2\cdot1+2\cdot119+560=800$ by (i) and
(iii), so the odd part is at least $800$ as well, and $e_{1}+e_{3}\ge400$ by
Serre duality. The Hochschild bound on the odd part is $2(16+328)=688$.
\end{proof}

\begin{figure}[tbp]
\centering
  \includegraphics[width=\textwidth]{fig_extprofile}
\caption{\Cref{thm:mumfordprofile}. Left: the lower bounds $r_{k}$ for
$\dim\Ext^{k}(E,E)$ when $\ch(E)=N\omega$, one bar per degree, symmetric about
$k=4$ as \Cref{lem:rankduality} predicts; the bar at $k=2$ is the $119$ of
\Cref{thm:mumfordrigid}, made of $28$ classes in $H^{0,2}$, $63$ in
$H^{1}(T)$ and $28$ bivectors. On the hyperplane $L_{0}=0$ the bars at $2$,
$4$ and $6$ drop to $110$, $548$ and $110$ at a sample point. Right: $\chi(E,E)=0$ balances the
even and the odd degrees, so the odd side, which Hochschild cohomology fills
only to $688$, carries at least $800$.}
\label{fig:extprofile}
\end{figure}

\begin{proposition}[Divisor classes do not reach the exceptional classes]
\label{prop:mumforddivisor}
Let $t\in C$ and let $\omega$ be a rational exceptional class. Then $\omega_{t}$
is not in the $\QQ$-span of products of divisor classes of $X_{t}\times
X_{t}$: that span meets the rational Hodge classes of $H^{2}\otimes H^{2}$
exactly in the classes with $\alpha\in\QQ$. Consequently no perfect complex
whose Chern character lies in the subring generated by divisor classes, in
particular no bounded complex of direct sums of line bundles, has
$\ch(E)=N\omega$ with $N\ne0$. At a CM point $c$ the fourfold $X_{c}$ has eight
Hodge classes in $H^{4}$, of which the products of divisor classes span six.
\end{proposition}

\begin{proof}
At a point with Hodge group $G$, the divisor classes of $X_{t}\times X_{t}$ are
the three classes of \Cref{prop:mumfordproduct}(ii), all invariant under
$\Sp(V,\psi)$, and so are their products; the $\Sp(V,\psi)$-invariant part of
$H^{2}\otimes H^{2}$ is spanned by $\pi_{0}$ and $\pi_{12}+\pi_{13}+\pi_{23}$,
the classes with $a_{1}=a_{2}=a_{3}$. At a CM point the Hodge group is a
maximal torus of $G$ (proof of \Cref{thm:mumfordcm}), so the Hodge classes
are the monomials of weight zero, and item (XLV) computes over $\QQ$: $16$
divisor classes, $132$ Hodge classes in degree four, products of divisor
classes spanning $100$ of them, and a span meeting that of the four $\pi$ in
exactly the classes with $a_{1}=a_{2}=a_{3}$. For a rational class these
coordinates are the conjugates of $\alpha$, which by \Cref{lem:conjugates} are
equal only if $\alpha\in\QQ$. The Chern character of a complex in the subring
generated by divisor classes has its component in $H^{4}$ in the span of
products of two divisor classes, which excludes $N\omega$. The last sentence
is the same computation for $X_{c}$.
\end{proof}

At a CM point the Hodge classes of $X_{c}$ in $H^{4}$ that are not products
of divisor classes are among the classes that \Cref{thm:mumfordcm} pulls back
along homomorphisms $X_{c}\times X_{c}\to X_{c}$; they are algebraic by
Markman's theorem \cite{Mar25a}. Li proves,
independently and more generally, that every Hodge class on every power of
every complex CM abelian fourfold is algebraic \cite[Theorem~1.3]{Li26},
which gives \Cref{thm:mumfordcm} for every power $X_{c}^{n}$ and not only the
square. \Cref{prop:mumforddivisor} says that an object with $\ch(E)=N\omega$
cannot have its Chern character in the subring generated by divisor classes,
so it cannot be assembled from line bundles alone. The next theorem shows that
no sequence of the standard operations of the derived category, applied to
line bundles on powers of $X_{t}$, gets any further.

\begin{theorem}[What line bundles generate]\label{thm:lefschetzclosure}
Let $t\in C$, $X=X_{t}$, and let $L(X)\subseteq\Sp(V,\psi)$ be the Lefschetz
group of $X$, the centraliser of $\End^{0}(X)$ in $\Sp(V,\psi)$. Let
$\mathcal{A}$ be the class of abelian varieties isogenous to a power of $X$;
through an isogeny $X^{n}\to B$ the group $L(X)$ acts diagonally on
$H^{*}(B,\QQ)$, and the action does not depend on the isogeny. Let
$\mathcal{L}$ be the smallest class of objects of the categories
$D^{b}(B)$, $B\in\mathcal{A}$, that contains the line bundles and is closed
under shifts, cones, derived tensor products, derived duals, translations,
the functors $Lf^{*}$ and $Rf_{*}$ for homomorphisms $f$ between members of
$\mathcal{A}$, and the Fourier--Mukai functors $\Phi_{K}\colon D^{b}(B)\to
D^{b}(B')$ with kernel $K\in\mathcal{L}$ on $B\times B'$.
\begin{enumerate}[label=\textup{(\roman*)},leftmargin=2.6em,itemsep=2pt]
\item Every object of $\mathcal{L}$ has an $L(X)$-invariant Chern character.
\item $L(X)=\Sp(V,\psi)$ if $t$ is not a CM point, and $L(X)$ is a maximal
  torus of $\Sp(V,\psi)$ if it is. In both cases the $L(X)$-invariant classes
  on $X\times X$ form the subring generated by divisor classes, of dimensions
  $1,3,6,10,15,10,6,3,1$, respectively $1,16,100,304,454,304,100,16,1$, in
  the degrees $0,2,\dots,16$.
\item No rational exceptional class $\omega_{t}$ is $L(X)$-invariant. Hence
  for no object $E\in\mathcal{L}$ on $X\times X$ and no $N\ne0$ is
  $\ch(E)-N\omega_{t}$ invariant under $L(X)$; in particular
  $\ch(E)\ne N\omega_{t}$.
\item The $\Sp(V,\psi)$-module generated by $\pi_{12}-\pi_{13}$ is
  irreducible, of highest weight $2\varpi_{2}$ and dimension $308$, and it
  meets the span of the four $\pi$ in the plane of the classes with
  $a_{0}=0$ and $a_{1}+a_{2}+a_{3}=0$.
\end{enumerate}
\end{theorem}

\begin{proof}
(i) Two isogenies $X^{n}\to B$ differ by an element of $\End^{0}(X^{n})=
M_{n}(\End^{0}(X))$, which commutes with $L(X)$ acting diagonally, so the
action is well defined. For the same reason a homomorphism $f$ between members
of $\mathcal{A}$ is $L(X)$-equivariant on $H^{1}$, so $f^{*}$ is equivariant
on $H^{*}$, and so is $f_{*}$, which is $f^{*}$ conjugated by Poincar\'e
duality, since $L(X)\subseteq\SL(V)$ fixes the orientation classes.
Translations act trivially on cohomology. The Chern character is additive on
distinguished triangles and multiplicative on derived tensor products,
$\ch(E^{\vee})=\sum_{i}(-1)^{i}\ch_{i}(E)$, $\ch(Lf^{*}E)=f^{*}\ch(E)$, and
$\ch(Rf_{*}E)=f_{*}\ch(E)$ by the Grothendieck--Riemann--Roch theorem
\cite[Theorem~15.2]{Ful98}, the Todd class of an abelian variety being $1$;
and $\Phi_{K}=Rp_{2*}(p_{1}^{*}(\,\cdot\,)\otimes^{L}K)$. So the objects with
$L(X)$-invariant Chern character form a class with all the closure
properties, and it remains to see that it contains the line bundles, that is,
that $c_{1}(\mathcal{O}(D))$ is $L(X)$-invariant. Pulled back along an
isogeny, which is injective on rational cohomology, $c_{1}$ becomes a Hodge
class of degree two on some $X^{n}$, an element of the invariants of the
Hodge group in ${\textstyle\bigwedge^{2}}(V^{\vee}\otimes\QQ^{n})=
{\textstyle\bigwedge^{2}}V^{\vee}\otimes S^{2}\QQ^{n}\oplus S^{2}V^{\vee}
\otimes{\textstyle\bigwedge^{2}}\QQ^{n}$. If $t$ is not a CM point the Hodge
group is $G$, whose invariants in ${\textstyle\bigwedge^{2}}V$ are the multiples of $\psi$
and in $S^{2}V$ are zero, since every summand of $S^{2}(V_{1}\otimes V_{2}
\otimes V_{3})$ contains a factor $S^{2}V_{i}$; so the Hodge classes of degree
two are $\Sp(V,\psi)$-invariant. If $t$ is a CM point, a monomial of degree two
in a weight basis has weight zero for the Hodge torus exactly when its two
weights are opposite, and then it has weight zero for $L(X)$, by the
description in (ii).

(ii) If $t$ is not a CM point, $\End^{0}(X)=\QQ$ and $L(X)=\Sp(V,\psi)$. At a
CM point $\End^{0}(X)\otimes\CC$ is the algebra of diagonal matrices in a
weight basis (proof of \Cref{thm:mumfordcm}), and $\psi$ pairs each weight
with its opposite only, so $L(X)_{\CC}$ is the torus of diagonal matrices
$\operatorname{diag}(\lambda_{w})$ with $\lambda_{w}\lambda_{-w}=1$, of rank
four. For $t$ not CM, item (XLVI) of \Cref{sec:verification} bounds the
$\Sp(V,\psi)$-invariants of ${\textstyle\bigwedge^{*}}(V\oplus V)$ from above by the stated
dimensions and shows that products of the three divisor classes span spaces
of those dimensions. At a CM point the invariants are spanned by the
monomials of weight zero for $L(X)$; a set of weights $\pm e_{j}$ with sum
zero is a union of opposite pairs, so each such monomial is, up to sign, a
product of monomials of degree two and weight zero, which are divisor
classes, and item (XLVI) counts them: their numbers are the coefficients of
$(1+4y+y^{2})^{4}$.

(iii) If $t$ is not a CM point, the $\Sp(V,\psi)$-invariant combinations of the
$\pi$ are those with $a_{1}=a_{2}=a_{3}$ (\Cref{prop:mumforddivisor}), and
the coefficients of $\omega_{t}$ are the distinct conjugates of $\alpha$
(\Cref{lem:conjugates}). At a CM point use the map $x\mapsto E_{x}$ of the
proof of \Cref{lem:oneexceptional}, which is $\Sp(V,\psi)$-equivariant: it
carries $\omega_{t}$ to an element $(r,\beta)$ of $\QQ\times F$ with
$\beta\notin\QQ$, acting on $U_{ij}$ by the conjugate $\sigma_{k}(\beta)$
(\Cref{prop:mumfordrm}(ii)), and $\omega_{t}$ is $L(X)$-invariant only if
$E_{\omega_{t}}$ preserves the weight spaces of $L(X)$ in ${\textstyle\bigwedge^{2}}V$. Take
weight vectors $v=x\otimes y\otimes z$ and $w=x\otimes y'\otimes z'$ of
weights $(1,1,1)$ and $(1,-1,-1)$ for the Hodge torus. They are not
opposite, so $v\wedge w$ spans a weight space of $L(X)$, every weight of
$L(X)$ on $V$ having multiplicity one. But
$v\wedge w=(x\cdot x)\otimes\bigl(y\cdot y'\otimes z\wedge z'+y\wedge y'
\otimes z\cdot z'\bigr)$ has nonzero components in $U_{12}$ and in $U_{13}$,
on which $E_{\omega_{t}}$ acts by $\sigma_{3}(\beta)\ne\sigma_{2}(\beta)$. So
the line is not preserved. Item (XLVI) confirms the conclusion in exact
arithmetic: the components of $\pi_{12}$, $\pi_{13}$, $\pi_{23}$ of nonzero
weight for $L(X)$ span a space of dimension two, with the relation
$(1,1,1)$ only. The rest of (iii) follows from (i).

(iv) Item (XLVI) generates the module by applying the $36$ elements
$x\mapsto\psi(u,x)v+\psi(v,x)u$ of $\mathfrak{sp}(V,\psi)$ until the span
stops growing: it has dimension $308$ and highest weight $(2,2,0,0)$ for the
torus of diagonal matrices in $\Sp(V,\psi)$, and $308$ is, by Weyl's
dimension formula, the dimension of
the irreducible module of highest weight $2\varpi_{2}$; and adding the four
$\pi$ raises its dimension by two.
\end{proof}

\begin{remark}[Where a construction has to leave $\mathcal{L}$]
\label{rem:leavel}
\Cref{thm:lefschetzclosure} extends \Cref{prop:mumforddivisor} from sums of
line bundles to everything the standard operations make out of them, at every
point of the curve, including the CM points, where the exceptional classes
are algebraic: the cycles that represent them there are not Chern characters
of objects of $\mathcal{L}$, and neither are the two Hodge classes of $X_{c}$
in degree four that are not products of divisor classes. At a point that is
not CM every Chern character of an object of $\mathcal{L}$ lies in the
isotypic component of the trivial representation of $\Sp(V,\psi)$, while by
(iv) an object with $\ch(E)=N\omega$ must have a nonzero component in the one
of highest weight $2\varpi_{2}$. Every construction of such an object
therefore contains a step that does not commute with the Lefschetz group on
cohomology: passing to a direct summand, a subobject or a Jordan--H\"older
factor, choosing a point of a moduli space, or deforming from a member with
more endomorphisms. The last is what semiregularity does along the curve
(\Cref{thm:mumfordnumerical}).
\end{remark}

\begin{lemma}[The summands of a semiregular object]\label{lem:semiregsum}
Let $Z$ be a smooth projective variety and $E_{1},E_{2}$ perfect complexes on
$Z$. On $\Ext^{2}(E_{1}\oplus E_{2},E_{1}\oplus E_{2})=\bigoplus_{i,j}
\Ext^{2}(E_{j},E_{i})$ the semiregularity map is
$\sigma_{E_{1}\oplus E_{2}}(\zeta)=\sigma_{E_{1}}(\zeta_{11})
+\sigma_{E_{2}}(\zeta_{22})$. So $\sigma_{E_{1}\oplus E_{2}}$ is injective if
and only if $\sigma_{E_{1}}$ and $\sigma_{E_{2}}$ are, and
$\Ext^{2}(E_{1},E_{2})=\Ext^{2}(E_{2},E_{1})=0$.
\end{lemma}

\begin{proof}
The Atiyah class is natural, so $\mathrm{At}(E_{1}\oplus E_{2})$ is the
diagonal sum of $\mathrm{At}(E_{1})$ and $\mathrm{At}(E_{2})$. In
$\sigma(\zeta)=\sum_{q}\frac1{q!}\operatorname{tr}(\zeta\circ\mathrm{At}^{q})$
of \Cref{thm:BF} the composite $\zeta\circ\mathrm{At}^{q}$ has diagonal blocks
$\zeta_{ii}\circ\mathrm{At}(E_{i})^{q}$, and the trace of a block matrix of
morphisms is the sum of the traces of its diagonal blocks.
\end{proof}

\begin{corollary}[No Hochschild bypass, and no free splitting]
\label{cor:mumfordsplit}
Let $c\in C$, $Y=X_{c}\times X_{c}$, $\omega=\omega_{a}$ a rational
exceptional class with $L_{0}(a)\ne0$, and $E$ a perfect complex on $Y$ with
$\ch(E)=N\omega$, $N\ne0$.
\begin{enumerate}[label=\textup{(\roman*)},leftmargin=2.6em,itemsep=2pt]
\item $E$ extends to the first order deformation of $D^{b}(Y)$ given by
  $x\in HH^{2}(Y)=HT^{2}(Y)$ only if $x\in\CC\kappa$. No $B$-field, no
  Poisson bivector and no commutative deformation off the curve carries it.
\item If $\sigma_{E}$ is injective and $E=E_{1}\oplus E_{2}$ with both
  summands nonzero, then $\Ext^{2}(E_{1},E_{2})=\Ext^{2}(E_{2},E_{1})=0$, each
  $\sigma_{E_{i}}$ is injective, each $E_{i}$ extends to first order along the
  curve, and $\kappa\lrcorner\ch(E_{i})=0$. If moreover
  $\dim\Ext^{2}(E,E)=119$, then $\dim\Ext^{2}(E_{i},E_{i})=
  120-\dim\Ann_{HT^{2}}(\ch E_{i})$, and
  \[
    \Ann_{HT^{2}}(\ch E_{1})+\Ann_{HT^{2}}(\ch E_{2})=HT^{2}(Y),\qquad
    \Ann_{HT^{2}}(\ch E_{1})\cap\Ann_{HT^{2}}(\ch E_{2})=\CC\kappa .
  \]
\end{enumerate}
\end{corollary}

\begin{proof}
(i) The kernel of $\operatorname{ev}_{E}$ on $HH^{2}(Y)$ is the space of
deformations of $D^{b}(Y)$ along which $E$ deforms to first order
\cite{Tod09}, \cite[Section~8.3]{Mar25a}. If $\operatorname{ev}_{E}(x)=0$ then
$N\,x\lrcorner\,\omega=\sigma_{E}(\operatorname{ev}_{E}(x))=0$, so
$x\in\CC\kappa$ by \Cref{thm:mumfordrigid}(ii). The three kinds of
deformation named are the three summands of $HT^{2}(Y)$.

(ii) The first two assertions are \Cref{lem:semiregsum}. By
\Cref{thm:mumfordnumerical}(ii), $\operatorname{ev}_{E}(\kappa)=0$, and
$\operatorname{ev}_{E}$ is the diagonal sum of $\operatorname{ev}_{E_{1}}$
and $\operatorname{ev}_{E_{2}}$, so each vanishes at $\kappa$; then
$\kappa\lrcorner\ch(E_{i})=\sigma_{E_{i}}(\operatorname{ev}_{E_{i}}(\kappa))
=0$. Put $e(x)=120-\dim\Ann_{HT^{2}}(x)$, the rank of contraction into $x$.
It is subadditive, and by \Cref{thm:mumfordprofile}(i)
$\dim\Ext^{2}(E_{i},E_{i})\ge e(\ch E_{i})$. The cross terms vanish, so
\[
  119=\dim\Ext^{2}(E,E)=\sum_{i}\dim\Ext^{2}(E_{i},E_{i})
  \ge e(\ch E_{1})+e(\ch E_{2})\ge e(N\omega)=119,
\]
and every inequality is an equality. Hence
$\dim\Ann(\ch E_{1})+\dim\Ann(\ch E_{2})=121$. The intersection of the two
annihilators lies in $\Ann(N\omega)=\CC\kappa$ and contains $\kappa$, so it is
$\CC\kappa$, and the sum has dimension $120$.
\end{proof}

Part (ii) is the Mumford form of \Cref{thm:nosum}: a semiregular object
attaining the bound either does not split, or splits as two pieces whose
annihilators fill $HT^{2}(Y)$ and meet only in the direction of the curve, the
analogue of the two branches $P$ and $Q$ of \Cref{rem:twobranches}.

\begin{proposition}[Tate classes at every closed point of every reduction]
\label{prop:modp}
There are a finitely generated subring $R\subset\CC$, a model $W_{R}\to
C_{R}$ of $W\to C$ over $R$, with $C_{R}$ smooth and projective over $R$ with
geometrically connected fibres and $W_{R}$ an abelian scheme, projective over
$C_{R}$, and a finite surjective morphism $C'_{R}\to C_{R}$ from a second curve of the same
kind, with the following properties. Put $Y'=(W_{R}\times_{C_{R}}W_{R})
\times_{C_{R}}C'_{R}$, with projection $\pi'\colon Y'\to C'_{R}$, and let
$\omega$ be a rational exceptional class. For every prime $\ell$ there is a
global section $\omega_{\ell}$ of the lisse sheaf
$R^{4}\pi'_{*}\QQ_{\ell}(2)$ on $C'_{R}[1/\ell]$ whose restriction to
$C'_{R}\otimes_{R}\CC$ is the pull-back of $\omega\otimes1$; write
$\omega_{\ell}(\bar{x})$ for its value at a geometric point over a point $x$.
\begin{enumerate}[label=\textup{(\roman*)},leftmargin=2.6em,itemsep=2pt]
\item At a closed point $x$ of $C'_{R}[1/\ell]$ the residue field
  $\kappa(x)$ is finite, $Y'_{x}$ is the square of an abelian fourfold over
  it, and $\omega_{\ell}(\bar{x})$ is fixed by $\Gal(\bar{\kappa}/\kappa(x))$.
\item At such a point $\omega_{\ell}(\bar{x})$ is a $\QQ_{\ell}$-linear
  combination of classes of algebraic cycles of codimension two on $Y'_{x}$,
  defined over $\kappa(x)$.
\item Let $s$ be a closed point of $\Spec R[1/\ell]$ and $\eta_{s}$ the generic
  point of the reduction $C'_{s}$. Then $\omega_{\ell}(\bar\eta_{s})$ is fixed
  by the Galois group of the function field $\kappa(\eta_{s})$, and it is not
  a $\QQ_{\ell}$-linear combination of products of divisor classes of
  $Y'_{\bar\eta_{s}}$.
\end{enumerate}
Every closed point $x$ of a reduction $C_{s}$ of $C$ is the image of a closed
point $x'$ of $C'_{s}$, and $Y'_{x'}=Y_{x}\otimes_{\kappa(x)}\kappa(x')$; so
(i) and (ii) hold at every closed point of every reduction of $C$, after a
finite extension of its residue field.
\end{proposition}

\begin{proof}
\emph{The models.} The family $W\to C$ is defined over a subfield
$k_{1}\subset\CC$ finitely generated over $\QQ$, say as $W_{1}\to C_{1}$. Let
$K_{1}$ be the function field of $C_{1}$ and fix an embedding over $k_{1}$ of
an algebraic closure $\bar{K}_{1}$ into $\CC$; it sends the geometric generic
point of $C_{1}$ to a point $t\in C$, and $Y_{t}=X_{t}\times X_{t}$ is the base
change of $Y_{\bar{K}_{1}}$. The class $\omega_{t}$ is a Hodge class on an
abelian variety, hence absolutely Hodge, so it comes from an absolutely Hodge
class of $Y_{\bar{K}_{1}}$, and $\Gal(\bar{K}_{1}/K_{1})$ acts on these
through a finite quotient \cite[Theorem~2.11 and Proposition~2.9]{Del82}. Let
$K'\subset\bar{K}_{1}$ be a finite extension of $K_{1}$ on which the action is
trivial, $k'$ the algebraic closure of $k_{1}$ in $K'$, and $C'$ the smooth
projective curve over $k'$ with function field $K'$; it is geometrically
connected over $k'$ and maps finitely onto $C_{1}\otimes_{k_{1}}k'$.
Spreading out $W_{1}\to C_{1}$, $C'$ and this map over a finitely generated
subring of $k'$, and inverting one element of it, gives a normal $R$ and
models with the stated properties.

\emph{The section.} By smooth and proper base change
$R^{4}\pi'_{*}\QQ_{\ell}(2)$ is lisse on $C'_{R}[1/\ell]$. Its stalk at the
geometric generic point $\Spec\bar{K}_{1}$ is $H^{4}(Y_{\bar{K}_{1}},
\QQ_{\ell}(2))=H^{4}(Y_{t},\QQ)\otimes\QQ_{\ell}$, and $\omega_{t}\otimes1$
is the $\ell$-adic component of an absolutely Hodge class, so it is fixed by
$\Gal(\bar{K}_{1}/K')$. The scheme $C'_{R}[1/\ell]$ is normal and integral
with function field $K'$, so $\Gal(\bar{K}_{1}/K')$ maps onto its fundamental
group \cite[Tag~0BQI]{Stacks}, and the vector extends to a global section
$\omega_{\ell}$. On the connected curve $C'_{R}\otimes_{R}\CC$ it agrees with
the pull-back of $\omega\otimes1$, the two being global sections of the same
local system that agree at a point over $t$.

(i) and (ii). A scheme of finite type over $\ZZ$ has finite residue fields at
its closed points. A global section of a lisse sheaf is fixed by the
fundamental group of every point, which is (i); and (ii) is Li's theorem
\cite[Theorem~1.1]{Li26}, with $n=r=2$, for the abelian fourfold that is the
fibre of $W_{R}$ at the image of $x$, base changed to $\kappa(x)$.

(iii) The first assertion holds for the same reason as (i). For the second,
let $\bar\eta$ be a geometric generic point of $C'_{R}$ and $\bar{s}$ a
geometric point over $s$. The morphism $C'_{R}\to\Spec R$ is flat and proper
with geometrically connected and geometrically reduced fibres, so the
specialisation map $\pi_{1}(C'_{\bar\eta})\to\pi_{1}(C'_{\bar{s}})$ is
surjective \cite[Tag~0C0P]{Stacks}, and it is compatible with the actions of
the two groups on the first cohomology of the fibres of $W_{R}$, identified by
cospecialisation. The image of $\pi_{1}(C'_{\bar\eta})$ contains that of
$\pi_{1}(C'_{\CC})$, the closure of the image of the topological fundamental
group of a finite cover of $C$, whose Zariski closure contains
$G_{\QQ_{\ell}}$ because the connected monodromy group of $W\to C$ is $G$
(proof of \Cref{cor:mumfordlefschetz}). So the Zariski closure of the
monodromy of $W_{R}$ over $C'_{\bar{s}}$ contains $G_{\QQ_{\ell}}$. A divisor
class of $Y'_{\bar\eta_{s}}$ is defined over a finite extension of
$\kappa(\eta_{s})$, so it is fixed by an open subgroup of the Galois group,
whose image in $\pi_{1}(C'_{s})$ is open \cite[Tag~0BQI]{Stacks} and meets
$\pi_{1}(C'_{\bar{s}})$ in an open subgroup; the Zariski closure of the image
of an open subgroup has the same identity component, so the class is
$G_{\QQ_{\ell}}$-invariant. The $G$-invariants of ${\textstyle\bigwedge^{2}}(V\oplus V)$ are the
three divisor classes of \Cref{prop:mumfordproduct}(ii), which are
$\Sp(V,\psi)$-invariant, so every product of divisor classes of
$Y'_{\bar\eta_{s}}$ is $\Sp(V,\psi)$-invariant. The class
$\omega_{\ell}(\bar\eta_{s})$ corresponds to $\omega\otimes1$, whose
coefficients $a_{1},a_{2},a_{3}$ are the distinct conjugates of $\alpha$
(\Cref{lem:conjugates}), and the $\Sp(V,\psi)$-invariant classes among the
$\omega_{a}$ are those with $a_{1}=a_{2}=a_{3}$ (\Cref{prop:mumforddivisor}).
The last sentence of the statement holds because $C'_{R}\to C_{R}$ is finite
and surjective.
\end{proof}

At the generic point of each reduction the class is therefore a Tate class
over a global function field that no product of divisor classes explains, the
exact analogue of the situation over $\CC$, while at every closed point it is
explained by cycles. The next theorem says that a bound on those cycles is
all that is missing, with no passage through $p$-adic Hodge theory.

\begin{theorem}[The Mumford target is a degree bound in positive
characteristic]\label{thm:modpdensity}
Keep \Cref{prop:modp}, fix a prime $\ell$ and a relatively ample line bundle
on $Y'\to C'_{R}$, and for a finite set $\Phi$ of tuples of Hilbert
polynomials say that $\omega$ is \emph{$\Phi$-represented} at a point $x$ of
$C'_{R}[1/\ell]$ if $\omega_{\ell}(\bar{x})=\sum_{i}q_{i}\operatorname{cl}(Z_{i})$
for some $q_{i}\in\QQ_{\ell}$ and closed subschemes $Z_{i}\subset Y'_{\bar{x}}$
of pure codimension two whose Hilbert polynomials form a tuple in $\Phi$. The
following are equivalent.
\begin{enumerate}[label=\textup{(\alph*)},leftmargin=2.6em,itemsep=2pt]
\item $\omega_{t}$ is algebraic for every $t\in C$.
\item For some finite $\Phi$, the closed points of $C'_{R}[1/\ell]$ at which
  $\omega$ is $\Phi$-represented are Zariski dense in $C'_{R}$.
\item For some finite $\Phi$ and some Zariski dense set $S$ of closed points of
  $\Spec R[1/\ell]$, the class $\omega$ is $\Phi$-represented at infinitely
  many closed points of $C'_{s}$ for every $s\in S$.
\end{enumerate}
For a single closed point $s$ of $\Spec R[1/\ell]$, the class $\omega$ is
$\Phi$-represented at infinitely many closed points of $C'_{s}$, for some
finite $\Phi$, if and only if $\omega_{\ell}(\bar\eta_{s})$ is a
$\QQ_{\ell}$-linear combination of classes of algebraic cycles on
$Y'_{\bar\eta_{s}}$.
\end{theorem}

\begin{proof}
One standard property of the $\ell$-adic cycle class is used beyond those over
a field: if $\mathcal{O}$ is a discrete valuation ring in which $\ell$ is
invertible, $\mathcal{Y}\to\Spec\mathcal{O}$ is smooth and projective, and
$\mathcal{Z}\subset\mathcal{Y}$ is a closed subscheme flat over $\mathcal{O}$,
then the specialisation isomorphism of $\ell$-adic cohomology carries the
cycle class of the geometric generic fibre of $\mathcal{Z}$ to that of its
geometric special fibre \cite[Cycle]{SGA4h}. The specialisation isomorphism
also carries the values of a global section of a lisse sheaf to each other,
by its construction.

\emph{Step 1: the represented loci are closed.} For a tuple
$\underline{P}=(P_{1},\dots,P_{m})$ let $\cH_{\underline{P}}$ be the fibre
product over $C'_{R}[1/\ell]$ of the relative Hilbert schemes
$\Hilb^{P_{i}}(Y'/C'_{R}[1/\ell])$. It is projective over $C'_{R}[1/\ell]$
\cite{Gro61}, so it has finitely many connected components. On a component
$T$, with universal subschemes $\mathcal{Z}_{i}$, put $A(h)=\{q\in
\QQ_{\ell}^{m}:\sum_{i}q_{i}\operatorname{cl}(\mathcal{Z}_{i,\bar{h}})
=\omega_{\ell}(\bar{h})\}$ for $h\in T$. If $h'$ is a specialisation of $h$,
there is a discrete valuation ring mapping to $T$ with generic point to $h$ and
closed point to $h'$ \cite[Tag~054F]{Stacks}, and the property above, applied
to the pull-backs of $Y'$ and of the $\mathcal{Z}_{i}$, gives $A(h)=A(h')$.
Any two points of the connected noetherian scheme $T$ are joined by a chain of
specialisations and generisations, through the generic points of its
irreducible components, so $A$ is constant on $T$, say $A_{T}$. Let
$\Sigma_{\underline{P}}$ be the union of the images in $C'_{R}[1/\ell]$ of the
components $T$ with $A_{T}\ne\emptyset$. It is closed, as a finite union of
images of proper morphisms, and it is the set of points at which $\omega$ is
$\{\underline{P}\}$-represented: at a point $x$ of the image of such a $T$, a
closed point of the fibre of $T$ over $x$ has residue field finite over
$\kappa(x)$ and gives subschemes of $Y'_{\bar{x}}$. Put
$\Sigma_{\Phi}=\bigcup_{\underline{P}\in\Phi}\Sigma_{\underline{P}}$.

\emph{Step 2: (b) implies (a).} $\Sigma_{\Phi}$ is closed and contains a dense
set, so it is $C'_{R}[1/\ell]$ and contains the generic point $\xi$ of
$C'_{R}$, whose residue field is $K'$. So $\omega_{t}\otimes1=
\omega_{\ell}(\bar\xi)=\sum_{i}q_{i}\operatorname{cl}(Z_{i})$ with
$q\in\QQ_{\ell}^{m}$ and subschemes $Z_{i}$ of $Y_{\bar{K}_{1}}$, and so of
$Y_{t}$. The cycle class maps are compatible with the comparison isomorphism
$H^{4}(Y_{t},\QQ)\otimes\QQ_{\ell}\cong H^{4}(Y_{t},\QQ_{\ell})$
\cite[Section~1 and Example~2.1(a)]{Del82}, so the linear system
$\omega_{t}=\sum_{i}q_{i}[Z_{i}]$, which has rational coefficients, is
solvable over $\QQ_{\ell}$, hence over $\QQ$, and $\omega_{t}$ is algebraic.
The $Z_{i}$ are defined over a finite extension of $K'$, so they spread out to
subschemes, flat over a nonempty open subset \cite[Tag~052A]{Stacks}, of the
pull-back of $W\times_{C}W$ to a finite cover of $C$, and by the local
constancy of the cycle class used in the proof of \Cref{thm:closedlocus},
$\omega$ is algebraic at every point of a nonempty Zariski open subset of $C$,
which is uncountable. \Cref{cor:mumfordequiv} gives (a).

\emph{Step 3: (c) implies (b).} Let $D$ be the set of closed points at which
$\omega$ is $\Phi$-represented and $\bar{D}$ its closure. For $s\in S$ the set
$D\cap C'_{s}$ is infinite, so its closure is the irreducible curve $C'_{s}$,
and $C'_{s}\subseteq\bar{D}$. The points $u$ of $\Spec R[1/\ell]$ over which
the fibre of $\bar{D}$ has dimension at least one form a closed set, by the
upper semicontinuity of the dimension of the fibres of a proper morphism
\cite[Tag~0D4I]{Stacks}; it contains the dense set $S$, so it contains the
generic point. The generic fibre of $\bar{D}$ is then a closed subset of
dimension one of the irreducible curve $C'_{R}\otimes_{R}k'$, so it contains
$\xi$, and $\bar{D}=C'_{R}[1/\ell]$.

\emph{Step 4: (a) implies (c).} By (a) and Step 1, applied over $\bar{K}_{1}$
with the Betti class in place of the $\ell$-adic one, $\omega_{t}$ is the
class of a rational combination of subschemes of $Y_{\bar{K}_{1}}$, defined
over a finite extension $L$ of $K'$. Enlarge $R$ to a finitely generated
subring of the algebraic closure of $k'$ in $L$, let $C''_{R}\to C'_{R}$ be the
normalisation in $L$, and spread the subschemes out over $C''_{R}$; by generic
flatness \cite[Tag~052A]{Stacks} they are flat over a dense open subset $U$ of
$C''_{R}[1/\ell]$, with constant Hilbert polynomials. By Step 1 applied on $U$,
$\omega$ is represented by their fibres at every point of $U$. The image of $U$
in $C'_{R}[1/\ell]$ contains a dense open subset $U'$, every closed point of
$U'$ is the image of a closed point of $U$, and $U'$ meets $C'_{s}$ in a
nonempty open, hence infinite, subset for every $s$ in a dense open subset of
$\Spec R[1/\ell]$.

\emph{The single reduction.} If $\omega$ is $\Phi$-represented at infinitely
many closed points of $C'_{s}$, then $\Sigma_{\Phi}\cap C'_{s}$ is closed and
infinite, so it is $C'_{s}$ and contains $\eta_{s}$. Conversely, cycles at
$\bar\eta_{s}$ are defined over a finite extension of $\kappa(\eta_{s})$, and
spreading them out as in Step 4 over a finite cover of $C'_{s}$ represents
$\omega$, with fixed Hilbert polynomials, at every closed point of a nonempty
open subset of $C'_{s}$.
\end{proof}

\begin{remark}[What is left in positive characteristic]\label{rem:modpremains}
By \Cref{prop:modp}(ii) the class is represented at every closed point of
$C'_{R}[1/\ell]$, so \Cref{thm:modpdensity} turns the Mumford target into a
statement of uniformity alone: the Hilbert data of some representation must
stay in one finite set on a Zariski dense set of closed points, for instance
at infinitely many closed points of infinitely many reductions. The passage
to characteristic zero is then done by the specialisation of cycle classes
along $C'_{R}$, and neither the $p$-adic variational Hodge conjecture nor a
lifting of cycles is needed. Li's theorem gives no bound, and where its proof
carries cycles into a prescribed isogeny class it does so by an isogeny whose
degree it does not control \cite[proof of Theorem~1.2, Step~3]{Li26}. Within
one reduction the
bound is equivalent to the class being algebraic at the generic point of
$C'_{s}$, a Tate class over a global function field that is not a combination
of products of divisor classes (\Cref{prop:modp}(iii)); this is the
codimension two case of the variational Tate conjecture along $C'_{s}$, which
Morrow proves for divisors in its crystalline form \cite{Mor19}. From one
reduction alone the way back to characteristic zero would be the $p$-adic
variational Hodge conjecture for $W\times_{C}W$ over the Witt vectors of
$\kappa(s)$, with rational rather than $\ell$-adic coefficients and with the
crystalline class of the cycle matching the de Rham class of $\omega$; Bloch,
Esnault and Kerz lift such a class in $K_{0}\otimes\QQ$ of the special fibre
to the formal completion for $p>15$, the dimension being nine, and leave the
algebraisation open \cite[Theorem~1.3]{BEK14}. Neither the uniform bound nor
this conjecture is proved.
\end{remark}

The five routes of \Cref{rem:bypassaudit} that are not closed off are
\emph{specialisation}, \emph{reduction modulo $p$}, \emph{semiregularity},
\emph{Kuga--Satake} and \emph{motivated classes}. They are not five
obligations.

\begin{corollary}[The open routes are one statement]\label{cor:mumfordone}
Let $W\to C$ be as in \Cref{cor:mumfordlefschetz}, let $\omega$ be a rational
exceptional class, and let the \emph{Mumford target} be the first case of
\textup{(F2)}: the Hodge conjecture for $X_{t}\times X_{t}$ for every $t\in C$.
\begin{enumerate}[label=\textup{(\roman*)},leftmargin=2.6em,itemsep=2pt]
\item Each of the following is equivalent to the Mumford target: $\omega_{t}$ is
  algebraic for uncountably many $t\in C$; $B(W\times_{C}W)$ holds for
  $L_{W\times_{C}W}$; $\omega$ is represented with Hilbert data in one finite
  set at infinitely many points of $C$, for instance at infinitely many CM
  points; $\omega$ is represented with Hilbert data in one finite set on a
  Zariski dense set of closed points of the arithmetic model of
  \Cref{thm:modpdensity}.
\item Each of the following implies the Mumford target: a perfect complex $E$
  at one point of $C$ with $\ch(E)=N\omega$, $N\ne0$, no other component, and
  $\dim\Ext^{2}(E,E)=119$; the algebraicity, for uncountably many $t$, of the
  Kuga--Satake class of the surface $S_{\lambda}$ of \Cref{thm:mumfordks}
  attached to $X_{t}$; a rational map as in \Cref{prop:hksquare}(iv), under
  its hypotheses, from $X_{t}\times X_{t}$ for uncountably many $t$.
\item There is a countable set $\Sigma_{0}\subset C$ containing the CM points
  such that the Mumford target holds as soon as the Hodge conjecture holds for
  $X_{c}\times X_{c}$ at one point $c\notin\Sigma_{0}$.
\end{enumerate}
\end{corollary}

\begin{proof}
(i) The first two are \Cref{cor:mumfordequiv}, with \Cref{lem:oneexceptional}
for one class in place of two, the third is \Cref{cor:mumfordbounded} and the
fourth is \Cref{thm:modpdensity}. (ii) The first is
\Cref{thm:mumfordnumerical}(iii). The second and the third give the Hodge
conjecture for $X_{t}\times X_{t}$ at each of the uncountably many points, by
\Cref{thm:mumfordks}(ii) and \Cref{prop:hksquare}(iv), and then (i) applies.
(iii) In the proof of \Cref{cor:mumfordbounded} the set of points at which
$\omega$ is algebraic is the union of the countably many closed sets
$\Sigma_{\underline{P},\underline{q}}$, each finite or equal to $C$. Let
$\Sigma_{0}$ be the union of the finite ones and of the CM points. If
$\omega_{c}$ is algebraic at a point $c\notin\Sigma_{0}$, then $c$ lies in
some $\Sigma_{\underline{P},\underline{q}}$ that is not finite, which is $C$,
and (i) applies.
\end{proof}

\begin{remark}[Which routes are needed]\label{rem:mumfordneeded}
Of the five routes, specialisation, reduction modulo $p$ and motivated
classes are the Mumford target itself in three other forms, by
\Cref{cor:mumfordone}(i) and, for the motivated classes,
\Cref{cor:mumfordequiv}; semiregularity and Kuga--Satake are stronger inputs
that imply it. So one route suffices and the other four can be dropped, and
item (XLIX) of \Cref{sec:verification} computes this from the rules of the
paper: every minimal set of open statements from which the target follows
has one element, and the statements equivalent to it over what is proved
are exactly the four forms of (i) and the target. Nothing in the paper
proves any of them, and each of the three forms is the whole of the problem:
a bound that no known construction gives, or a case of \textup{(L)} that is
equivalent to the target. The two stronger inputs ask for an object and a
cycle that nothing in the paper constructs.

Nor is the target, once proved, a step that finishes the conjecture. It is a
special case of the conjecture, and so it is needed in the sense that every
proof of the conjecture proves it. But it is only the first case of (F2), and
a proof of it leaves (F1), every other case of (F2), and (F3) open. The last
of these is not a remainder of the conjecture but the conjecture itself
(\Cref{prop:f3ishc}); its weaker form \textup{(F3$'$)} is open as well
(\Cref{rem:f3primeopen}). Of the five minimal sufficient sets of
\Cref{thm:closure}(ii), only the route through this paper,
$\{\mathrm{P2},\mathrm{F2},\mathrm{F3}'\}$, has the target as an input,
as the first case of (F2). On the other four it is not an input: (F3) is the
conjecture, and \textup{(L)} implies \textup{(V)} (\Cref{prop:lefvhc}),
which carries the algebraicity at the CM points of \Cref{thm:mumfordcm} to
every member. Item (XLIX) checks both: that (F3), being equivalent to the
conjecture, yields the target, and that \textup{(L)} and \textup{(V)} do.
\end{remark}

\subsection{Summary}

\begingroup
\renewcommand{\arraystretch}{1.25}
\begin{longtable}{@{}>{\raggedright\arraybackslash}p{0.30\linewidth}>{\raggedright\arraybackslash}p{0.30\linewidth}>{\raggedright\arraybackslash}p{0.30\linewidth}@{}}
\caption{The routes considered, and the boundary of each result.}
\label{tab:scope}\\
\toprule
\textbf{Route} & \textbf{Status here} & \textbf{Not covered} \\
\midrule
\endfirsthead
\multicolumn{3}{@{}l}{\footnotesize\itshape Table \thetable, continued.}\\[2pt]
\toprule
\textbf{Route} & \textbf{Status here} & \textbf{Not covered} \\
\midrule
\endhead
\midrule
\multicolumn{3}{r@{}}{\footnotesize\itshape continued on the next page}\\
\endfoot
\bottomrule
\endlastfoot
Ambient object, restricted to $\Av$, transported by $\Phi_{\cP}$
  & closed by \Cref{thm:divisor}
  & ambient varieties with exceptional classes \\
Any $K$-equivariant construction with the norm character
  & closed by \Cref{thm:character}
  & constructions that are not eigenclasses \\
Secant dimension count
  & no obstruction, \Cref{cor:nonumerical}
  & none \\
Semiregularity dimension count
  & no obstruction, \Cref{cor:countnotbind}
  & none \\
Semiregularity of an object whose Chern character sits on the Weil line
  & closed by \Cref{thm:rankzero} and \Cref{cor:kernelfamily} for direct
    sums, and for every object by \Cref{thm:p2false}
  & objects of positive rank; the weaker hypothesis of
    \cite[Question 11.4]{Mar25b} \\
Induction through hyperplane sections
  & closed by \Cref{thm:hyperplane}, \Cref{thm:returntrip}
  & correspondences, degenerations, coniveau \\
Kuga--Satake, from a weight two structure of K3 type
  & not obstructed, \Cref{prop:ksnotrestriction}
  & available only at $n=1$, \Cref{cor:ksboundary} \\
Divisors, on a member with quaternionic multiplication
  & succeeds, \Cref{thm:quatweil}
  & fails at full Hodge group, \Cref{prop:divfails} \\
Propagation off the base locus
  & reduced to one bound or one rank, \Cref{thm:degreebound},
    \Cref{thm:semiregclosure}
  & neither hypothesis proved, \Cref{rem:tworemain} \\
Objects with effective Chern character at a base point
  & closed by \Cref{thm:positivity}
  & objects indecomposable in $D^{b}$, \Cref{rem:whatitmustbe} \\
The weakened criterion for a secant object of positive rank
  & reduced to one identity on the factor, \Cref{thm:factor}; the identity
    fails for Markman's candidate, \Cref{thm:candidatefails}, and at every
    transversal double point, \Cref{cor:localobstruction}
  & objects whose support has no such point, \Cref{rem:candidatemeaning} \\
A secant object of rank one with smooth support
  & invariants forced and the discriminant bounded by $9$,
    \Cref{thm:smoothinvariants}, \Cref{cor:fourdiscriminants}
  & existence for $d\in\{1,3,5,7\}$, and one vanishing in
    $H^{1}(S,N_{S/X})$, \Cref{rem:smoothremains} \\
The local part of the weakened criterion
  & it is a condition on depth: no obstruction at a Cohen--Macaulay point,
    \Cref{thm:cmvanishing}, \Cref{cor:cmdichotomy}
  & Cohen--Macaulay supports, where the criterion becomes one global
    vanishing \\
A support that is a complete intersection
  & closed for every $b$ and every $d$, \Cref{thm:noci}
  & supports of higher Hilbert--Burch rank, \Cref{thm:burchsearch} \\
A support with a Hilbert--Burch resolution by line bundles or by simple
semi-homogeneous bundles
  & closed for every rank, twist and $d$, \Cref{thm:nosplit},
    \Cref{thm:noblocks}
  & resolutions over bundles that are not projectively flat \\
The secant route, granted every demand it makes
  & reaches $\delta=\delta_{0}$, and every $\delta$ by \Cref{prop:descent},
    \Cref{thm:closure}(iv)
  & the CM fields of degree at least four, which give the quadratic ones by
    \Cref{prop:descent}(ii) \\
Propagation from a base point
  & closes every imaginary quadratic family at once, \Cref{thm:degreebound},
    \Cref{thm:semiregclosure}
  & (P2) not proved, \Cref{cor:fourremain} \\
The bounded form (P1), read on the transported cycle
  & closed: the transport spreads each component by $c^{2n}/m_{i}$ and the
    data is unbounded whenever a component has finite stabiliser,
    \Cref{lem:multiplicity}, \Cref{thm:p1forces}
  & nothing; the other reading is closed too \\
The bounded form (P1), read on the minimal representative
  & closed: equivalent to the conjecture for the family,
    \Cref{thm:p1equivalent}
  & nothing; it is a restatement, so it is not a route \\
The semiregularity form (P2), for a given representative
  & forces the evaluation map to vanish on the $4n^{2}$-dimensional
    annihilator, and excludes every representative that is a sheaf on a
    smooth germ of codimension two of its support, \Cref{thm:p2forces},
    \Cref{cor:p2shape}
  & representatives that are genuine complexes along every such component;
    members with an abelian subvariety tangent to an eigenspace \\
The same form, tested by a number
  & $\dim\Ext^{2}(E,E)=2n(2n-1)$ at one base point would suffice,
    \Cref{thm:p2numerical}, and no representative attains it,
    \Cref{thm:p2false}
  & nothing \\
The same form, for a representative in codimension $n$
  & closed at every member of every family: the support must have a
    component of codimension less than $n$, and at a general member every
    component of codimension at least two has multiplicity zero,
    \Cref{lem:finiteproduct}, \Cref{thm:p2support}
  & representatives glued along components of lower codimension with
    alternating length zero \\
The Hochschild action on such a representative
  & the annihilator is the mixed part of a splitting of $HH^{1}$ and the
    action factors through a ring of dimension $2^{2n+1}-1$ in one of
    dimension $2^{4n}$, \Cref{thm:hhsplitting}, \Cref{cor:hhfactor}
  & whether a complex with that action exists \\
The shape of such a representative
  & an indecomposable summand carries the class; at $n=2$, and at $n=3$
    granting the compatibility of \Cref{cor:hhfactor}(iii), its
    endomorphism algebra has dimension at least three, so no simple object
    serves, \Cref{thm:p2endo}
  & nothing: no object of any shape meets the pure form, \Cref{thm:p2false} \\
The pure form of (P2), Chern character exactly a Weil class
  & false for every imaginary quadratic field and every $n\ge2$: a
    semiregular such complex would deform to very general non-algebraic Weil
    tori, whose coherent sheaves have no Chern character below the top
    degree, \Cref{prop:weiltorussheaves}, \Cref{thm:p2false}
  & the corrected form (P2), with powers of the polarisation allowed in the
    Chern character, \Cref{rem:p2prime} \\
The reduction to Weil classes, and to abelian varieties
  & false as a statement about Hodge rings in dimension eight,
    \Cref{prop:mumfordproduct}; impossible Hodge-theoretically,
    \Cref{prop:noabeliantype}
  & the algebraicity of the classes beyond the Weil lines, which is (F2), and
    (F3), which is the whole conjecture, \Cref{prop:f3ishc} \\
The two classes on the self-product of a Mumford fourfold
  & a real multiplication by a cubic field on the primitive part of $H^{2}$,
    algebraic once the Kuga--Satake class of one K3 surface of Picard number
    thirteen is, \Cref{prop:mumfordrm}, \Cref{thm:mumfordks}
  & that Kuga--Satake class, \Cref{rem:mumfordks} \\
Bypasses for the same two classes
  & degeneration and larger families closed off by rigidity; semiregularity
    reduced to $\dim\Ext^{2}=119$ at one CM point; specialisation reduced to
    a degree bound, \Cref{thm:mumfordrigid}, \Cref{thm:mumfordnumerical},
    \Cref{cor:mumfordbounded}
  & an object with $\dim\Ext^{2}=119$, or the degree bound,
    \Cref{rem:bypassaudit} \\
Objects built from line bundles, for the same two classes
  & closed at every point of the curve: the Chern characters stay invariant
    under the Lefschetz group, \Cref{thm:lefschetzclosure}
  & steps that break that symmetry, \Cref{rem:leavel} \\
Twistor deformation of the square, for the same two classes
  & closed: the component of the Hodge locus among all complex tori is the
    Mumford curve, which contains no twistor line, \Cref{prop:notwistor}
  & the Kuga--Satake class, which moves in a seven-dimensional family \\
The five routes left open, for the same two classes
  & one statement: three of them are equivalent to the target and two imply
    it, and every minimal set of open statements yielding the target has one
    element, \Cref{cor:mumfordone}, item (XLIX)
  & the target itself, which is only the first case of (F2),
    \Cref{rem:mumfordneeded} \\
Pull-back from a hyperk\"ahler manifold, for the same two classes
  & the square is holomorphic symplectic and the classes are the twisted dual
    forms of the part of $H^{2}$ its symplectic class generates; a symplectic
    rational map to some $S_{\mu}^{[n]}$, $n\ge5$, suffices, and no K3
    surface or hyperk\"ahler eightfold serves, \Cref{prop:hksquare}
  & such a map, \Cref{rem:hkroute} \\
Reduction modulo $p$, for the same two classes
  & represented by cycles at every closed point of every reduction; the
    target is equivalent to a bound on their Hilbert data on a Zariski dense
    set, \Cref{prop:modp}, \Cref{thm:modpdensity}
  & the bound, \Cref{rem:modpremains} \\
The Weil classes of every CM field
  & the same propagation problem, with a CM base point in every family,
    \Cref{thm:cmbasepoint}, \Cref{thm:cmpropagation}
  & (P2) not proved for any field \\
What is left, as a whole
  & five minimal sets of open statements: $\{\mathrm{F3}\}$ and
    $\{\mathrm{L},\mathrm{M}\}$, each equivalent to the conjecture, and
    three through the conjecture modulo abelian varieties,
    $\{\mathrm{L},\mathrm{F3}'\}$, $\{\mathrm{V},\mathrm{F3}'\}$ and this
    route, $\{\mathrm{P2},\mathrm{F2},\mathrm{F3}'\}$,
    \Cref{thm:closure}, \Cref{prop:f3ishc}, \Cref{cor:fourremain},
    \Cref{rem:bullseye}
  & the Lefschetz standard conjecture, (M), (V), (F2), (F3$'$), all open \\
The conjecture modulo abelian varieties, (F3$'$)
  & holds on the class $\mathcal{A}$ of varieties whose cohomology is
    reached from abelian varieties by algebraic correspondences (curves,
    abelian varieties, products, surjective images), in degrees $0$, $2$,
    $2n-2$, $2n$, and in dimension at most three, \Cref{prop:f3prime}
  & degree four on fourfolds outside $\mathcal{A}$ and beyond, for
    instance squares of K3 surfaces without an algebraic Kuga--Satake
    correspondence, \Cref{rem:f3primeopen} \\
\end{longtable}
\endgroup
